 \documentclass[11pt]{amsart}
\usepackage[dvipsnames]{xcolor}
\usepackage{amsfonts,amssymb,amsmath,amscd,amstext}
\usepackage{mathrsfs}
%\usepackage{mathabx}
\usepackage[colorlinks=true,linkcolor=teal,citecolor=purple,urlcolor=Plum]{hyperref}
\usepackage[utf8]{inputenc}
%\usepackage{microtype}
%\usepackage{graphicx}
%%% ATTIVA IL SEGUENTE PER IMMAGINE %%%
\usepackage[final]{graphicx}
\usepackage{aligned-overset}
%\usepackage{lmodern}
%\usepackage{listings}
%\lstset{
 %   language=Mathematica,
%    basicstyle=\ttfamily\small,  % Stile di base: testo monospaziato e dimensione ridotta
 %   backgroundcolor=\color{gray!10},  % Sfondo grigio chiaro
 %   frame=single,                  % Riquadro attorno al codice
 %   rulecolor=\color{black},       % Colore del bordo del riquadro
%    keywordstyle=\color{blue},     % Parole chiave in blu
%    stringstyle=\color{red},       % Stringhe in rosso
%    commentstyle=\color{green!60!black}, % Commenti in verde
%    showstringspaces=false,        % Nasconde i simboli degli spazi nelle stringhe
%}
%\usepackage{changes}
\usepackage{float}
\usepackage{comment}
\usepackage[noabbrev,capitalize,nameinlink]{cleveref}
\usepackage{mathdots}
\usepackage[a4paper,top=2.3cm,bottom=2.3cm,left=1.4cm,right=1.4cm]{geometry}
\renewcommand{\leq}{\leqslant}
\renewcommand{\geq}{\geqslant}
\renewcommand{\le}{\leqslant}
\renewcommand{\ge}{\geqslant}
\newcommand{\ptl}{\partial}
\newcommand{\td}{T d^\hh}
\newcommand{\hhh}{{\mathcal{H}}}
\newcommand{\hn}{\mathbb{H}^n}
\newcommand{\A}{\mathbf{A}}
\newcommand{\B}{\mathbf{B}}
\newcommand{\C}{\mathbf{C}}
\newcommand{\D}{\mathbf{D}}
\newcommand{\norm}[1]{\| #1 \|}
\newcommand{\g}[1]{g_\varepsilon\left( #1 \right)}
\newcommand{\normale}{N}
\newcommand{\no}{\normale}
\newcommand{\noh}{\normale^\hh}
\newcommand{\rifun}{\mathcal{Q}}
\newcommand{\penafun}{\mathcal{P}}
\newcommand{\n}[1]{\nabla^\varepsilon_{ #1 }}
\newcommand{\tmc}{{\mathscr{H}}}
\newcommand{\timc}{{\left(\mathscr{H}^\hhh\right)^{-1}}}
\newcommand{\hnn}{\ell}
\newcommand{\gh}[1]{g_\hh\left( #1 \right)}
\newcommand{\vh}{\nu}
\newcommand{\fis}{\tilde g_{S,n,k}}
\newcommand{\fisd}{g_{S,k,\omega}}
\newcommand{\LMD}{\mathbf{\Lambda}}
\newcommand{\GT}{\mathbf{\tilde{G}}}
\newcommand{\rr}{{\mathbb{R}}}
\newcommand{\cc}{{\mathbb{C}}}
\newcommand{\hh}{{\mathbb{H}}}
\newcommand{\nn}{{\mathbb{N}}}
\newcommand{\sph}{{\mathbb{S}}}
\renewcommand{\MR}[1]{%
  \href{https://mathscinet.ams.org/mathscinet-getitem?mr=#1}{MR#1}%
}
\newcommand{\la}{\lambda}
\newcommand{\ka}{\kappa}
\newcommand{\Sg}{\Sigma}
\newcommand{\Ru}{\mathbb{R}}
\newcommand{\para}{\omega}
\newcommand{\f}{g}
\newcommand{\sg}{\sigma}
\newcommand{\Om}{\Omega}
\newcommand{\eps}{\varepsilon}
\newcommand{\var}{\varphi}
\newcommand{\ga}{\gamma}
\newcommand{\Ga}{\Gamma}
\newcommand{\area}{\sigma}
\newcommand{\ns}{\nabla^S}
\newcommand{\vol}{V}
\renewcommand{\v}{v^\eps}
\newcommand{\vt}{\v_{2n+1}}
\newcommand{\co}{\cos(\theta)}
\newcommand{\si}{\sin(\theta)}
\newcommand{\Es}{\mathrm{G}}
\newcommand{\Ej}{\mathrm{K}}
\newcommand{\kk}{\kappa}
\newcommand{\mh}{\mathcal{H}}
\newcommand{\Co}[1]{{\bf C}^{#1}}
\newcommand{\ep}{\epsilon}
\newcommand{\df}{d_{K_,\partial\Om}}
\newcommand{\dfz}{d_{K_0,\partial\Om}}
\newcommand{\mA}{\mathcal{A}}
\newcommand{\Lap}{\mathcal{L}}
\newcommand{\s}{\mathcal S}
\newcommand{\scu}{\longrightarrow}
\newcommand{\so}[1]{W^{1,p}(#1)}
\newcommand{\soz}[1]{W_0^{1,p}(#1)}
\newcommand{\lp}[1]{L^{p}(#1)}
\renewcommand{\j}{\mathcal{C}}
\newcommand{\lpl}[1]{L^{p}_{loc}(#1)}
\newcommand{\z}{W^{\varphi}}
\newcommand{\sol}[1]{W^{1,p}_{loc}(#1)}
\newcommand{\x}{\mathfrak{X}}
\newcommand{\normalp}[1]{\Vert#1\Vert_{L^p(A')}}
\newcommand{\anso}[1]{W^{1,p}_X(#1)}
\renewcommand{\th}{\tilde h}
\newcommand{\W}{\mathbf Y}
\newcommand{\Y}{\mathcal X}
\newcommand{\haus}{\mathcal{H}}
\newcommand{\ansol}[1]{W^{1,p}_{X,loc}(#1)}
\newcommand{\cl}[1]{\overline{#1}}
\newcommand{\N}{\mathcal N}
\newcommand{\ciuno}[1]{C^{1}(\overline{#1})}
\newcommand{\lul}{L^1_{loc}(U)}
\newcommand{\ch}{\overline{co}}
\renewcommand{\k}{\kappa}
\newcommand{\sub}{\partial_{X,N}}

\newcommand{\E}{\mathbf{E}}
\newcommand{\X}{\mathbf{X}}
\newcommand{\lagrange}{\mathscr{L}}
\newcommand{\Z}{\mathbf{Z}}
\renewcommand{\n}{\mathbf{N}}
\newcommand{\G}{\mathcal G}
\newcommand{\En}{\mathrm{E}}
\newcommand{\adel}{A(\delta)}
\newcommand{\hold}[1]{C_X^{0,\alpha}(#1)}
\newcommand{\holdl}[1]{C_{X,loc}^{0,\alpha}(#1)}
\newcommand{\barR}{\overline{\mathbb{R}}}
\newcommand{\average}{{\mathchoice {\kern1ex\vcenter{\hrule height.4pt
				width 6pt
				depth0pt} \kern-9.7pt} {\kern1ex\vcenter{\hrule height.4pt width 4.3pt
				depth0pt}
			\kern-7pt} {} {} }}
\newcommand{\ave}{\average\int}

\definecolor{champagne}{rgb}{0.97, 0.91, 0.81}

\newcommand{\simone}[1]{\todo[inline,color=champagne]{Simone: #1}}
\definecolor{asparagus}{rgb}{0.53, 0.66, 0.42}
\DeclareMathOperator{\ric}{Ric}
\DeclareMathOperator{\graf}{graph}
\DeclareMathOperator{\divv}{div}
\DeclareMathOperator{\diver}{div}
\DeclareMathOperator{\trace}{trace}
\DeclareMathOperator{\intt}{int}
\DeclareMathOperator{\hess}{Hess}
\newcommand{\jacobi}{\mathcal{J}}
\DeclareMathOperator{\diam}{diam}
\DeclareMathOperator{\lip}{Lip}
\DeclareMathOperator{\supp}{supp}
\DeclareMathOperator{\spann}{span}
\DeclareMathOperator{\jac}{Jac}
\DeclareMathOperator{\rank}{rank}
\DeclareMathOperator{\tor}{Tor}
\DeclareMathOperator{\R}{R}
\DeclareMathOperator{\sym}{sym}
\DeclareMathOperator{\torn}{Tor_{\nabla}}
\DeclareMathOperator{\tors}{Tor_{\nabla^S}}
\newtheorem{theorem}{Theorem}[section]
\newtheorem{proposition}[theorem]{Proposition}
\newtheorem{definition}[theorem]{Definition}
\newtheorem{lemma}[theorem]{Lemma}
\newtheorem{corollary}[theorem]{Corollary}
\theoremstyle{definition}
\newtheorem{remark}[theorem]{Remark}
\newtheorem{example}[theorem]{Example}
\renewcommand{\theexample}{\!}
\newtheorem{problem}{Problem}
\renewcommand{\theproblem}{\!\!}
\newcommand{\V}{{\mathbb{V}}}

\theoremstyle{remark}

\renewcommand{\labelenumi}{\arabic{enumi}.}

\newenvironment{enum}{\begin{enumerate}
\renewcommand{\labelenumi}{(\roman{enumi})}}{\end{enumerate}}

\numberwithin{equation}{section}


%\setcounter{tocdepth}{1}

%\thanks{Key words. Bernstein problem, Heisenberg group, ruled hypersurface, minimal hypersurface, second fundamental form, totally geodesic.\\
%\indent MSC. 53C17, 49Q05}

%\author[G. Giovannardi]{Gianmarco Giovannardi}
%\address[Gianmarco Giovannardi]{Dipartimento di Matematica e Informatica "U. Dini", Università degli Studi di Firenze, Viale Morgani 67/A, 50134, Firenze, Italy}
%\email{gianmarco.giovannardi@unifi.it}

%\author[M.~Fogagnolo]{Mattia Fogagnolo}
%\address[M.~Fogagnolo]{Dipartimento di Matematica ``Tullio Levi-Civita'', Università degli Studi di Padova\protect\newline
%\indent  via Trieste 63, 35131 Padova (PD), Italy}
%\email{mattia DOT fogagnolo AT unipd DOT it}

%\author[A.~Pinamonti]{Andrea Pinamonti} \address[A.~Pinamonti]{Dipartimento di Matematica, Università di Trento\protect\newline
%\indent Via Sommarive, 14, 38123 Povo TN} 
%\email{andrea DOT pinamonti AT unitn DOT it}

%\author[S.~Verzellesi]{Simone Verzellesi}
%\address[S. Verzellesi]{Department of Decision Sciences and BIDSA, Bocconi University \protect\newline
%\indent Via R\"ontgen 1, 20136, Milano (MI), Italy}
%\email{simone DOT verzellesi AT unibocconi DOT it}
\author[S.~Verzellesi]{Simone Verzellesi}
\address[S.~Verzellesi]{Dipartimento di Matematica ``Tullio Levi-Civita'', Università degli Studi di Padova, via Trieste 63, 35131 Padova (PD), Italy}
\email{simone DOT verzellesi AT unipd DOT it}

%\author[F.~Serra Cassano]{Francesco Serra Cassano} \address[Francesco Serra Cassano]{Dipartimento di Matematica\\ Università di Trento\\ Via Sommarive, 14, 38123 Povo TN} \email{francesco.serracassano@unitn.it}



%\author[K.~Zambanini]{Kilian Zambanini} \address[Kilian Zambanini]{Dipartimento di Matematica\\ Università di Trento\\ Via Sommarive, 14, 38123 Povo TN} \email{kilian.zambanini@unitn.it}

\title[Variation formulas for mean curvature functionals]{Variation formulas for mean curvature functionals}



%\author[S.~Verzellesi]{Simone Verzellesi}
%\address[S.~Verzellesi]{Dipartimento di Matematica, Università di Trento, via Sommarive 14, 38123 Povo (TN), Italy}
%\email{simone.verzellesi@unitn.it}

\date{\today}

\subjclass{53C42, 49Q10, 53C17}
\keywords{Variation formulas, mean curvature, Riemannian manifolds, Heisenberg group}
\thanks{\textit{Acknowledgments.} S. Verzellesi thanks R. Monti, J. Pozuelo, M. Ritoré and D. Vittone for fruitful discussion about the addressed topics.
\textit{Memberships and funding information.} S. Verzellesi is member of the Istituto Nazionale di Alta Matematica (INdAM), Gruppo Nazionale per l'Analisi Matematica, la Probabilità e le loro Applicazioni (GNAMPA). 
S.~Verzellesi is supported by the University of Padova. S.~Verzellesi received funding through INdAM-GNAMPA 2026 Project \emph{Variational, Geometric, and Analytic Perspectives on Regularity}, CUP E53C25002010001}

\bibliographystyle{abbrv} 
\begin{document}
\begin{abstract}

We establish first and second variation formulas for general mean curvature functionals under arbitrary variations in arbitrary Riemannian manifolds. We obtain corresponding formulas in the sub-Riemannian Heisenberg group through a Riemannian approximation scheme.
\end{abstract}
\maketitle
%\tableofcontents
\section{Introduction}
%%%%%% BOZZA1%%% %%%%
%Geometric inequalities are among the most powerful mechanisms through which the shape of space is reflected in the behaviour of hypersurfaces. In Euclidean geometry, the isoperimetric inequality, the Minkowski inequality, and the Heintze--Karcher inequality form a remarkably coherent picture: the same distinguished objects, namely round spheres, appear simultaneously as isoperimetric regions, as equality cases for curvature inequalities, and as stable critical points of natural geometric variational problems. This unity is one of the classical signatures of the rigidity of the Euclidean setting.
%
%The purpose of this paper is to investigate what remains of this picture in the sub-Riemannian Heisenberg group. The answer turns out to be substantially more subtle than a direct analogy with the Euclidean theory would suggest. In the Heisenberg group, the interaction between the horizontal distribution, characteristic points, and anisotropic dilations changes the variational nature of curvature functionals in a fundamental way. In particular, the natural candidates inherited from the standard sub-Riemannian geometry---the Pansu sphere and the Korányi sphere---do not play the same universal role as Euclidean spheres. Although the Pansu sphere is the expected isoperimetric profile and the Korányi sphere is the boundary of the canonical homogeneous ball, neither of them is, in general, an equilibrium configuration for the shape optimization problems naturally associated with Minkowski- and Heintze--Karcher-type inequalities.
%
%This observation is the starting point of the present work. We initiate a variational theory for hypersurface functionals driven by the horizontal mean curvature in the Heisenberg group, with particular emphasis on the problem of minimizing the total horizontal mean curvature under a horizontal area constraint. More precisely, for a smooth closed hypersurface \(S\) we consider quantities of the form
%\[
%    \int_S f(H_H)\,d\sigma_H,
%\]
%where \(H_H\) denotes the horizontal mean curvature and \(d\sigma_H\) the sub-Riemannian perimeter measure. The case \(f(r)=r\) leads to the total horizontal mean curvature, whose Euclidean analogue is governed by the Minkowski inequality; the case \(f(r)=r^{-1}\) is related instead to Heintze--Karcher-type phenomena. The Heisenberg setting forces one to confront two intertwined difficulties: the degeneracy at characteristic points and the lack of a direct Riemannian first- and second-variation calculus for the relevant horizontal quantities.
%
% first contribution is therefore structural. We derive first and second variation formulas for general mean-curvature-driven functionals. We first establish the corresponding formulas in arbitrary Riemannian manifolds, and then pass to the Heisenberg group through a Riemannian approximation scheme. This provides a flexible variational framework for horizontal mean curvature functionals and yields, in particular, explicit Euler--Lagrange and stability operators adapted to the sub-Riemannian geometry. Since characteristic points are unavoidable for closed hypersurfaces in the first Heisenberg group, we introduce suitable non-characteristic notions of stationarity and stability, in which variations are localized away from the characteristic set. This perspective is strong enough to detect the correct critical objects, while still retaining the singular features of the ambient geometry.
%
%We then specialize the theory to rotationally invariant surfaces in \(\mathbb H^1\). In this class the variational problem becomes explicitly tractable, but remains far from trivial. The total horizontal mean curvature and the horizontal area reduce to one-dimensional functionals on the profile curve, and the Euler--Lagrange equation can be solved completely. The outcome is a new one-parameter family of closed rotationally invariant critical surfaces, which we call \emph{Pansu--Minkowski spheres}. These surfaces interpolate, in a precise sense, between the familiar constant-horizontal-mean-curvature geometry of the Pansu sphere and the curvature optimization problem dictated by the Minkowski functional.
%
%More precisely, we prove that every closed, rotationally invariant, horizontally mean convex surface in \(\mathbb H^1\) which is stationary for the total horizontal mean curvature under horizontal-area-preserving non-characteristic variations is, up to vertical translations and intrinsic dilations, one of these Pansu--Minkowski spheres. Thus the variational problem selects a family of profiles different from the Korányi geometry and, except at a limiting endpoint, distinct from the standard Pansu sphere. Among this family, a single member is distinguished by the normalized total mean curvature: the optimal Pansu--Minkowski sphere uniquely minimizes the scale-invariant functional
%\[
%    \sigma_H(S)^{-2/3}\int_S H_H\,d\sigma_H
%\]
%within the class of Pansu--Minkowski spheres.
%The stability picture is equally revealing. We show that Pansu--Minkowski spheres are stable, and in fact satisfy a quantitative positivity estimate for the second variation, under rotationally invariant non-characteristic perturbations. This stability is strong enough to imply a local minimality result for the total horizontal mean curvature under fixed horizontal area, within a natural class of rotationally invariant perturbations. On the other hand, we prove that this stability is genuinely symmetry-dependent: once arbitrary non-characteristic perturbations are allowed, Pansu--Minkowski spheres become unstable. This dichotomy suggests that the correct sub-Riemannian analogue of the Euclidean Minkowski problem is not merely a question of identifying equality cases, but rather one of understanding which symmetry, regularity, and admissibility classes are compatible with a stable curvature theory.
%The paper is organized as follows. In Section~2 we recall the basic geometry of the Heisenberg group, the horizontal calculus on hypersurfaces, and the relevant notions of horizontal mean curvature and perimeter measure. Section~3 is devoted to rotationally invariant surfaces in \(\mathbb H^1\), to the reduction of curvature functionals to the profile curve, and to the proof that neither the Pansu sphere nor the Korányi sphere is an optimal configuration for the Minkowski- and Heintze--Karcher-type variational problems considered here. In Section~4 we establish the first and second variation formulas for horizontal mean-curvature-driven functionals and formulate the corresponding notions of constrained stationarity and stability. Section~5 contains the classification of rotationally invariant critical points and the construction of the Pansu--Minkowski spheres, including the identification of the optimal member of the family. Section~6 proves the stability and instability results, while Section~7 establishes local minimality under rotationally invariant perturbations. The appendices contain the Riemannian variation formulas and the Riemannian approximation argument leading to the sub-Riemannian formulas used throughout the paper.

Let $M$ be a Riemannian manifold and let $S\subseteq M$ be a smooth, embedded, closed, two-sided hypersurface. In this paper, we consider geometric functionals of the form
\begin{equation}
    \label{intro_riemfun}
\mathcal H_f(S)=\int_S f(H)\,d S,
\end{equation}
where $H$ denotes the mean curvature of $S$ and $f:\mathbb R\to\mathbb R$ is a smooth function. Our main purpose is to establish first and second variation formulas of $\mathcal H_f$ under arbitrary smooth variations. We refer to \Cref{maintheoremappendixA} for a summary of the main outcomes in this regard.

\medskip
Under additional constraints on the ambient manifold, on the class of variations or on the curvature functional, such variation formulas are well-known and appear in several different contexts. When $f\equiv 1$, \eqref{intro_riemfun} boils down to the area functional, for which variation formulas are available in great generality \cite{MR4676392}. Other natural choices of $f$ specialize \eqref{intro_riemfun} to relevant curvature energies, such as the \emph{Willmore energy}, the \emph{total mean curvature} and the \emph{total inverse mean curvature}. Variation formulas for these and more general curvature functionals have typically been addressed within very specific frameworks, for instance allowing for restricted classes of variations \cite{MR3962031,MR4927775}, or imposing additional geometric constraints on the ambient manifold \cite{MR3666332,MR3962031,MR4927775,MR341351}. However, to the best of our knowledge, a general formula simultaneously allowing an arbitrary mean curvature functional, arbitrary variations, and an arbitrary ambient manifold does not seem to be available. Our first purpose is to provide such a formula, both at first and (most importantly) at second order.

\medskip
Variation formulas are natural tools when addressing geometric variational problems. The first variation identifies the corresponding Euler-Lagrange equation, while the second variation provides the relevant stability information. Besides the obvious role played by such variational tools in connection with the isoperimetric problem \cite{MR86338,MR0731682,MR0917854,MR3921314}, more general mean curvature functionals arise e.g. in connection with the establishment of sharp geometric inequalities of Minkowski \cite{MR3544938,MR2522433,MR1511220}, Willmore \cite{MR4037467,MR4941125,MR291994,MR278240}, and Heintze-Karcher \cite{MR533065, MR1173047,MR996826,MR3702549} type. Accordingly, general variation formulas provide a unified framework for these problems, while allowing the resulting identities to be adapted to the specific geometry under consideration.

\medskip
The derivation of the variation formulas, inspired by the approach proposed in \cite{MR4676392} and carried out in \Cref{appendiceriem}, is based on a direct analysis of the pointwise evolution of the relevant geometric quantities. After introducing the basic preliminaries about hypersurfaces and variations, we establish a computational lemma for iterated covariant derivatives, which provides the main technical tool for the subsequent calculations. We then describe how the basic geometric quantities associated with the evolving hypersurface change along the variation. More precisely, we derive first and second order formulas for the induced Jacobian, the unit normal, the second fundamental form, and the mean curvature. While the corresponding first-order identities are well-known, their second-order counterparts, in the general setting considered here, constitute the main computational step of the argument. These pointwise identities are then combined with the area formula and integrated over the reference hypersurface, leading to the first and second variation formulas for \eqref{intro_riemfun}. Throughout the argument, no assumptions are imposed on the curvature of the ambient manifold, and the variation is allowed to have both tangential and normal components.

\medskip
As a second main result, we exploit the Riemannian machinery to establish corresponding \emph{sub-Riemannian} variation formulas. Namely, we focus on the sub-Riemannian  \emph{Heisenberg group} $\mathbb H^n$, a suitable stratified Lie group which provides the prototypical model within sub-Riemannian \cite{MR3971262,MR2363343} and \emph{pseudohermitian} \cite{MR2165405} geometries. In this framework, a general mean curvature functional takes the form 
\begin{equation}\label{tmcfunsubriemderfgt5ryhythyh}
    \tmc^\hhh_f(S)=\int_S f\left(H^\hhh\right)\,d\sigma^\hhh,
\end{equation} 
where $S$ is an embedded, closed hypersurface, $H^\hhh$ is its \emph{horizontal mean curvature} and $\sigma^\hhh$ is its \emph{sub-Riemannian surface measure}. We refer to \Cref{appendicesubriem} for the related definitions, and to \Cref{teoremavariazionigeneralimainsubriem} for the statement of the variation formulas.

\medskip
Typically, sub-Riemannian geometry is determined by a distinguished distribution of admissible directions, the \emph{horizontal distribution} $\hhh$, endowed with a metric defined \emph{a priori} only along such directions. Accordingly, both intrinsic and extrinsic geometric quantities are built from such horizontal geometry. The horizontal geometry of a hypersurface may collapse when the horizontal distribution coincides with its tangent space, i.e. in the presence of \emph{characteristic points}. Nevertheless, a sub-Riemannian structure can be \emph{a posteriori} approximated by a suitable family of Riemannian metrics whose associated geometries converge, in a suitable sense, to the sub-Riemannian one.
Our approach to the sub-Riemannian problem is indeed inspired by the above construction. Namely, following a well-established approach \cite{MR2774306,MR2262784,MR4761954,MR2043961,MR5029815}, we approximate the sub-Riemannian structure $\left(\hh^n,\hhh,\langle\cdot,\cdot\rangle\right)$ by means of a sequence of collapsing Riemannian structures $$\big(\hh^n,\langle\cdot,\cdot\rangle_\eps\big),\qquad \eps>0.$$
For each approximating metric, the Riemannian variation formulas apply directly. The sub-Riemannian formulas are then obtained by a careful limiting argument. To this aim, we restrict to consider variations which do not move characteristic points. This procedure requires a detailed analysis of the convergence of the geometric quantities entering the Riemannian formulas. We therefore study both ambient objects, such as the curvature tensors of the approximating metrics, and quantities associated with the hypersurface, in particular those built from its second fundamental form. A notable feature of the approximation is that, although some of these quantities may diverge individually as $\varepsilon\to 0$, the specific combinations in which they enter the variation formulas remain convergent through appropriate cancellations. This phenomenon is well known in the Riemannian approximation of first and second variations of area \cite{MR2723818,MR5029815}. Here we show that it persists for the more involved quantities arising from the second variation of the mean curvature.
\section{Variation formulas for Riemannian mean curvature functionals}\label{appendiceriem}
In this section, we establish first and second variation formulas for \eqref{intro_riemfun}. Much of the notation and several results are borrowed from \cite{MR4676392}. We refer to \cite{MR1138207} for a general account of Riemannian geometry.
\subsection{Preliminaries} 
Here and hereafter, $M$ is a fixed $(n+1)$-dimensional Riemannian manifold for some $n\geq 1$, $\langle\cdot,\cdot\rangle$ is its Riemannian metric and $\nabla$ is its Levi-Civita connection.
In the following, Einstein's summation convention is assumed.
Denote by $\R$ both the $(3,1)$ and the $(4,0)$ Riemann tensor, namely
\begin{equation*}
    \R(\A ,\B )\C =\nabla_{\A}\nabla_{\B} \C-\nabla_{\B}\nabla_{\A} \C-\nabla_{[ \A,\B ]}\C,\quad \R(\A,\B,\C,\D)=\left\langle \R(\A,\B)\C,\D\right\rangle,\quad \A,\B,\C,\D\in\Gamma(TM).
\end{equation*}
If $\A,\C\in\Gamma(TM)$ are fixed, denote by $\j$ the $(2,0)$-tensor field defined by 
\begin{equation*}
   \j(\A ,\C)(\B,\D)\coloneqq\R(\A ,\B ,\C,\D),\qquad \B,\D\in\Gamma(TM).
\end{equation*}
%In this way,
%\begin{equation*}
 %   \ric(\A ,\C)=-\trace\j(\A ,\C ),\qquad  \A,\C \in\Gamma(TM).
%\end{equation*}
We recall that, if $ \A,\B ,\C,\D \in\Gamma(TM)$, then
    \begin{align*}
        \left(\nabla_{\A}\ric\right)(\B,\C )=-\trace\left(\nabla_{\A} \R\right)(\B,\cdot,\C ,\cdot),\qquad
        \left(\nabla_{\A}\R\right)(\B,\C,\D,\D)=0.
    \end{align*}
%\begin{proof}
 %   Let $\E_1,\ldots,\E_{n+1}$ be a local orthonormal frame. The thesis follows since
 %   \begin{equation*}
  %      \begin{split}
    %        \sum_{i=1}^{n+1}R\left(\B,\nabla_{\A}\E_i,\C ,\E_i\right)&+\sum_{i=1}^{n+1}R\left(\B,\E_i,\C ,\nabla_{\A}\E_i\right)\\
    %        &=\sum_{i,j=1}^{n+1}\left\langle \nabla_{\A}  \E_i , \E_j \right\rangle R\left(\B, \E_j ,\C , \E_i \right)+\sum_{i,j=1}^{n+1}\left\langle \nabla_{\A}  \E_i , \E_j \right\rangle R\left(\B, \E_i ,\C , \E_j \right)\\
    %        &=\sum_{i,j=1}^{n+1}\left(\left\langle \nabla_{\A}  \E_i , \E_j \right\rangle +\left\langle \nabla_{\A}  \E_j , \E_j \right\rangle\right)R\left(\B, \E_j ,\C , \E_i \right)\\
     %       &=\sum_{i,j=1}^{n+1} A\left\langle   \E_i , \E_j \right\rangle R\left(\B, \E_j ,\C , \E_i \right)\\
     %       &=0.
    %    \end{split}
   % \end{equation*}
%\end{proof}
%\begin{lemma}
%    Let $ \A,\B ,\C ,\D\in\Gamma(TM)$. Then
   % \begin{equation*}
  %      \left(\nabla_{\A}\R\right)(\B,\C,\D,\D)=0.
  %  \end{equation*}
%\end{lemma}
%\begin{proof}
  %  By a direct computation,
   % \begin{equation*}
   %     \begin{split}
    %        \left(\nabla_{\A}\R\right)(\B,\C,\D,\D)&=\A\left(R\left(\B,\C,\D,\D\right)\right)-R\left(\nabla_\A \B,\C,\D,\D\right)-R\left( \B,\nabla_\A\C,\D,\D\right)\\
   %         &\quad-R\left( \B,\C,\nabla_\A\D,\D\right)-R\left( \B,\C,\D,\nabla_\A\D\right)\\
    %        &=0.
   %     \end{split}
  %  \end{equation*}
%\end{proof}
%\begin{lemma}
 %   Let \footnote{Qesto lemma per ora non lo ho usato, ma secondo me serve a semplificare la parte non normale della variazione seconda di H. Solo che ora mi sto concentrando su quella normale.}$\A,\B\in\Gamma(TM)$. Then
 %   \begin{equation*}
   %     \left\langle \nabla \A,\left(\nabla\B\right)^t\right\rangle+\ric(\A,\B)=\divv \left(\nabla_\A \B\right)-\A\divv \B.
   % \end{equation*}
%\end{lemma}
%\begin{proof}
 %   Fix $p\in M$ and a geodesic frame $\E_1,\ldots,\E_{n+1}$ at $p$. Then
 %   \begin{equation*}
  %      \begin{split}
   %         \left\langle \nabla \A,\left(\nabla\B\right)^t\right\rangle&+\ric(\A,\B)=\sum_{i,j=1}^{n+1}\left\langle\nabla_{\E_i}\A,\E_j\right\rangle \left\langle\nabla_{\E_j}\B,\E_i\right\rangle+\ric(\A,\B)\\
     %       &\quad=\sum_{i,j=1}^{n+1}\E_i\left\langle\A,\E_j\right\rangle \left\langle\nabla_{\E_j}\B,\E_i\right\rangle+\ric(\A,\B)\\
     %         &\quad=\sum_{i,j=1}^{n+1}\E_i\left(\left\langle\A,\E_j\right\rangle \left\langle\nabla_{\E_j}\B,\E_i\right\rangle\right)-\sum_{i,j=1}^{n+1}\left\langle\A,\E_j\right\rangle \E_i\left\langle\nabla_{\E_j}\B,\E_i\right\rangle+\ric(\A,\B)\\
      %        &\quad=\sum_{i=1}^{n+1}\E_i\left\langle\nabla_{\A}\B,\E_i\right\rangle-\sum_{i,j=1}^{n+1}\left\langle\A,\E_j\right\rangle \left\langle\nabla_{\E_i}\nabla_{\E_j}\B,\E_i\right\rangle+\ric(\A,\B)\\
         %     &\quad=\divv \left(\nabla_\A \B\right)-\sum_{i,j=1}^{n+1}\left\langle\A,\E_j\right\rangle \left\langle\nabla_{\E_j}\nabla_{\E_i}\B,\E_i\right\rangle-\sum_{i,j=1}^{n+1}\left\langle\A,\E_j\right\rangle \R(\E_i,\E_j,\B,\E_j)+\ric(\A,\B)\\
        %      &\quad=\divv \left(\nabla_\A \B\right)-\A\divv \B.
      %  \end{split}
  %  \end{equation*}
%\end{proof}
 Let $p\in M$. Let $x_1,\ldots,x_{n+1}$ be local coordinates in $M$ near $p$.
 Write 
 \begin{equation*}
     \R\left(\frac{\partial }{\partial x_i},\frac{\partial}{\partial x_j}\right)\frac{\partial }{\partial x_k}=\R_{ijk}^l\frac{\partial}{\partial x_l},\qquad  \left(\nabla_{\frac{\partial}{\partial x_\alpha}}\R\right)\left(\frac{\partial }{\partial x_i},\frac{\partial}{\partial x_j}\right)\frac{\partial }{\partial x_k}=\left(\nabla \R\right)_{\alpha i j k}^l\frac{\partial}{\partial x_l},\qquad\alpha,i,j,k=1,\ldots,n+1.
 \end{equation*}
 Notice that 
 \begin{equation}\label{riemanincoordinate}
     \R_{ijk}^l=\frac{\partial\Gamma_{jk}^l}{\partial x_i}-\frac{\partial \Gamma_{ik}^l}{\partial x_j}+\Gamma_{jk}^m\Gamma_{im}^l-\Gamma_{ik}^m\Gamma_{jm}^l,\qquad i,j,k,l=1,\ldots,n+1,
 \end{equation}
where $\Gamma^k_{ij}$ are the Christoffel symbols. In particular, when $x_1,\ldots,x_{n+1}$ are normal coordinates centered at $p$,
\begin{equation}\label{normcoordcristo}
   \Gamma_{ij}^{k}(p)=0\qquad\text{for any $i,j,k=1,\ldots,n+1$,} 
\end{equation}
 so that 
\begin{equation}\label{rieminnorm}
    \R_{ijk}^l(p)=\frac{\partial \Gamma_{jk}^l}{\partial x_i}(p)-\frac{\partial \Gamma_{ik}^l}{\partial x_j}(p)\qquad\text{for any $i,j,k,l=1,\ldots,n+1$}
\end{equation}
and 
\begin{equation}\label{deririeminnorm}
    \left(\nabla \R\right)_{\alpha ijk}^l(p)=\frac{\partial ^2\Gamma_{jk}^l}{\partial x_\alpha\partial x_i}(p)-\frac{\partial^2 \Gamma_{ik}^l}{\partial x_\alpha\partial x_j}(p)\qquad\text{for any $\alpha,i,j,k,l=1,\ldots,n+1$}.
\end{equation}
    \subsection{Hypersurfaces}
    Throughout this section, $S\subseteq M$ is a smooth, embedded, closed, two-sided hypersurface. Denote by $\n$ a (globally defined) unit normal to $S$. If $\A\in\Gamma(TM)$, denote by $\A^T$ its orthogonal projection onto $TS$. Denote by $A$ and $h$ the shape operator and the second fundamental form, i.e.
    \begin{equation*}
        A\left(\A\right)=\nabla_\A\n,\qquad h(\A,\B)=\left\langle A(\A),\B\right\rangle,\qquad\A,\B\in\Gamma(TS).    \end{equation*}
    Denote by $H$ the (non-averaged) mean curvature of $S$, and by $\nabla^S$ the Levi-Civita connection of $S$. If $B$ is a $(2,0)$-tensor field on $S$, set   \begin{equation*}
       B^t(\A,\B)\coloneqq B(\B,\A),\qquad \langle L_B(\A),\B\rangle=B(\A,\B),\qquad \A,\B\in\Gamma(TS).
    \end{equation*} 
    For any $k\in\mathbb N_+$, define 
\begin{equation*}
    B^k(\A,\B)=\left\langle L_B^k(\A),\B\right\rangle,\qquad \A,\B\in \Gamma(TS).
\end{equation*}
If $B$ is symmetric, then $B^k$ is symmetric. If $p\in S$ and $e_1,\ldots,e_n$ is an orthonormal basis of $T_p S$, we write
\begin{equation*}\label{atok}
   \left(B^k\right)_{ij}\coloneqq B^k(e_i,e_j)=\sum_{l_1,\ldots,l_{k-1}=1}^nB_{il_1}B_{l_!l_2}\cdots B_{l_{k-1}l_{k-1}}B_{l_{k-1}j},\qquad i,j=1,\ldots, n.
    \end{equation*}
  %  Accordingly, if $C$ is another symmetric, $(2,0)$-tensor field on $S$, we set
  %  \begin{equation*}
   %     \left(B\circ C\right)(\A,\B)\coloneqq\left\langle L_B(\A),L_C(\B)\right\rangle,\qquad \A,\B\in\Gamma(TS).
  %  \end{equation*}
  %  In this way,
  %  \begin{equation}\label{compoduetensori}
   %     \left(B\circ C\right)_{ij}=\sum_{k=1}^nB_{ik}C_{kj},\qquad i,j=1,\ldots,n.
   % \end{equation}
  %  Notice that 
  %  \begin{equation*}
  %     \trace\left(B\circ C\right)=\trace\left(C\circ B\right).
 %   \end{equation*}
%Denote by $A: \Gamma(TS)\to \Gamma(TS)$ the shape operator associated with $S$. For any $k\in\mathbb N_+$, define 
%\begin{equation*}
 %   h^k(X,Y)=\left\langle A^k(X),Y\right\rangle,\qquad X,Y\in \Gamma(TS).
%\end{equation*}
%Notice that $h^k$ is symmetric. Moreover, if $p\in S$ and $e_1,\ldots,e_n$ is an orthonormal basis of $T_p S$, we write
%\begin{equation*}
%    \left(h^k\right)_{ij}=h^k(e_i,e_j),\ldots i,j=1,\ldots,n. 
%\end{equation*}
%In particular,
%\begin{equation}\label{atok}
 %   \left( h^2\right)_{ij}=\sum_{k=1}^nh_{ik}h_{jk},\qquad  \left( h^3\right)_{ij}=\sum_{k,l=1}^nh_{ik}h_{kl}h_{lj},\ldots i,j=1,\ldots,n.
  %  \end{equation}
 %   More generally, if $B$ is a $(2,0)$ symmetric tensor field on $S$, and defining $\tilde B C:\Gamma(TS)\to\Gamma(TS)$ by 
  %  \begin{equation*}
 %       \langle \tilde B(X),Y\rangle=B(X,Y),\qquad X,Y\in\Gamma(TS),
 %   \end{equation*}
 %   we define the $(2,0)$-tensor field $B^k$ by
  %  \begin{equation*}
 %       B^k(X,Y)=\left\langle \tilde B^k\left(X\right),Y\right\rangle,\qquad X,Y\in\Gamma(TS).
%    \end{equation*}
  %  Notice that 
  %  \begin{equation*}
  %     \trace\left(B\circ C\right)=\trace\left(C\circ B\right).
 %   \end{equation*}
 If $\A,\B,\C\in\Gamma(TS)$,
 %the \emph{Codazzi equation} and
 the \emph{traced Codazzi equation} reads as  
     \begin{align}
      %      \left(\nabla^S_\B h\right)(\A,\C)&=\left(\nabla^S_\A h\right)(\B,\C)+\R(\A,\B,\C,\n),  \nonumber  \\  
            \left(\divv^S h\right)(\A)&=\A H+\ric(\A,\n)\label{tracedcodazzi}
        \end{align}
        We will make use of the following consequence of the traced Codazzi equation \eqref{tracedcodazzi}.
%    \begin{proposition}[Codazzi equations]
    %    Let $\A,\B,\C\in \Gamma(TS)$. Then 
   %     \begin{align}
   %         \left(\nabla^S_\B h\right)(\A,\C)&=\left(\nabla^S_\A h\right)(\B,\C)+\R(\A,\B,\C,\n),  \nonumber  \\  
     %       \left(\divv^S h\right)(\A)&=\A H+\ric(\A,\n)\label{tracedcodazzi}.
     %   \end{align}
  %  \end{proposition}
    \begin{lemma}
        Let $\A\in\Gamma(TM)$. Then
        \begin{equation}\label{corollarioditracedcodazzi}
            \left\langle h,\nabla^S\A\right\rangle+\ric(\A,\n)=\divv^S A\left(\A^T\right)-\A^T H+\left\langle\A,\n\right\rangle\left(|h|^2+\ric(\n,\n)\right).
        \end{equation}
            \end{lemma}
        \begin{proof}
            Let $p\in S$. Let $\E_1,\ldots,\E_n$ be a geodesic frame of $S$ at $p$. Then
      \begin{equation*}              \begin{split}                 \left\langle h,\nabla^S\A\right\rangle
      %&=\sum_{i,j=1}^n h(\E_i,\E_j)\left\langle\nabla_{\E_i}\A,\E_j\right\rangle\\
      &=\sum_{i,j=1}^n h(\E_i,\E_j)\E_i\left\langle\A,\E_j\right\rangle-\sum_{i,j=1}^n h(\E_i,\E_j)\left\langle\A,\nabla_{\E_i}\E_j\right\rangle\\
      &=\sum_{i,j=1}^n \E_i\left(h(\E_i,\E_j)\left\langle\A,\E_j\right\rangle\right)-\sum_{i,j=1}^n\E_i\left(h(\E_i,\E_j)\right)\left\langle\A,\E_j\right\rangle-\left\langle\A,\n\right\rangle\sum_{i,j=1}^n h(\E_i,\E_j)\left\langle\n,\nabla_{\E_i}\E_j\right\rangle\\
      &=\sum_{i=1}^n \E_i\left(h\left(\E_i,\A^T\right)\right)-\sum_{i,j=1}^n\left(\nabla^S_{\E_i}h\right)(\E_i,\E_j)\left\langle\A,\E_j\right\rangle+\left\langle\A,\n\right\rangle|h|^2\\
      &=\divv^S A\left(\A^T\right)-\left(\divv^S h\right)\left(\A^T\right)+\left\langle\A,\n\right\rangle|h|^2.
             \end{split}
      \end{equation*}  
      Therefore
      \begin{equation*}
          \begin{split}
               \left\langle h,\nabla^S\A\right\rangle+\ric(\A,\n)&= \divv^S A\left(\A^T\right)-\left(\divv^S h\right)\left(\A^T\right)+\ric\left(\A^T,\n\right)+\left\langle\A,\n\right\rangle\left(|h|^2+\ric(\n,\n)\right)\\
               \overset{\eqref{tracedcodazzi}}&{=}\divv^S A\left(\A^T\right)-\A^T H+\left\langle\A,\n\right\rangle\left(|h|^2+\ric(\n,\n)\right).
          \end{split}
      \end{equation*}
      \end{proof}
  \subsection{Variations}\label{subsec_variations}
    A variation is a smooth map $\Phi:I\times M\to M$, where $I\subseteq\rr$ is any open neighborhood of $0$, such that:
\begin{itemize}
    \item $p\mapsto \Phi(t,p)$ is a diffeomorphism for any $t\in I$;
    \item $\Phi(0,p)=p$ for any $p\in M$. 
\end{itemize}
We adopt the notation $\Phi_t(p)\coloneqq\Phi(t,p)$. A variation is \emph{compactly supported} if $\Phi_t(p)=p$ outside a compact set $K(\Phi)$. In the following, we tacitly assume that a variation is compactly supported. This is clearly not restrictive when computing variations of a closed hypersurface. Nevertheless, closedness is assumed just for the sake of simplicity: the forthcoming variation formulas would clearly hold as well provided the variation is supported far from the boundary of the reference hypersurface. Define the time-dependent vector field $\Y$ by
\begin{equation}\label{def_xechis}
    \Y(t,q)=\left.\frac{\partial}{\partial s}\right|_{s=0}\left(\Phi_{s+t}\circ\Phi_t^{-1}\right)(q)\qquad\text{for any $t\in I$ and $q\in M$.}
\end{equation}
%Then, the \emph{velocity} $\X$ and the \emph{acceleration} $\X'$ of $\Phi$ are defined respectively by 
Define
\begin{equation*}
    \X(q)=\Y(0,q),\qquad \X'(q)=\left.\frac{\partial }{\partial t}\right|_{t=0}\Y(t,q),\qquad\Z(q)=\left(\X'+\nabla_\X\X\right)(q)\qquad\text{for any $q\in M$.}
\end{equation*}
The vector fields $\X$ and $\Z$ are known as \emph{variational velocity field} and \emph{variational acceleration field} of $\Phi$. We point out that, while the velocity depends only on the underlying differential structure, the acceleration does depend on the metric via the connection term $\nabla_\X\X$. 
%Given any couple of vector fields $\X,\Z$ it is always possible to construct a variation having $\X$ and $\Z$ respectively as velocity and acceleration. 
%a Typical instance is the exponential variation
%\begin{equation}\label{exponentialvariation}
%    \Phi_t(p)\coloneqq\exp_p\left(t\X(p)+\frac{t^2}{2}\Z(p)\right)\footnote{Controllare, non mi è chiaro se sia la formula giusta}.
%\end{equation}
Fix $\varphi\in C^\infty(S)$ and $ \X^T\in\Gamma(TS)$ such that 
\begin{equation*}
    \X|_S=\ \X^T+\varphi \n.
\end{equation*}
When $X^T\equiv 0$, $\Phi$ is called \emph{normal variation}. 
     Set $S_t\coloneqq \Phi_t(S)$ for any $t\in I.$ If $p\in S$ is fixed, set $\beta_p(t)\coloneqq \Phi_t(p)$ for every $t\in I$. If there is no ambiguity, we write $\beta=\beta_p$. By definition, $\beta(t)\in S_t$ for every $t\in I$. Moreover,

\begin{equation*}\label{xtempodefinisceflusso}
    \begin{split}
        \Y(t,\beta(t))=\Y(t,\Phi_t(p))=\left.\frac{\partial}{\partial s}\right|_{s=0}\Phi_{s+t}(p)=\left.\frac{\partial}{\partial s}\right|_{s=t}\Phi_{s}(p)=\dot\beta(t)\qquad\text{for any $t\in I$}.
    \end{split}
\end{equation*}
Therefore, $\Phi$ is the (time-dependent) flow of $\Y$. 
For any $t\in I$, denote by $\N(t,\cdot)$ a smooth choice of arbitrary local extensions, around $S_t$, of unit normals to $S_t$, in such a way that $\N(0,q)=\n(q)$ locally around $S$. Moreover, if $p\in S$, denote by $H_t\left(\Phi_t(p)\right)$ the mean curvature of $S_t$ at $\Phi_t(p)$. When $p\in S$ is fixed and local coordinates $x_1,\ldots,x_{n+1}$ around $p$ are given, we will write 
  \begin{equation*}
        \Y(t,q)=a^i(t,q)\left.\frac{\partial}{\partial x_i}\right|_{q},\qquad \N(t,q)=c^k(t,q)\left.\frac{\partial}{\partial x_k}\right|_{q}
    \end{equation*}
    and \begin{equation*}
        \X(q)=X^i(q)\left.\frac{\partial}{\partial x_i}\right|_{q},\qquad \X'(q)=\left(X'\right)^i(q)\left.\frac{\partial}{\partial x_i}\right|_{q},\qquad \n(q)=N^i(q)\left.\frac{\partial}{\partial x_i}\right|_{q}.
    \end{equation*}
    Notice that, by definition,
    \begin{equation*}
        X^i(q)=a^i(0,q),\qquad(X')^i(q)=\left.\frac{\partial }{\partial t}\right|_{t=0}a^i(t,q),\qquad N^i(q)=c^i(0,q).
    \end{equation*}
    We denote by $\frac{D}{dt}$ the covariant derivative operator, along $\beta$, induced by $\nabla$. We recall that
    %, if $F(t)=\left.f^j(t)\frac{\partial}{\partial x_j}\right|_{\beta(t)}$ locally around $p$, then
\begin{equation}\label{comefarecovdev}
    \frac{D}{dt}\left(\left.f^j(t)\frac{\partial}{\partial x_j}\right|_{\beta(t)}\right)=\dot f^i(t)\left.\frac{\partial}{\partial x_j}\right|_{\beta(t)}+a^i(t,\beta(t))f^j(t)\Gamma_{ij}^k(\beta(t))\left.\frac{\partial}{\partial x_k}\right|_{\beta(t)}.
\end{equation}
%In particular,
%\begin{equation*}
 %   \Z(p)=\left.\frac{D}{dt}\right|_{t=0}\Y(t,\beta_p(t)).
%\end{equation*}
  Finally, when $p\in S$ and $e\in T_p S$ are fixed, we set  
    \begin{equation}\label{estensionevettore}
        E(t)\coloneqq\left(d\Phi_t\right)|_p(e)\qquad\text{for every $t\in I$.} 
    \end{equation}
    $E$ is a vector field along $\beta$, and $E(t)\in T_{\beta(t)}S_t$ for every $t\in I$. In local coordinates around $p$, we write
      \begin{equation*}
        E(t)=b^j(t)\left.\frac{\partial}{\partial x_j}\right|_{\beta(t)}
    \end{equation*}
    Fix $p\in S$. Let $e_1,\ldots,e_n$ be an orthonormal basis of $T_p S$. The \emph{Jacobian} of $\Phi_t|_{S}$ at $p$ is defined by
    \begin{equation*}
        \jac\left(\Phi_t|_{S}\right)\left(p\right)\coloneqq\sqrt{\det \G(t)}
    \end{equation*}
    where the symmetric matrix $\G$ is defined by
    \begin{equation}\label{grammatrix}
         \G(t)_{ij}\coloneqq\langle E_i(t),E_j(t)\rangle,\qquad i,j=1,\ldots,n.
    \end{equation}
    If $p\in S$ and we fix local coordinates $x_1,\ldots, x_{n+1}$ in a neighborhood of $p$, say $U$, the continuity of $\Phi$ ensures that $\Phi_t(q)\in U$ for any $t$ small and any $q$ sufficiently close to $p$. We will tacitly assume this. 
\subsection{A computational lemma}
The following technical lemma is the main computational tool used below.
\begin{lemma}\label{fondamentalecomplemmavariazioniriemannianejhrnfi4rfjnfin}
    Let $p\in M$. Let $x_1,\ldots,x_{n+1}$ be local coordinates in $M$ near $p$.  Let $ \A,\B ,\C ,\D\in\Gamma(TM)$. Then, near $p$,
     \begin{equation}\label{nablatrecosenonnormale}
    \begin{split}
       \nabla_{\A}\nabla_{\B} \C
        &=A^\alpha B^\beta\left(\frac{\partial^2 C^\gamma}{\partial x_\alpha\partial x_\beta}\frac{\partial}{\partial x_\gamma}+\frac{\partial C^\gamma}{\partial x_\alpha}\Gamma_{\beta\gamma}^\delta\frac{\partial}{\partial x_\delta}+C^\gamma\frac{\partial \Gamma_{\alpha\gamma}^\delta}{\partial x_\beta}\frac{\partial}{\partial x_\delta}+\frac{\partial C^\gamma}{\partial x_\beta}\Gamma_{\alpha\gamma}^\delta\frac{\partial}{\partial x_\delta}+C^\gamma\Gamma_{\alpha\gamma}^\delta\Gamma_{\beta\delta}^\eta\frac{\partial}{\partial x_\eta}\right)\\
        &\quad+A^\alpha\frac{\partial B^\beta}{\partial x_\alpha}\left(\frac{\partial C^\gamma}{\partial x_\beta}\frac{\partial}{\partial x_\gamma}+C^\gamma\Gamma_{\beta\gamma}^\delta\frac{\partial}{\partial x_\delta}\right)+\R(\A ,\B )\C .
    \end{split}
\end{equation}
Assume in addition that $x_1,\ldots,x_{n+1}$ are normal coordinates centered at $p$. Then, at $p$, 
\begin{equation}\label{nablatrecose}
        \begin{split}
            \nabla_{\A}\nabla_{\B}\C&=\R(\A ,\B )\C +A^\alpha B^\beta C^\gamma\frac{\partial \Gamma_{\alpha\gamma}^\delta}{\partial x_\beta} \frac{\partial}{\partial x_\delta}+A^\alpha B^\beta\frac{\partial^2 C^\gamma}{\partial x_\alpha \partial x_\beta} \frac{\partial}{\partial x_\gamma}+A^\alpha\frac{\partial B^\beta}{\partial x_\alpha}\frac{\partial C^\gamma}{\partial x_\beta}\frac{\partial}{\partial x_\gamma}
        \end{split}
    \end{equation}
    and
    \begin{equation}\label{nablaquattrocose}
    \begin{split}          &\nabla_{\A}\nabla_{\B}\nabla_{\C}\D=\left(\nabla_{\B} \R\right)(\A ,\C )\D+\R(\A ,\B )\nabla_{\C} \D+\R\left(\A ,\C \right)\nabla_{\B} \D+\R\left(\A ,\nabla_{\B} \C\right)\D+\R\left(\nabla_\A \B,\C\right)\D\\
          &\quad+A^\alpha B^\beta\left(\frac{\partial ^2C^\delta}{\partial x_\alpha\partial x_\beta}\frac{\partial D^\gamma}{\partial x_\delta} +\frac{\partial C^\delta}{\partial x_\beta}\frac{\partial^2 D^\gamma}{\partial x_\alpha\partial x_\delta}+\frac{\partial C^\delta}{\partial x_\alpha}\frac{\partial^2 D^\gamma}{\partial x_\beta\partial x_\delta} +  C^\delta\frac{\partial^3 D^\gamma}{\partial x_\alpha\partial x_\beta\partial x_\delta} \right)\frac{\partial}{\partial x_\gamma}\\
          &\quad+A^\alpha B^\beta\left(C^\eta\frac{\partial D^\delta}{\partial x_\eta}\frac{\partial \Gamma_{\alpha\delta}^\gamma}{\partial x_\beta} +\frac{\partial C^\delta  }{\partial x_\beta}D^\eta\frac{\partial \Gamma^\gamma_{\alpha\eta}}{\partial x_\delta} + C^\delta\frac{\partial D^\eta}{\partial x_\beta}\frac{\partial \Gamma^\gamma_{\alpha\eta}}{\partial x_\delta}+\frac{\partial C^\delta }{\partial x_\alpha}D^\eta\frac{\partial \Gamma^\gamma_{\delta \eta}}{\partial x_\beta} +C^\delta \frac{\partial D^\eta}{\partial x_\alpha}\frac{\partial \Gamma^\gamma_{\delta \eta}}{\partial x_\beta} \right)\frac{\partial}{\partial x_\gamma}\\
          &\quad+A^\alpha B^\beta C^\delta D^\eta\frac{\partial^2 \Gamma^\gamma_{\alpha \eta}}{\partial x_\beta\partial x_\delta} \frac{\partial}{\partial x_\gamma}+A^\alpha\frac{\partial B^\beta}{\partial x_\alpha}\left(\frac{\partial C^\delta}{\partial x_\beta}\frac{\partial D^\gamma}{\partial x_\delta}+ C^\delta\frac{\partial^2 D^\gamma}{\partial x_\beta\partial x_\delta}+ C^\delta D^\eta\frac{\partial \Gamma^\gamma_{\beta \eta}}{\partial x_\delta}\right)\frac{\partial}{\partial x_\gamma}.
        \end{split}
    \end{equation}
\end{lemma}
\begin{proof}
    By a direct computation, near $p$,
    \begin{equation*}
    \begin{split}
       \nabla_{\A}\nabla_{\B} \C&=A^\alpha\nabla_{\frac{\partial}{\partial x_\alpha}}\left(B^\beta\nabla_{\frac{\partial}{\partial x_\beta}} \left(C^\gamma\frac{\partial}{\partial x_\gamma}\right)\right)\\
        &=A^\alpha B^\beta\nabla_{\frac{\partial}{\partial x_\alpha}}\left(\nabla_{\frac{\partial}{\partial x_\beta}} \left(C^\gamma\frac{\partial}{\partial x_\gamma}\right)\right)+A^\alpha\frac{\partial B^\beta}{\partial x_\alpha}\nabla_{\frac{\partial}{\partial x_\beta}} \left(C^\gamma\frac{\partial}{\partial x_\gamma}\right)\\
        &=A^\alpha B^\beta \nabla_{\frac{\partial}{\partial x_\alpha}}\left(\frac{\partial C^\gamma}{\partial x_\beta} \frac{\partial}{\partial x_\gamma}+C^\gamma\Gamma_{\beta\gamma}^\delta \frac{\partial}{\partial x_\delta}\right)+A^\alpha\frac{\partial B^\beta}{\partial x_\alpha}\left(\frac{\partial C^\gamma}{\partial x_\beta}\frac{\partial}{\partial x_\gamma}+C^\gamma\Gamma_{\beta\gamma}^\delta\frac{\partial}{\partial x_\delta}\right)\\
        &=A^\alpha B^\beta\left(\frac{\partial^2 C^\gamma}{\partial x_\alpha\partial x_\beta}\frac{\partial}{\partial x_\gamma}+\frac{\partial C^\gamma}{\partial x_\alpha}\Gamma_{\beta\gamma}^\delta\frac{\partial}{\partial x_\delta}+C^\gamma\frac{\partial \Gamma_{\beta\gamma}^\delta}{\partial x_\alpha}\frac{\partial}{\partial x_\delta}+\frac{\partial C^\gamma}{\partial x_\beta}\Gamma_{\alpha\gamma}^\delta\frac{\partial}{\partial x_\delta}+C^\gamma\Gamma_{\beta\gamma}^\delta\Gamma_{\alpha\delta}^\eta\frac{\partial}{\partial x_\eta}\right)\\
        &\quad+A^\alpha\frac{\partial B^\beta}{\partial x_\alpha}\left(\frac{\partial C^\gamma}{\partial x_\beta}\frac{\partial}{\partial x_\gamma}+C^\gamma\Gamma_{\beta\gamma}^\delta\frac{\partial}{\partial x_\delta}\right).
    \end{split}
\end{equation*}
Since
\begin{equation*}
\begin{split}
     A^\alpha B^\beta C^\gamma\left(\frac{\partial\Gamma_{\beta\gamma}^\delta}{\partial x_\alpha}\frac{\partial}{\partial x_\delta}+\Gamma_{\beta\gamma}^\delta\Gamma_{\alpha\delta}^\eta\frac{\partial}{\partial x_\eta}\right)&= A^\alpha B^\beta C^\gamma\left(\frac{\partial\Gamma_{\beta\gamma}^\delta}{\partial x_\alpha}+\Gamma_{\beta\gamma}^\eta\Gamma_{\alpha\eta}^\delta\right)\frac{\partial}{\partial x_\delta}\\
     &\overset{\eqref{riemanincoordinate}}{=}\R(\A ,\B )\C +A^\alpha B^\beta C^\gamma\left(\frac{\partial\Gamma_{\alpha\gamma}^\delta}{\partial x_\beta}+\Gamma_{\alpha\gamma}^\eta\Gamma_{\beta\eta}^\delta\right)\frac{\partial}{\partial x_\delta},
\end{split}
\end{equation*}
then \eqref{nablatrecosenonnormale} follows. 
%\begin{equation*}
  %  \begin{split}
   %    \nabla_{\A}\nabla_{\B} C&=A^\alpha\nabla_{\frac{\partial}{\partial x_\alpha}}\left(B^\beta\nabla_{\frac{\partial}{\partial x_\beta}} \left(C^\gamma\frac{\partial}{\partial x_\gamma}\right)\right)\\
 %       &=A^\alpha B^\beta\nabla_{\frac{\partial}{\partial x_\alpha}}\left(\nabla_{\frac{\partial}{\partial x_\beta}} \left(C^\gamma\frac{\partial}{\partial x_\gamma}\right)\right)+A^\alpha\frac{\partial B^\beta}{\partial x_\alpha}\nabla_{\frac{\partial}{\partial x_\beta}} \left(C^\gamma\frac{\partial}{\partial x_\gamma}\right)\\
   %     &=A^\alpha B^\beta\nabla_{\frac{\partial}{\partial x_\alpha}}\left(C^\gamma\Gamma_{\beta\gamma}^\delta \frac{\partial}{\partial x_\delta}\right)+A^\alpha B^\beta\nabla_{\frac{\partial}{\partial x_\alpha}}\left(\frac{\partial C^\gamma}{\partial x_\beta} \frac{\partial}{\partial x_\gamma}\right)+A^\alpha\frac{\partial B^\beta}{\partial x_\alpha}\frac{\partial C^\gamma}{\partial x_\beta}\frac{\partial}{\partial x_\gamma}\\
   %     &=A^\alpha B^\beta C^\gamma\frac{\partial \Gamma_{\beta\gamma}^\delta}{\partial x_\alpha} \frac{\partial}{\partial x_\delta}+A^\alpha B^\beta\frac{\partial^2 C^\gamma}{\partial x_\alpha \partial x_\beta} \frac{\partial}{\partial x_\gamma}+A^\alpha\frac{\partial B^\beta}{\partial x_\alpha}\frac{\partial C^\gamma}{\partial x_\beta}\frac{\partial}{\partial x_\gamma}\\
   %     &=\R(\A ,\B )\C +A^\alpha B^\beta C^\gamma\frac{\partial \Gamma_{\alpha\gamma}^\delta}{\partial x_\beta} \frac{\partial}{\partial x_\delta}+A^\alpha B^\beta\frac{\partial^2 C^\gamma}{\partial x_\alpha \partial x_\beta} \frac{\partial}{\partial x_\gamma}+A^\alpha\frac{\partial B^\beta}{\partial x_\alpha}\frac{\partial C^\gamma}{\partial x_\beta}\frac{\partial}{\partial x_\gamma}.
  %      &=-\R(B,A)\C +A^\alpha B^\beta C^\gamma\frac{\partial \Gamma_{\alpha\gamma}^\delta}{\partial x_\beta} \frac{\partial}{\partial x_\delta}+A^\alpha B^\beta\frac{\partial^2 C^\gamma}{\partial x_\alpha \partial x_\beta} \frac{\partial}{\partial x_\gamma}+A^\alpha\frac{\partial B^\beta}{\partial x_\alpha}\frac{\partial C^\gamma}{\partial x_\beta}\frac{\partial}{\partial x_\gamma}\\
  %  \end{split}
%\end{equation*}
Moreover, \eqref{nablatrecose} follows by \eqref{nablatrecosenonnormale} and \eqref{normcoordcristo}. Finally, at $p$,
    \begin{equation*}
        \begin{split}
            \nabla_{\A}&\nabla_{\B}\nabla_{\C}\D=\nabla_{\A}\nabla_{\B}\left(\nabla_{\C}\D\right)\\
            &\overset{\eqref{nablatrecose}}{=}\R(\A ,\B )\nabla_{\C} \D+A^\alpha B^\beta \left(\nabla_{\C} \D\right)^\gamma\frac{\partial \Gamma_{\alpha\gamma}^\delta}{\partial x_\beta} \frac{\partial}{\partial x_\delta}+A^\alpha B^\beta\frac{\partial^2 \left(\nabla_{\C} \D\right)^\gamma}{\partial x_\alpha \partial x_\beta} \frac{\partial}{\partial x_\gamma}+A^\alpha\frac{\partial B^\beta}{\partial x_\alpha}\frac{\partial \left(\nabla_{\C} \D\right)^\gamma}{\partial x_\beta}\frac{\partial}{\partial x_\gamma}.
        \end{split}
    \end{equation*}
    Fix $\alpha,\beta,\gamma=1,\ldots,n+1$.
    Recall that
    \begin{equation*}
        \left(\nabla_{\C} \D\right)^\gamma=C^\delta\frac{\partial D^\gamma}{\partial x_\delta}+C^\delta D^\eta\Gamma^\gamma_{\delta\eta}.
    \end{equation*}
    Then, at $p$,
    \begin{equation*}
        \begin{split}
            \frac{\partial \left(\nabla_{\C} \D\right)^\gamma}{\partial x_\beta}&= \frac{\partial C^\delta}{\partial x_\beta}\frac{\partial D^\gamma}{\partial x_\delta}+C^\delta\frac{\partial^2 D^\gamma}{\partial x_\beta\partial x_\delta}+\frac{\partial C^\delta  }{\partial x_\beta}D^\eta\Gamma^\gamma_{\delta\eta}+C^\delta\frac{\partial D^\eta}{\partial x_\beta}\Gamma^\gamma_{\delta\eta}+C^\delta D^\eta\frac{\partial \Gamma^\gamma_{\delta \eta}}{\partial x_\beta}\\
            \overset{\eqref{normcoordcristo}}&{=}\frac{\partial C^\delta}{\partial x_\beta}\frac{\partial D^\gamma}{\partial x_\delta}+C^\delta\frac{\partial^2 D^\gamma}{\partial x_\beta\partial x_\delta}+C^\delta D^\eta\frac{\partial \Gamma^\gamma_{\delta \eta}}{\partial x_\beta},
        \end{split}
    \end{equation*}
    and moreover
    \begin{equation*}
        \begin{split}
              \frac{\partial^2 \left(\nabla_{\C} \D\right)^\gamma}{\partial x_\alpha\partial x_\beta}&\overset{\eqref{normcoordcristo}}{=}\frac{\partial ^2C^\delta}{\partial x_\alpha\partial x_\beta}\frac{\partial D^\gamma}{\partial x_\delta}+\frac{\partial C^\delta}{\partial x_\beta}\frac{\partial^2 D^\gamma}{\partial x_\alpha\partial x_\delta}+\frac{\partial C^\delta}{\partial x_\alpha}\frac{\partial^2 D^\gamma}{\partial x_\beta\partial x_\delta}+C^\delta\frac{\partial^3 D^\gamma}{\partial x_\alpha\partial x_\beta\partial x_\delta}\\
              &\quad+\frac{\partial C^\delta  }{\partial x_\beta}D^\eta\frac{\partial \Gamma^\gamma_{\delta\eta}}{\partial x_\alpha}+C^\delta\frac{\partial D^\eta}{\partial x_\beta}\frac{\partial \Gamma^\gamma_{\delta\eta}}{\partial x_\alpha}+\frac{\partial C^\delta }{\partial x_\alpha}D^\eta\frac{\partial \Gamma^\gamma_{\delta \eta}}{\partial x_\beta}+C^\delta \frac{\partial D^\eta}{\partial x_\alpha}\frac{\partial \Gamma^\gamma_{\delta \eta}}{\partial x_\beta}+C^\delta D^\eta\frac{\partial^2 \Gamma^\gamma_{\delta \eta}}{\partial x_\alpha\partial x_\beta}.
        \end{split}
    \end{equation*}
    Therefore
    \begin{equation*}
             \begin{split}          \nabla_{\A}&\nabla_{\B}\nabla_{\C}\D=\R(\A ,\B )\nabla_{\C} \D+A^\alpha B^\beta C^\eta\frac{\partial D^\gamma}{\partial x_\eta}\frac{\partial \Gamma_{\alpha\gamma}^\delta}{\partial x_\beta} \frac{\partial}{\partial x_\delta}\\
          &\quad+A^\alpha B^\beta\left(\frac{\partial ^2C^\delta}{\partial x_\alpha\partial x_\beta}\frac{\partial D^\gamma}{\partial x_\delta} +\frac{\partial C^\delta}{\partial x_\beta}\frac{\partial^2 D^\gamma}{\partial x_\alpha\partial x_\delta}+\frac{\partial C^\delta}{\partial x_\alpha}\frac{\partial^2 D^\gamma}{\partial x_\beta\partial x_\delta} +  C^\delta\frac{\partial^3 D^\gamma}{\partial x_\alpha\partial x_\beta\partial x_\delta} \right)\frac{\partial}{\partial x_\gamma}\\
          &\quad+\underbrace{A^\alpha B^\beta\left(\frac{\partial C^\delta  }{\partial x_\beta}D^\eta\frac{\partial \Gamma^\gamma_{\delta\eta}}{\partial x_\alpha} + C^\delta\frac{\partial D^\eta}{\partial x_\beta}\frac{\partial \Gamma^\gamma_{\delta\eta}}{\partial x_\alpha}\right)\frac{\partial}{\partial x_\gamma}}_{\mathrm{I}} +A^\alpha B^\beta\left(\frac{\partial C^\delta }{\partial x_\alpha}D^\eta\frac{\partial \Gamma^\gamma_{\delta \eta}}{\partial x_\beta} +C^\delta \frac{\partial D^\eta}{\partial x_\alpha}\frac{\partial \Gamma^\gamma_{\delta \eta}}{\partial x_\beta} \right)\frac{\partial}{\partial x_\gamma}\\
          &\quad+\underbrace{A^\alpha B^\beta C^\delta D^\eta\frac{\partial^2 \Gamma^\gamma_{\delta \eta}}{\partial x_\alpha\partial x_\beta} \frac{\partial}{\partial x_\gamma}}_{\mathrm{II}}+A^\alpha\frac{\partial B^\beta}{\partial x_\alpha}\left(\frac{\partial C^\delta}{\partial x_\beta}\frac{\partial D^\gamma}{\partial x_\delta}+ C^\delta\frac{\partial^2 D^\gamma}{\partial x_\beta\partial x_\delta}\right)\frac{\partial}{\partial x_\gamma}+ \underbrace{A^\alpha\frac{\partial B^\beta}{\partial x_\alpha}C^\delta D^\eta\frac{\partial \Gamma^\gamma_{\delta \eta}}{\partial x_\beta}\frac{\partial}{\partial x_\gamma}}_{\mathrm{III}}.
        \end{split}
    \end{equation*}
    Since 
    \begin{align*}
       \mathrm{I}\overset{\eqref{rieminnorm}}&{=}\R\left(\A ,\nabla_{\B} \C\right)\D+A^\alpha B^\beta\frac{\partial C^\delta}{\partial x_\beta}D^\eta\frac{\partial \Gamma^\gamma_{\alpha\eta}}{\partial x_\delta}\frac{\partial}{\partial x_\gamma}+\R\left(\A ,\C \right)\nabla_{\B} \D+A^\alpha B^\beta C^\delta\frac{\partial D^\eta}{\partial x_\beta}\frac{\partial\Gamma^\gamma_{\alpha\eta}}{\partial x_\delta}\frac{\partial}{\partial x_\gamma},\\
       \mathrm{II}&=A^\alpha B^\beta C^\delta D^\eta\frac{\partial^2 \Gamma^\gamma_{\delta \eta}}{\partial x_\beta\partial x_\alpha} \frac{\partial}{\partial x_\gamma}\overset{\eqref{deririeminnorm}}{=}\left(\nabla_{\B}\R\right)(\A,\C )\D+A^\alpha B^\beta C^\delta D^\eta\frac{\partial^2 \Gamma^\gamma_{\alpha \eta}}{\partial x_\beta\partial x_\delta} \frac{\partial}{\partial x_\gamma},\\
       \mathrm{III}\overset{\eqref{rieminnorm}}&{=}\R\left(\nabla_\A \B,\C\right)\D+A^\alpha\frac{\partial B^\beta}{\partial x_\alpha}C^\delta D^\eta\frac{\partial \Gamma^\gamma_{\beta \eta}}{\partial x_\delta}\frac{\partial}{\partial x_\gamma}
    \end{align*}
    then \eqref{nablaquattrocose} follows.
  %  \begin{equation*}
     %   \begin{split}
     %       \nabla_{\A}&\nabla_{\B}\nabla_{\C} D=\sum_{\alpha,\beta=1}^{n+1}A^\alpha\nabla_{\frac{\partial}{\partial x_\alpha}}\left(B^\beta \nabla_{\frac{\partial}{\partial x_\beta}}\nabla_{C}D\right)\\
     %       &=\sum_{\alpha,\beta,\gamma,\delta=1}^{n+1}A^\alpha B^\beta\nabla_{\frac{\partial}{\partial x_\alpha}}\left( \nabla_{\frac{\partial}{\partial x_\beta}}\left(C^\gamma\nabla_{\frac{\partial}{\partial x_\gamma}}\left(D^\delta\frac{\partial}{\partial x_\delta}\right)\right)\right)+\sum_{\beta,\gamma,\delta=1}^{n+1}\left(\nabla_{\A} B\right)^\beta \nabla_{\frac{\partial}{\partial x_\beta}}\left(C^\gamma\nabla_{\frac{\partial}{\partial x_\gamma}}\left(D^\delta\frac{\partial}{\partial x_\delta}\right)\right)\\
     %       &=\sum_{\alpha,\beta,\gamma,\delta=1}^{n+1}A^\alpha B^\beta\nabla_{\frac{\partial}{\partial x_\alpha}}\left( \nabla_{\frac{\partial}{\partial x_\beta}}\left(C^\gamma\nabla_{\frac{\partial}{\partial x_\gamma}}\left(D^\delta\frac{\partial}{\partial x_\delta}\right)\right)\right)+\nabla_{\nabla_{\A} B}\nabla_{C}D.
    %    \end{split}
   % \end{equation*}
 %   First, 
  %  \begin{equation*}
   %     \begin{split}
   %         \sum_{\beta,\gamma,\delta=1}^{n+1}\left(\nabla_{\A} B\right)^\beta \nabla_{\frac{\partial}{\partial x_\beta}}&\left(C^\gamma\nabla_{\frac{\partial}{\partial x_\gamma}}\left(D^\delta\frac{\partial}{\partial x_\delta}\right)\right)=\nabla_{\nabla_{\A} B}\nabla_{C}D\\
     %       &\overset{\eqref{nablatrecose}}{=}\R\left(\nabla_{\A} B,\C \right) D+\sum_{\alpha,\beta,\gamma,\delta=1}^{n+1}\left(\nabla_{\A} B\right)^\alpha C^\beta D^\gamma\frac{\partial \Gamma_{\alpha\gamma}^\delta}{\partial x_\beta} \frac{\partial}{\partial x_\delta}\\
     %       &+\sum_{\alpha,\beta,\gamma=1}^{n+1}\left(\nabla_{\A} B\right)^\alpha C^\beta\frac{\partial^2 D^\gamma}{\partial x_\alpha \partial x_\beta} \frac{\partial}{\partial x_\gamma}+\sum_{\beta,\gamma=1}^{n+1}\left(\nabla_{\nabla_{\A} B} C\right)^\beta\frac{\partial D^\gamma}{\partial x_\beta}\frac{\partial}{\partial x_\gamma}.
   %     \end{split}
   % \end{equation*}
 %   Fix $\alpha,\beta=1,\ldots,n+1$. Then
   % \begin{equation*}
   %     \begin{split}
    %        \sum_{\gamma,\delta=1}^{n+1}&\nabla_{\frac{\partial}{\partial x_\alpha}}\left( \nabla_{\frac{\partial}{\partial x_\beta}}\left(C^\gamma\nabla_{\frac{\partial}{\partial x_\gamma}}\left(D^\delta\frac{\partial}{\partial x_\delta}\right)\right)\right)\\
      %      & =\sum_{\gamma,\delta=1}^{n+1}\nabla_{\frac{\partial}{\partial x_\alpha}}\left( C^\gamma\nabla_{\frac{\partial}{\partial x_\beta}}\left(\nabla_{\frac{\partial}{\partial x_\gamma}}\left(D^\delta\frac{\partial}{\partial x_\delta}\right)\right)\right)+\sum_{\gamma,\delta=1}^{n+1}\nabla_{\frac{\partial}{\partial x_\alpha}}\left( \frac{\partial C^\gamma}{\partial x_\beta}\nabla_{\frac{\partial}{\partial x_\gamma}}\left(D^\delta\frac{\partial}{\partial x_\delta}\right)\right)\\
      %      &=\sum_{\gamma,\delta=1}^{n+1}C^\gamma\nabla_{\frac{\partial}{\partial x_\alpha}}\left( \nabla_{\frac{\partial}{\partial x_\beta}}\left(\nabla_{\frac{\partial}{\partial x_\gamma}}\left(D^\delta\frac{\partial}{\partial x_\delta}\right)\right)\right)\\
       %     &\quad+\sum_{\gamma,\delta=1}^{n+1}\frac{\partial C^\gamma}{\partial x_\alpha}\nabla_{\frac{\partial}{\partial x_\beta}}\left(\nabla_{\frac{\partial}{\partial x_\gamma}}\left(D^\delta\frac{\partial}{\partial x_\delta}\right)\right)+\sum_{\gamma,\delta=1}^{n+1}\nabla_{\frac{\partial}{\partial x_\alpha}}\left( \frac{\partial C^\gamma}{\partial x_\beta}\nabla_{\frac{\partial}{\partial x_\gamma}}\left(D^\delta\frac{\partial}{\partial x_\delta}\right)\right)\\
     %   \end{split}
  %  \end{equation*}
    \end{proof}    \subsection{Pointwise variations} In this section we deduce the pointwise evolution of the relevant geometric quantities. 
    The first lemma (cf. \cite[Lemma 1.13]{MR4676392} and \cite[Lemma 1.22]{MR4676392}) describes the behavior of $E$.
    \begin{lemma}\label{lemmaevolutiondie}
    Let $p\in S.$ Let $e\in T_p S$. Then 
    \begin{align}
        \left.\frac{D}{dt}\right|_{t=0}E(t)&=\nabla_{e}\X;\label{prounoaux}\\
        \left.\frac{D^2}{dt^2}\right|_{t=0}E(t)&=\nabla_{e}\Z+\R(\X,e)\X.\label{produeaux}
    \end{align}
\end{lemma}
By \Cref{lemmaevolutiondie}, we can compute the evolution of the Jacobian of $\Phi$. \Cref{lemmaderivatejacobiano} is surely well-known (cf. \cite[Theorem 1.11]{MR4676392} and \cite[Theorem 1.21]{MR4676392}). We include its proof for the sake of completeness.
\begin{lemma}\label{lemmaderivatejacobiano}
    Let $p\in S$. Then
    \begin{align}
        \left.\frac{d}{dt}\right|_{t=0}\jac\left(\Phi_t|_{S}\right)\left(p\right)&=\varphi H+\divv^S\X^T\label{derijacuno},\\
        \left.\frac{d^2}{dt^2}\right|_{t=0}\jac\left(\Phi_t|_{S}\right)\left(p\right)&=\left(\divv^S\X\right)^2-\left\langle\nabla^S\X,\left(\nabla^S\X\right)^t\right\rangle-\ric(\X,\X)-\R(\X^T,\n,\X^T,\n)\label{derijacdue}\\
        &\quad+\left|\nabla^S\varphi\right|^2-2h\left(\nabla^S\varphi,\X^T\right)+h^2\left(\X^T,\X^T\right)+\divv^S\Z\nonumber.
    \end{align}
    In particular, if $\Phi$ is a normal variation, 
    \begin{equation}\label{derjduenormal}
    \begin{split}
         \left.\frac{d^2}{dt^2}\right|_{t=0}\jac\left(\Phi_t|_{S}\right)\left(p\right)&=\varphi^2H^2-\varphi^2|h|^2-\varphi^2\ric(\n,\n)+\left|\nabla^S\varphi\right|^2+\divv^S\Z.
        \end{split}
    \end{equation}
\end{lemma}
\begin{proof}
    Since $\G(t)$ has full rank for every $t$ sufficiently small, Jacobi's formula grants that, for any such $t$,
    \begin{equation*}
        \frac{d \det \G(t)}{d t}=\det\G(t)\trace\left(\G(t)^{-1}\frac{d \G(t)}{dt}\right), 
    \end{equation*}
    whence
    \begin{equation*}
          \frac{d}{dt}\jac\left(\Phi_t|_{S}\right)\left(p\right)=\frac{1}{2}\jac\left(\Phi_t|_{S}\right)\left(p\right)\trace\left(\G(t)^{-1}\frac{d \G(t)}{dt}\right).
    \end{equation*}
   Recalling that $\G(0)_{ij}=\delta_{ij}$ for $i,j=1,\ldots,n$ by construction,
    \begin{align*}
         \left. \frac{d}{dt}\right|_{t=0}\jac\left(\Phi_t|_{S}\right)\left(p\right)=\frac{1}{2}\trace\left(\left.\frac{d }{dt}\right|_{t=0}\G(t)\right)\overset{\eqref{prounoaux}}{=}\sum_{i=1}^n\left\langle \nabla_{e_i}\X,e_i\right\rangle=\divv^S\X.
    \end{align*}
    Since 
    \begin{equation*}
        \divv^S\X=\divv^S\left(\varphi\n\right)+\divv^S\X^T=\varphi H+\divv^S\X^T,
    \end{equation*}
    \eqref{derijacuno} follows. Moreover, since \begin{equation}\label{derivatadimatriceinversa}
        \frac{d\G(t)^{-1}}{dt}=-\G(t)^{-1}\frac{d\G(t)}{dt}\G(t)^{-1},
    \end{equation}
    then
    \begin{equation*}
        \begin{split}
             \left. \frac{d^2}{dt^2}\right|_{t=0}\jac\left(\Phi_t|_{S}\right)\left(p\right)&\overset{\eqref{derijacuno}}{=}\left(\divv^S\X\right)^2+\underbrace{\frac{1}{2}\trace\left(-\left(\left.\frac{d }{dt}\right|_{t=0}\G(t)\right)^2\right)}_{\mathrm{I}}+\underbrace{\frac{1}{2}\trace\left(\left.\frac{d^2 }{dt^2}\right|_{t=0}\G(t)\right)}_{\mathrm{II}}.
        \end{split}
    \end{equation*}
    First,
    \begin{equation*}
        \begin{split}
            \mathrm{I}=
            \overset{\eqref{prounoaux}}{=}-\frac{1}{2}\sum_{i,j=1}^n\left(\left\langle\nabla_{e_i}\X,e_j\right\rangle+\left\langle \nabla_{e_j}\X,e_i\right\rangle\right)^2=-\sum_{i,j=1}^n\left\langle\nabla_{e_i}\X,e_j\right\rangle^2-\sum_{i,j=1}^n\left\langle\nabla_{e_i}\X,e_j\right\rangle\left\langle \nabla_{e_j}\X,e_i\right\rangle.
        \end{split}
    \end{equation*}
    Moreover,
    \begin{equation*}
        \begin{split}
            \mathrm{II}&\overset{\eqref{prounoaux},\eqref{produeaux}}{=}\sum_{i=1}^n\left\langle \nabla_{e_i}\X,\nabla_{e_i}\X\right\rangle+\sum_{i=1}^n\left\langle\nabla_{e_i}\Z+\R(\X,e_i)\X,e_i\right\rangle\\
            &=\sum_{i,j=1}^n\left\langle \nabla_{e_i}\X,e_j\right\rangle^2+\sum_{i=1}^n\left\langle \nabla_{e_i}\X,\n\right\rangle^2+\divv^S\Z-\ric(\X,\X)-\R(\X,\n,\X,\n)\\
            &=\sum_{i,j=1}^n\left\langle \nabla_{e_i}\X,e_j\right\rangle^2+\sum_{i=1}^n\left(e_i(\varphi)-\left\langle A\left(\X^T\right),e_i\right\rangle\right)^2+\divv^S\Z-\ric(\X,\X)-\R(\X^T,\n,\X^T,\n)\\
            &=\sum_{i,j=1}^n\left\langle \nabla_{e_i}\X,e_j\right\rangle^2+\left|\nabla^S\varphi\right|^2-2h\left(\nabla^S\varphi,\X^T\right)+h^2\left(\X^T,\X^T\right)+\divv^S\Z-\ric(\X,\X)-\R(\X^T,\n,\X^T,\n).\\
        \end{split}
    \end{equation*}
    In this way, \eqref{derijacdue} follows. Finally, if $\Phi$ is a normal variation, then $\X=\varphi\n$ and $\X^T=0$. In particular, $\nabla^S\X=\left(\nabla^S\X\right)^t=\varphi h$, and \eqref{derjduenormal} by \eqref{derijacdue}.
\end{proof}
Define $V(t)\coloneqq \N(t,\beta(t))$. By definition, $V$ is a vector field along $\beta$, and in particular $V(t)$ is normal to $S_t$ at $\beta(t)$ for every $t\in I$. $V$ evolves as follows.
\begin{lemma}
    Let $p\in S$. 
   % Set $V(t)=\N(t,\beta(t))$ for every $t\in\rr$. 
    Let $e_1,\ldots,e_n$ be an orthonormal basis of $T_p S$. Then
    \begin{align}
        &\,V'(p)\coloneqq \left.\frac{D}{dt}\right|_{t=0}V(t)=-\nabla^S\varphi+ A\left( \X^T\right) ; \label{lemma125ritore}\\
          &V''(p)\coloneqq \left.\frac{D^2}{dt^2}\right|_{t=0}V(t)=\left(-\left|\nabla^S\varphi\right|^2-\left|A\left(\X ^T\right)\right|^2+2 h\left( \X^T,\nabla^S\varphi\right)\right)\n\label{secondordernormalbehave}\\
         &\quad+\sum_{j=1}^n\left(-\left \langle \n ,\nabla_{e_j}\Z+\R(\X,e_j)\ \X^T\right\rangle+2\left\langle \nabla^S\varphi- A\left( \X^T\right) ,\varphi A(e_j)+\nabla_{e_j}\ \X^T\right\rangle\right)e_j.\nonumber
    \end{align}
    Moreover, fix local coordinates $x_1,\ldots,x_{n+1}$ in $M$ near $p$, and set 
    \begin{equation*}
        \n'(q)=\frac{\partial c^k}{\partial t}(0,q)\left.\frac{\partial}{\partial x_k}\right|_q,\qquad \n''(q)=\frac{\partial^2 c^k}{\partial t^2}(0,q)\left.\frac{\partial}{\partial x_k}\right|_q\qquad\text{locally around $p$.}
    \end{equation*}
 Then, when $q\in S$ is close to $p$,
 \begin{align}
        \n'(q)&=V'(q)-\nabla_\X \n, \label{piccolopezzodercovnorm}\\
        \n''(q)&    =V''(q)-2\nabla_\X\n'-\nabla_{\X'}\n-\nabla_\X\nabla_\X \n\label{pezzocondtquadrockappa}
 \end{align}
\end{lemma}
\begin{proof}
   First,  \eqref{lemma125ritore} follows by \cite[Lemma 1.25]{MR4676392}. We prove \eqref{secondordernormalbehave}. As $|V(t)|\equiv 1$,
%\begin{equation*}
  %  \left\langle\frac{D}{dt}V(t),V(t)\right\rangle=\frac{d}{dt}\left\langle V(t),V(t)\right\rangle- \left\langle\frac{D}{dt}V(t),V(t)\right\rangle=- \left\langle\frac{D}{dt}V(t),V(t)\right\rangle,
%\end{equation*}
then
\begin{equation}\label{solitacosa}
      \left\langle\frac{D}{dt}V(t),V(t)\right\rangle\equiv 0.
\end{equation}
In particular, by \eqref{lemma125ritore} and \eqref{solitacosa},
\begin{equation*}
    \begin{split}
         \left\langle\left.\frac{D^2}{dt^2}\right|_{t=0}V(t),\left .\n\right|_p\right\rangle&=\left.\frac{d}{dt}\right|_{t=0}\left\langle\frac{D}{dt}V(t),V(t)\right\rangle-\left|V'(p)\right|^2
         =-\left|\nabla^S\varphi\right|^2-\left|A\left(\X ^T\right)\right|^2+2 h\left( \X^T,\nabla^S\varphi\right).
    \end{split}
\end{equation*}
Let $f\in T_p S$. Set $F(t)\coloneqq \left(d\Phi_t\right)|_p(f)$. Recall that $F(t)$ is a vector field along $\beta$, and $F(t)\in T_{\beta(t)}S_t$. Therefore
\begin{equation*}
    \begin{split}
         \left\langle\left.\frac{D^2}{dt^2}\right|_{t=0}V(t),f\right\rangle&=\left.\frac{d}{dt}\right|_{t=0}\left\langle\frac{D}{dt}V(t),F(t)\right\rangle-\left\langle\left.\frac{D}{dt}\right|_{t=0}V(t),\left.\frac{D}{dt}\right|_{t=0}F(t)\right\rangle\\   
         &=-\left.\frac{d}{dt}\right|_{t=0}\left\langle V(t),\frac{D}{dt}F(t)\right\rangle-\left\langle\left.\frac{D}{dt}\right|_{t=0}V(t),\left.\frac{D}{dt}\right|_{t=0}F(t)\right\rangle\\  
         &=-\left\langle \n,\left.\frac{D^2}{dt^2}\right|_{t=0}F(t)\right\rangle-2\left\langle\left.\frac{D}{dt}\right|_{t=0}V(t),\left.\frac{D}{dt}\right|_{t=0}F(t)\right\rangle\\  
         \overset{\eqref{prounoaux},\eqref{produeaux},\eqref{lemma125ritore}}&{=}-\left \langle \n ,\nabla_{f}\left(\nabla_\X\X+\X'\right)+\R(\X,f)\X\right\rangle+2\left\langle \nabla^S\varphi- A\left( \X^T\right) ,\nabla_f \X\right\rangle\\
         &=-\left \langle \n,\nabla_{f}\Z+\R(\X,f)\ \X^T\right\rangle+2\left\langle \nabla^S\varphi- A\left(\ \X^T\right) ,\varphi A(f)+\nabla_f\ \X^T\right\rangle.
    \end{split}
\end{equation*}
In this way, \eqref{secondordernormalbehave} follows. Next we prove \eqref{piccolopezzodercovnorm}. Indeed, if $q\in S$ is close to $p$,
    \begin{equation*}
    \begin{split}
          \left.\frac{D}{dt}\right|_{t=0}\N(t,\beta_q(t))&\overset{\eqref{comefarecovdev}}{=}\frac{\partial c^k}{\partial t}(0,q)\left.\frac{\partial}{\partial x_k}\right|_q+X^i(q)\frac{\partial N^k}{\partial x_i}(q)\left.\frac{\partial}{\partial x_k}\right|_q+X^i(q)N^k(q)\Gamma_{ik}^l(q)\left.\frac{\partial}{\partial x_l}\right|_q=\n'+\nabla_\X \n.
    \end{split}
    \end{equation*}
 Finally, for $q$ close to $p$,
\begin{equation*}
\begin{split}
    & \left.\frac{D^2}{dt^2}\right|_{t=0}\mathcal N(t,\beta_q(t))=\left.\frac{D^2}{dt^2}\right|_{t=0}\left(c^k(t,\beta_q(t))\left.\frac{\partial}{\partial x_k}\right|_{\beta_q(t)}\right)\\
    &\quad=\left.\frac{D}{dt}\right|_{t=0}\left(\frac{\partial c^k}{\partial t}(t,\beta_q(t))\left.\frac{\partial}{\partial x_k}\right|_{\beta_q(t)}+a^i(t,\beta_q(t))\frac{\partial c^k}{\partial x_i}(t,\beta_q(t))\left.\frac{\partial}{\partial x_k}\right|_{\beta_q(t)}\right)\\
    &\qquad+\left.\frac{D}{dt}\right|_{t=0}\left(a^i(t,\beta_q(t))c^k(t,\beta_q(t))\Gamma_{ik}^l(\beta_q(t))\left.\frac{\partial}{\partial x_l}\right|_{\beta_q(t)}\right)\\
    &\quad=\n''(q)+2X^i(q)\frac{\partial^2 c^k}{\partial x_i\partial t}(0,q)\left.\frac{\partial}{\partial x_k}\right|_{q}+2X^i(q)\frac{\partial c^k}{\partial t}(0,q)\Gamma_{ik}^l(q)\left.\frac{\partial}{\partial x_l}\right|_{q}+(X')^i(q)\frac{\partial N^k}{\partial x_i}(q)\left.\frac{\partial}{\partial x_k}\right|_{q}\\
    &\qquad+X^l(q)\frac{\partial X^i}{\partial x_l}(q)\frac{\partial N^k}{\partial x_i}(q)\left.\frac{\partial}{\partial x_k}\right|_{q}+X^i(q)X^l(q)\frac{\partial ^2N^k}{\partial x_l\partial x_i}(q)\left.\frac{\partial}{\partial x_k}\right|_{q}+2X^i(q)X^l(q)\frac{\partial N^k}{\partial x_i}(q)\Gamma_{lk}^m(q)\left.\frac{\partial}{\partial x_m}\right|_{q}\\
    &\qquad+\left(X'\right)^i(q)N^k(q)\Gamma_{ik}^l(q)\left.\frac{\partial}{\partial x_l}\right|_{q}+X^m(q)\frac{\partial X^i}{\partial x_m}(q)N^k(q)\Gamma_{ik}^l(q)\left.\frac{\partial}{\partial x_l}\right|_{q}\\
    &\qquad+X^i(q)X^m(q)N^k(q)\frac{\partial \Gamma_{ik}^l}{\partial x_m}(q)\left.\frac{\partial}{\partial x_l}\right|_{q}+X^i(q)X^m(q)N^k(q)\Gamma^l_{ik}(q)\Gamma^s_{ml}\left.\frac{\partial}{\partial x_s}\right|_q,
\end{split}
\end{equation*}
whence \eqref{pezzocondtquadrockappa} follows by \eqref{nablatrecosenonnormale}.
\end{proof}
Next, we describe the evolution of the shape operator. The following proof relies crucially on \Cref{fondamentalecomplemmavariazioniriemannianejhrnfi4rfjnfin}.
\begin{lemma}\label{lemmaevoluzioneshapeoperatoririem3858943}
    Let $p\in S.$ Let $e\in T_p S$. Then
    \begin{align}
          \left.\frac{D}{dt}\right|_{t=0}\left(\left.\nabla_{E(t)}\N(t,\cdot)\right|_{\beta(t)}\right)&=\nabla_eV'+\R(\X,e)\n ,\label{protreaux}\\
            \left.\frac{D^2}{dt^2}\right|_{t=0}\left(\left.\nabla_{E(t)}\N(t,\cdot)\right|_{\beta(t)}\right)&=\nabla_e V''+\R\left(\X,\nabla_e\X\right)\n+2\R(\X,e)V'+\R\left(\Z,e\right)\n+\left(\nabla_\X \R\right)(\X,e)\n\label{proquattroaux},
    \end{align}
    where $V'$ and $V''$ are given respectively by \eqref{lemma125ritore} and \eqref{secondordernormalbehave}.
\end{lemma}
\begin{proof}
  %  \textcolor{red}{First way, che non riesco a finire perché non capisco alcune cose}
 %   Fix $s\in\rr$. define $\Psi^s:\rr\times M\to M$ by $\Psi^s_t(p)=\left(\Phi_{t+s}\circ\Phi^{-1}_s\right)(p)$. Notice that $\Psi^s$ is a smooth variation of $S$, and its velocity $X^s$ is given by \eqref{velocityofpsis}. Set $S^s_t\coloneqq \Psi^s_t(\Phi_s(S))=\Phi_{t+s}(S)$. Notice that $N^s_t\coloneqq N_{t+s}$ is a smooth choice of unit normals to $S^s_t$. Set $p^s=\Phi_s(p)$,  $e^s=d\Phi_s|_p(e)$ and $E^s_t=d\Psi^s_t|_{p^s}(e^s)$. Notice that $\nabla_{E^s_t}N^s_t$ is a vector field along the curve $t\mapsto\Psi^s_t(\Phi_s(p))=\beta(t+s)$. Set $Y_s=X^s|_{\beta(s)}$. By the third property,
 %   \begin{equation}\label{usaterzaformula}
 %        \left.\frac{D}{dt}\nabla_{E^s_t}N^s_t\right|_{t=0}=\R(Y^s,e^s)N^s+\nabla_{e^s}\nabla_{X^s}N^s.
 %   \end{equation}
 %   Since $E_t^s=E_{t+s}$ and $N_t^s=N_{t+s}$, we conclude that 
 %   \begin{equation*}
 %       \left.\frac{D}{dt}\nabla_{E_t}N_t\right|_{t=s}= \left.\frac{D}{dt}\nabla_{E_{t+s}}N_{t+s}\right|_{t=0}\overset{\eqref{usaterzaformula}}{=}\R(Y_s,E_s)N_s+\nabla_{E_s}\nabla_{X^s}N_s.
 %   \end{equation*}
 %   In particular,
 %   \begin{equation*}
   %      \left.\frac{D^2}{dt^2}\right|_{t=0}\nabla_{E_t}N_t= \left.\frac{D}{dt}\right|_{t=0}\left(\R(Y_t,E_t)N_t+\nabla_{E_t}\nabla_{Y_t}N_t\right).
%    \end{equation*}     
 %   Consider local normal coordinates $x_1,\ldots, x_{n+1}$ in a neighborhood of $p$, say $U$, in such a way that $\Gamma_{ij}^{k}(p)=0$, where $\Gamma^k_{ij}$ are the Christoffel symbols with respect to $x_1,\ldots,x_{n+1}$. By continuity of $\Phi$, there exists $\varepsilon>0$ and a neighborhood $V\subseteq U$ of $p$ such that $\Phi_t(q)\in U$ for any $t\in(-\eps,\eps)$ and any $q\in V$. 
 %   Set $\Phi^i(t,p)=\left(x_i\circ\Phi_t\right)(p)$. 
  %  Locally, we may set 
  %  \begin{equation*}
  %      Y_t=\sum_{i=1}^{n+1}a_i(t)\left.\frac{\partial}{\partial x_i}\right|_{\beta(t)},\quad E_t=\sum_{j=1}^{n+1}b_j(t)\left.\frac{\partial}{\partial x_j}\right|_{\beta(t)},\quad N_t=\sum_{k=1}^{n+1}c_k(t)\left.\frac{\partial}{\partial x_k}\right|_{\beta(t)},\quad F_t=\sum_{l=1}^{n+1}d_l(t)\left.\frac{\partial}{\partial x_l}\right|_{\beta(t)}.
   % \end{equation*}
  %  \begin{equation*}
  %      \left. X^t\right|_q=\sum_{i=1}^{n+1}f_i(t,q)\left.\frac{\partial}{\partial x_i}\right|_{q},\qquad \left. X\right|_q=\sum_{i=1}^{n+1}f_i(q)\left.\frac{\partial}{\partial x_i}\right|_{q}
%    \end{equation*}
 %   for some smooth functions $f_1,\ldots,f_{n+1}$, where with some abuse of notation we have set $f_j(p)=f_j(0,p)$. 
 % Owing respectively to the proof of \cite[Proposition 1.23]{MR4676392} and to \cite[Lemma 1.25]{MR4676392},
%
 % \begin{equation}\label{devcovdiyen}
  %    \begin{split}
   %         \left.\frac{D}{dt}\right|_{t=0}Y_t&=X'+\nabla_X X,\\
   %   \left.\frac{D}{dt}\right|_{t=0}N^t&= A\left( \X^T\right) -\nabla^S\varphi.\\
   %   \end{split}
 % \end{equation}
   %Moreover,
  %  \begin{equation*}
  %       E_t=\sum_{i=1}^{n+1}g_i(t)\left.\frac{\partial}{\partial x_i}\right|_{\beta(t)},\qquad N_t=\sum_{i=1}^{n+1}h_i(t)\left.\frac{\partial}{\partial x_i}\right|_{\beta(t)}.
 %   \end{equation*}
 %Define
% \begin{equation*}
 %    R_{ijkl}(t)=\R\left(\left.\frac{\partial}{\partial x_i}\right|_{\beta(t)},\left.\frac{\partial}{\partial x_j}\right|_{\beta(t)},\left.\frac{\partial}{\partial x_k}\right|_{\beta(t)},\left.\frac{\partial}{\partial x_l}\right|_{\beta(t)}\right)\qquad\text{for any $i,j,k,l=1,\ldots,n+1$.}
% \end{equation*}
%If $Z_t=\displaystyle\sum_{i=1}^{n+1}z_i(t)\left.\frac{\partial}{\partial x_i}\right|_{\beta(t)}$ is any vector field along $\beta$, since $x_1,\ldots,x_{n+1}$ are normal coordinates at $p$, then
%\begin{equation}\label{covdivdernorm}
%    \left.\frac{D}{dt}\right|_{t=0}Z_t=\sum_{i=1}^n\dot z_i(0)\frac{\partial}{\partial x_i}+\sum_{i=1}^n z_i(0)\nabla_{X}\frac{\partial}{\partial x_i}=\sum_{i=1}^n\dot z_i(0)\frac{\partial}{\partial x_i}.
%\end{equation}
%Therefore
%\begin{equation*}
 %   \begin{split}
 %       &\left\langle\left.\frac{D}{dt}\right|_{t=0}\R(Y_t,E_t)N_t,F_t\right\rangle = \left.\frac{d}{dt}\right|_{t=0} \R\left(Y_t,E_t,N_t,F_t\right)- \R\left(X,e,N,\nabla_f X\right)\\
 %       &\quad=\sum_{i,j,k,l=1}^{n+1}\left.\frac{d}{dt}\right|_{t=0}\left(a_i(t)b_j(t) c_k(t)d_l(t) R_{ijkl}(t)\right)- \R\left(X,e,N,\nabla_f X\right)\\
 %       &\quad=\sum_{i,j,k,l=1}^{n+1}\left(\dot a_i(0)b_j(0) c_k(0)d_l(0)+ a_i(0)\dot b_j(0) c_k(0)d_l(0) +a_i(0)b_j(0)\dot  c_k(0)d_l(0)\right)R_{ijkl}(0)\\
  %      &\qquad+\sum_{i,j,k,l=1}^{n+1} a_i(0)b_j(0) c_k(0)\dot d_l(0)R_{ijkl}(0)+\sum_{i,j,k,l=1}^{n+1} a_i(0)b_j(0) c_k(0) d_l(0)\dot R_{ijkl}(0)- \R\left(X,e,N,\nabla_f X\right)\\
  %      &\quad\overset{\eqref{covdivdernorm},\eqref{devcovdiyen}}{=}\R\left(X'+\nabla_X X,e,N,f\right)+\R\left(X,\nabla_e X,N,f\right)+\R\left(X,e, A\left( \X^T\right) -\nabla^S\varphi,f\right)\\
  %      & \qquad+\left(\nabla_X \R\right)\left(X,e,N,f\right).
 %   \end{split}
%\end{equation*}
%Moreover,
%\begin{equation*}
 %   \begin{split}
  %      \left\langle\left.\frac{D}{dt}\right|_{t=0}\nabla_{E_t}\nabla_{Y_t}N_t,F_t\right\rangle&=\left.\frac{d}{dt}\right|_{t=0} \left\langle\nabla_{E_t}\nabla_{Y_t}N_t,F_t\right\rangle-\left\langle \nabla_{e}\nabla_{X}N,\nabla_f X\right\rangle\\
  %      &=
  %  \end{split}
%%\end{equation*}
%\textcolor{purple}{Second way}

Consider local normal coordinates $x_1,\ldots, x_{n+1}$ centered at $p$.  Since $E(t)\in T_{\beta(t)}S_t,$ then
    \begin{equation*}
        \begin{split}
            \left.\nabla_{E(t)}\N(t,\cdot)\right|_{\beta(t)}&=b^j(t)\frac{\partial c^k}{\partial x_j}(t,\beta(t))\left.\frac{\partial}{\partial x_k}\right|_{\beta(t)}+b^j(t)c^k(t,\beta(t))\Gamma_{jk}^l(\beta(t))\left.\frac{\partial}{\partial x_l}\right|_{\beta(t)}.
        \end{split}
    \end{equation*}
    Therefore
    \begin{equation}\label{derivataprimashape}
        \begin{split}
           &\frac{D}{dt}\Big(\left.\nabla_{E(t)}\N(t,\cdot)\right|_{\beta(t)}\Big)=\dot b^j(t)\frac{\partial c^k}{\partial x_j}(t,\beta(t))\left.\frac{\partial}{\partial x_k}\right|_{\beta(t)}+b^j(t)\frac{\partial^2 c^k}{\partial t\partial x_j}(t,\beta(t))\left.\frac{\partial}{\partial x_k}\right|_{\beta(t)}\\
           &\quad+a^i(t,\beta(t))b^j(t)\frac{\partial ^2c^k}{\partial x_i\partial x_j}(t,\beta(t))\left.\frac{\partial}{\partial x_k}\right|_{\beta(t)}+a^i(t,\beta(t)) b^j(t)\frac{\partial c^k}{\partial x_j}(t,\beta(t))\Gamma_{ik}^l(\beta(t))\left.\frac{\partial}{\partial x_l}\right|_{\beta(t)}\\
           &\quad+\dot b^j(t)c^k(t,\beta(t))\Gamma_{jk}^l(\beta(t))\left.\frac{\partial}{\partial x_l}\right|_{\beta(t)}+b^j(t)\frac{\partial c^k}{\partial t}(t,\beta(t))\Gamma_{jk}^l(\beta(t))\left.\frac{\partial}{\partial x_l}\right|_{\beta(t)}\\
           &\quad+a^i(t,\beta(t))b^j(t)\frac{\partial c^k}{\partial x_i}(t,\beta(t))\Gamma_{jk}^l(\beta(t))\left.\frac{\partial}{\partial x_l}\right|_{\beta(t)}+a^i(t,\beta(t))b^j(t)c^k(t,\beta(t))\frac{\partial\Gamma_{jk}^l}{\partial x_i}(\beta(t))\left.\frac{\partial}{\partial x_l}\right|_{\beta(t)}\\
           &\quad+a^i(t,\beta(t))b^j(t)c^k(t,\beta(t))\Gamma_{jk}^l(\beta(t))\Gamma_{i l}^m(\beta(t))\left.\frac{\partial}{\partial x_m}\right|_{\beta(t)}.
        \end{split}
    \end{equation}
    In particular, by \eqref{normcoordcristo},
      \begin{equation*}
        \begin{split}
           \left.\frac{D}{dt}\right|_{t=0}\Big(\left.\nabla_{E(t)}\N(t,\cdot)\right|_{\beta(t)}\Big)&=\dot b^j(0)\frac{\partial N^k}{\partial x_j}(p)\left.\frac{\partial}{\partial x_k}\right|_{p}+e^j\frac{\partial^2 c^k}{\partial t\partial x_j}(0,p)\left.\frac{\partial}{\partial x_k}\right|_{p}        \\
&\quad+X^i(p)e^j\frac{\partial ^2N^k}{\partial x_i\partial x_j}(p)\left.\frac{\partial}{\partial x_k}\right|_{p}+X^i(p)e^jN^k(p)\frac{\partial\Gamma_{jk}^l}{\partial x_i}(p)\left.\frac{\partial}{\partial x_l}\right|_{p}.
    %       &\overset{\eqref{rieminnorm}}{=}\sum_{i,k=1}^{n+1}\dot b^i(0)\frac{\partial c^k}{\partial x_i}(0,p)\left.\frac{\partial}{\partial x_k}\right|_{p}\\
       %    &\quad+\sum_{j,k=1}^{n+1}b^j(0)\frac{\partial^2 c^k}{\partial x_j\partial t}(0,p)\left.\frac{\partial}{\partial x_k}\right|_{p}+\sum_{i,j,k=1}^{n+1}a^i(0,p)b^j(0)\frac{\partial ^2c^k}{\partial x_j\partial x_i}(0,p)\left.\frac{\partial}{\partial x_k}\right|_{p}\\         %  
        %   &\quad+\R(X,e) N|_{p}+\sum_{i,j,k,l=1}^{n+1}a^i(0,p)b^j(0)c^k(0,p)\frac{\partial\Gamma_{ik}^l}{\partial x_j}(p)\left.\frac{\partial}{\partial x_l}\right|_{p}.
        \end{split}
    \end{equation*}
    Observe that 
    \begin{equation}\label{comeebdotzero}
       \dot b^i(0)\left.\frac{\partial}{\partial x_i}\right|_{p}\overset{\eqref{normcoordcristo}}{=} \left.\frac{D}{dt}\right|_{t=0}E(t)\overset{\eqref{prounoaux}}{=}\nabla_e \X\overset{\eqref{normcoordcristo}}{=}\frac{\partial X^i}{\partial x_j}(p)e^j\left.\frac{\partial}{\partial x_i}\right|_p,
    \end{equation}
   whence   \eqref{protreaux} follows by
    \begin{equation*}
        \begin{split}
           \left.\frac{D}{dt}\right|_{t=0}\Big(\left.\nabla_{E(t)}\N(t,\cdot)\right|_{\beta(t)}\Big)&=\frac{\partial X^i}{\partial x_j}(p)e^j\frac{\partial N^k}{\partial x_i}(p)\left.\frac{\partial}{\partial x_k}\right|_{p}+e^j\frac{\partial^2 c^k}{\partial x_j\partial t}(0,p)\left.\frac{\partial}{\partial x_k}\right|_{p}        \\
&\quad+X^i(p)e^j\frac{\partial ^2N^k}{\partial x_i\partial x_j}(p)\left.\frac{\partial}{\partial x_k}\right|_{p}+X^i(p)e^jN^k(p)\frac{\partial\Gamma_{jk}^l}{\partial x_i}(p)\left.\frac{\partial}{\partial x_l}\right|_{p}\\
\overset{\eqref{nablatrecose}}&{=}e^j\frac{\partial^2 c^k}{\partial x_j\partial t}(0,p)\left.\frac{\partial}{\partial x_k}\right|_{p} +\nabla_e\nabla_\X \n-\R(e,\X)\n\\
\overset{\eqref{normcoordcristo}}&{=}\nabla_e\n'+\nabla_e\nabla_\X \n-\R(e,\X)\n\\
\overset{\eqref{piccolopezzodercovnorm}}&{=}\nabla_eV'-\R(e,\X)\n.
        \end{split}
    \end{equation*}  
%   Since
 %   \begin{equation*}
  %      \begin{split}
   %         \nabla_eN'=&\overset{\eqref{piccolopezzodercovnorm}}{=}\-\nabla_e\left(\nabla^S\varphi\right)-\nabla_e\left(\varphi\nabla_NN\right),
   %     \end{split}
   % \end{equation*}
 We prove \eqref{proquattroaux}. By \eqref{derivataprimashape} and \eqref{normcoordcristo},
    \begin{equation*}
        \begin{split}
             &\left.\frac{D^2}{dt^2}\right|_{t=0}\Big(\left.\nabla_{E(t)}\N(t,\cdot)\right|_{\beta(t)}\Big)=\ddot b^j(0)\frac{\partial N^k}{\partial x_j}(p)\left.\frac{\partial}{\partial x_k}\right|_{p}+2\dot b^j(0)\frac{\partial^2 c^k}{\partial x_j\partial t}(0,p)\left.\frac{\partial}{\partial x_k}\right|_{p}\\
             &\quad+2X^i(p)\dot b^j(0)\frac{\partial^2 N^k}{\partial x_i\partial x_j}(p)\left.\frac{\partial}{\partial x_k}\right|_{p}+e^j\frac{\partial^3 c^k}{\partial x_j\partial t^2}(0,p)\left.\frac{\partial}{\partial x_k}\right|_{p}+2X^i(p)e^j\frac{\partial^3 c^k}{\partial x_i\partial x_j\partial t}(0,p)\left.\frac{\partial}{\partial x_k}\right|_{p}\\
             &\quad +(X')^i(p)e^j\frac{\partial ^2N^k}{\partial x_i\partial x_j}(p)\left.\frac{\partial}{\partial x_k}\right|_{p}+X^l(p)\frac{\partial X^i}{\partial x_l}(p)e^j\frac{\partial ^2N^k}{\partial x_i\partial x_j}(p)\left.\frac{\partial}{\partial x_k}\right|_{p}+X^i(p)X^l(p)e^j\frac{\partial ^3N^k}{\partial x_l\partial x_i\partial x_j}(p)\left.\frac{\partial}{\partial x_k}\right|_{p}\\
             &\quad+X^i(p)X^m(p) e^j\frac{\partial N^k}{\partial x_j}(p)\frac{\partial \Gamma_{ik}^l}{\partial x_m}(p)\left.\frac{\partial}{\partial x_l}\right|_{p}+2X^i(p)\dot b^j(0)N^k(p)\frac{\partial \Gamma_{jk}^l}{\partial x_i}(p)\left.\frac{\partial}{\partial x_l}\right|_{p}\\
             &\quad +2X^i(p) e^j\frac{\partial c^k}{\partial t}(0,p)\frac{\partial \Gamma_{jk}^l}{\partial x_i}(p)\left.\frac{\partial}{\partial x_l}\right|_{p}+2X^i(p)X^m(p)e^j\frac{\partial N^k}{\partial x_i}(p)\frac{\partial \Gamma_{jk}^l}{\partial x_m}(p)\left.\frac{\partial}{\partial x_l}\right|_{p}\\
            &\quad+(X')^i(p)e^jN^k(p)\frac{\partial\Gamma_{jk}^l}{\partial x_i}(p)\left.\frac{\partial}{\partial x_l}\right|_{p}+X^m(p)\frac{\partial X^i}{\partial x_m}(p)e^jN^k(p)\frac{\partial\Gamma_{jk}^l}{\partial x_i}(p)\left.\frac{\partial}{\partial x_l}\right|_{p}\\
            &\quad+X^i(p)X^m(p)e^jN^k(p)\frac{\partial^2\Gamma_{jk}^l}{\partial x_m\partial x_i}(p)\left.\frac{\partial}{\partial x_l}\right|_{p}.
        \end{split}
    \end{equation*}
%    Observe that
 %   \begin{equation}\label{formacorddixprimoenablaxx}
    %    \begin{split}
      %      \left(\nabla_X X\right)(q)&=\sum_{i,m=1}^{n+1}a^m(0,q)\frac{\partial a^i}{\partial x_m}(0,q)\left.\frac{\partial}{\partial x_i}\right|_{q}+\sum_{i,l,m=1}^{n+1}a^i(0,q)a^m(0,q)\Gamma_{m,i}^l\left.\frac{\partial}{\partial x_l}\right|_{q}\qquad\text{for any $q\in V$.}
    %    \end{split}
  %  \end{equation}
       Notice that 
       \begin{equation*}
       \begin{split}
            \left.\frac{D^2}{dt^2}\right|_{t=0}E(t)\overset{\eqref{comefarecovdev}}&{=} \left.\frac{D}{dt}\right|_{t=0}\left(\dot b^j(t)\left.\frac{\partial}{\partial x_j}\right|_{\beta(t)}+a^i(t,\beta(t))b^j(t)\Gamma_{ij}^l(\beta(t))\left.\frac{\partial}{\partial x_l}\right|_{\beta(t)}\right)\\
            \overset{\eqref{normcoordcristo}}&{=}\ddot b^j(0)\left.\frac{\partial}{\partial x_j}\right|_{p}+X^i(p)X^m(p)e^j\frac{\partial \Gamma_{ij}^l}{\partial x_m}(p)\left.\frac{\partial}{\partial x_l}\right|_{p}\\
            &=\left(\ddot b^j(0)+X^i(p)X^m(p)e^l\frac{\partial \Gamma_{il}^j}{\partial x_m}(p)\right) \left.\frac{\partial}{\partial x_j}\right|_{p}.
       \end{split}
       \end{equation*}
       Therefore, by \eqref{normcoordcristo}, \eqref{nablatrecose}, \eqref{produeaux} and for any $j=1,\ldots,n+1$,      \begin{equation}\label{dersecetcond}
       \begin{split}
              &\ddot b^j(0)=\left(\nabla_{e}\nabla_\X\X\right)^j(p)+\left(\nabla_e \X'\right)^j(p)+\left(\R(\X,e)\X\right)^j(p)-X^i(p)X^m(p)e^l\frac{\partial \Gamma_{il}^j}{\partial x_m}(p)\\
              &\quad=e^lX^i(p)X^m(p)\frac{\partial\Gamma_{lm}^j}{\partial x_i}+ e^l X^i(p)\frac{\partial^2 X^j}{\partial x_l \partial x_i}(p)+e^l\frac{\partial X^i}{\partial x_l}(p)\frac{\partial X^j}{\partial x_i}(p)+e^l\frac{\partial (X')^j}{\partial x_l}(p)-X^i(p)X^m(p)e^l\frac{\partial\Gamma_{il}^j}{\partial x_m}(p)\\
              &\quad= e^l X^i(p)\frac{\partial^2 X^j}{\partial x_l \partial x_i}(p)+e^l\frac{\partial X^i}{\partial x_l}(p)\frac{\partial X^j}{\partial x_i}(p)+e^l\frac{\partial (X')^j}{\partial x_l}(p),
       \end{split}
       \end{equation}
       where the last equality follows by the symmetry of $\nabla$. 
 Then, by  \eqref{comeebdotzero} and \eqref{dersecetcond},
    \begin{equation*}
        \begin{split}
             &\left.\frac{D^2}{dt^2}\right|_{t=0}\Big(\left.\nabla_{E(t)}\N(t,\cdot)\right|_{\beta(t)}\Big)=e^l X^i(p)\frac{\partial^2 X^j}{\partial x_l \partial x_i}(p)\frac{\partial N^k}{\partial x_j}(p)\left.\frac{\partial}{\partial x_k}\right|_{p}+e^l\frac{\partial X^i}{\partial x_l}(p)\frac{\partial X^j}{\partial x_i}(p)\frac{\partial N^k}{\partial x_j}(p)\left.\frac{\partial}{\partial x_k}\right|_{p}\\
             &\quad+e^l\frac{\partial (X')^j}{\partial x_l}(p)\frac{\partial N^k}{\partial x_j}(p)\left.\frac{\partial}{\partial x_k}\right|_{p}+2\frac{\partial X^i}{\partial x_j}(p)e^j\frac{\partial^2 c^k}{\partial x_i\partial t}(0,p)\left.\frac{\partial}{\partial x_k}\right|_{p}\\
             &\quad+2X^i(p)\frac{\partial X^m}{\partial x_j}(p)e^j\frac{\partial^2 N^k}{\partial x_i\partial x_m}(p)\left.\frac{\partial}{\partial x_k}\right|_{p}+e^j\frac{\partial^3 c^k}{\partial x_j\partial t^2}(0,p)\left.\frac{\partial}{\partial x_k}\right|_{p}+2X^i(p)e^j\frac{\partial^3 c^k}{\partial x_i\partial x_j\partial t}(0,p)\left.\frac{\partial}{\partial x_k}\right|_{p}\\
             &\quad +(X')^i(p)e^j\frac{\partial ^2N^k}{\partial x_i\partial x_j}(p)\left.\frac{\partial}{\partial x_k}\right|_{p}+X^l(p)\frac{\partial X^i}{\partial x_l}(p)e^j\frac{\partial ^2N^k}{\partial x_i\partial x_j}(p)\left.\frac{\partial}{\partial x_k}\right|_{p}+X^i(p)X^l(p)e^j\frac{\partial ^3N^k}{\partial x_l\partial x_i\partial x_j}(p)\left.\frac{\partial}{\partial x_k}\right|_{p}\\
             &\quad+X^i(p)X^m(p) e^j\frac{\partial N^k}{\partial x_j}(p)\frac{\partial \Gamma_{ik}^l}{\partial x_m}(p)\left.\frac{\partial}{\partial x_l}\right|_{p}+2X^i(p)\frac{\partial X^m}{\partial x_j}(p)e^jN^k(p)\frac{\partial \Gamma_{mk}^l}{\partial x_i}(p)\left.\frac{\partial}{\partial x_l}\right|_{p}\\
             &\quad +2X^i(p) e^j\frac{\partial c^k}{\partial t}(0,p)\frac{\partial \Gamma_{jk}^l}{\partial x_i}(p)\left.\frac{\partial}{\partial x_l}\right|_{p}+2X^i(p)X^m(p)e^j\frac{\partial N^k}{\partial x_i}(p)\frac{\partial \Gamma_{jk}^l}{\partial x_m}(p)\left.\frac{\partial}{\partial x_l}\right|_{p}\\
            &\quad+(X')^i(p)e^jN^k(p)\frac{\partial\Gamma_{jk}^l}{\partial x_i}(p)\left.\frac{\partial}{\partial x_l}\right|_{p}+X^m(p)\frac{\partial X^i}{\partial x_m}(p)e^jN^k(p)\frac{\partial\Gamma_{jk}^l}{\partial x_i}(p)\left.\frac{\partial}{\partial x_l}\right|_{p}\\
            &\quad+X^i(p)X^m(p)e^jN^k(p)\frac{\partial^2\Gamma_{jk}^l}{\partial x_m\partial x_i}(p)\left.\frac{\partial}{\partial x_l}\right|_{p}.
        \end{split}
    \end{equation*}
    A careful comparison between the above expression and \eqref{nablatrecose} and \eqref{nablaquattrocose} grants that
    \begin{equation*}
        \begin{split}
            \left.\frac{D^2}{dt^2}\right|_{t=0}&\Big(\left.\nabla_{E(t)}\N(t,\cdot)\right|_{\beta(t)}\Big)=\nabla _e\nabla_\X\nabla_\X \n-\left(\nabla_\X \R\right)(e,\X)\n-2 \R(e,\X)\nabla_\X \n-\R\left(e,\nabla_\X \X\right)\n\\
            &\quad-\R\left(\nabla_e\X,\X\right)\n+2\nabla_e\nabla_\X\n'-2\R(e,\X)\n'+\nabla_e \n''+\nabla_e\nabla_{\X'}\n-\R(e,\X')\n.
        \end{split}
    \end{equation*}
    In addition,
    \begin{equation*}
        -2 \R(e,\X)\n'-2\R(e,\X)\nabla_\X \n\overset{\eqref{piccolopezzodercovnorm}}{=}-2\R(e,\X)V'
    \end{equation*}
    and 
    \begin{equation*}
        \nabla_e \n''+2\nabla_e\nabla_\X \n'+\nabla_e \nabla_{\X'}\n+\nabla_e\nabla_\X\nabla_\X \n\overset{\eqref{pezzocondtquadrockappa}}{=}\nabla_e V'',
    \end{equation*}
    whence
  \begin{equation*}
        \begin{split}
            \left.\frac{D^2}{dt^2}\right|_{t=0}\Big(\left.\nabla_{E(t)}\N(t,\cdot)\right|_{\beta(t)}\Big)&=-\left(\nabla_\X \R\right)(e,\X)\n-\R\left(e,\Z\right)\n+\nabla_e V''-2\R(e,\X)V'-\R\left(\nabla_e\X,\X\right)\n,
        \end{split}
    \end{equation*}
    and \eqref{proquattroaux} follows.
\end{proof}
\Cref{lemmaevoluzioneshapeoperatoririem3858943} allows to compute the pointwise evolution of the mean curvature at first and second-order.
\begin{lemma}\label{lemmavariazioneprimacurvaturaeseconda}
    Let $p\in S$. Then 
    \begin{align}
        \left.\frac{d}{dt}\right|_{t=0}H_t\left(\Phi_t(p)\right)&=-\Delta^S\varphi-\varphi\left(|h|^2+\ric(\n,\n)\right)+ \X^T H\label{derprimhinstatement}\\   
         \left.\frac{d^2}{dt^2}\right|_{t=0}H_t\left(\Phi_t(p)\right)&=2\left\langle h,\left(\nabla^S X\right)^2\right\rangle-2\left\langle \nabla^S X,\left(\nabla^SV'\right)^t\right\rangle-\left\langle h,\j(\X,\X)\right\rangle-\left\langle\nabla^S \X,\j(\n,\X)\right\rangle\label{dersechinstatement}\\
         &\quad+\left(-\left|\nabla^S\varphi\right|^2+2h\left(\nabla^S\varphi, \X^T\right)-h^2\left( \X^T, \X^T\right)\right)H+\divv^S\left( \tilde V''\right)^T\nonumber\\
         &\quad-2\ric(\X,V')-\R\left(\X^T,\n,V',\n\right)-\left(\nabla_\X \ric\right)(\X,\n)\nonumber\\
         &\quad-\Delta^S\left \langle\Z,\n\right\rangle-\left\langle\Z,\n\right\rangle\left(|h|^2+\ric(\n,\n)\right)+\Z^T H,\nonumber
    \end{align}
    where 
    \begin{equation*}
        \left(\tilde V''\right)^T\coloneqq \sum_{j=1}^n\left(\R(\X,e_j,\n,\X^T)+2\left\langle \nabla^S\varphi- A\left( \X^T\right) ,\varphi A(e_j)+\nabla_{e_j}\ \X^T\right\rangle\right)e_j
    \end{equation*}
    In particular, if $\Phi$ is a normal variation,
    \begin{equation}\label{secvarhnorm}
        \begin{split}
           \left.\frac{d^2}{dt^2}\right|_{t=0}H_t\left(\Phi_t(p)\right)&=\divv^S \left( A\left(\nabla^S\varphi^2\right)\right)+2\varphi\divv^S A\left(\nabla^S\varphi\right)-2\varphi\left\langle \nabla^S H,\nabla^S\varphi\right\rangle-\left|\nabla^S\varphi\right|^2H\\
                &\quad+2\varphi^2\trace\left(h^3\right)-2\varphi^2\left\langle h,\j(\n,\n)\right\rangle-\varphi^2\left(\nabla_\n \ric\right)(\n,\n)\\
         &\quad-\Delta^S\left \langle\Z,\n\right\rangle-\left\langle\Z,\n\right\rangle\left(|h|^2+\ric(\n,\n)\right)+\Z^T H.
        \end{split}
    \end{equation}
\end{lemma}
\begin{proof}
Formula \eqref{derprimhinstatement} is well-known (cf. \cite[Lemma 1.26]{MR4676392}). We prove it for the sake of completeness. 
    Let $e_1,\ldots,e_n$ be any orthonormal basis of $T_p S$. 
    %Denote by $\k_1,\ldots,\k_n$ the corresponding principal curvatures. 
    For $j=1,\ldots,n$, let $E_j(t)$ be as in \eqref{estensionevettore}. Define $\S(t)$ by
    \begin{equation*}
       \S(t)_{ij}=\left\langle\left.\nabla_{E_i(t)}\N(t,\cdot)\right|_{\beta(t)},E_j(t)\right\rangle,\qquad i,j=1,\ldots,n.
    \end{equation*}
    Then $\S(t)$ is symmetric, and $       H(\Phi_t(p))=\trace\left(\G(t)^{-1}\S(t)\right),$
   where $\G$ is defined in \eqref{grammatrix}. 
   Recalling \eqref{derivatadimatriceinversa},
    \begin{equation*}
        \frac{d}{dt}H_t(\Phi_t(p))=\trace\left(-\G(t)^{-1}\frac{d\G(t)}{dt}\G(t)^{-1}\S(t)+\G(t)^{-1}\frac{d\S(t)}{dt}\right).
    \end{equation*}
Since $e_1,\ldots,e_n$ is orthonormal, then $\G(0)_{ij}=\delta_{ij}$ for $i,j=1,\ldots,n$, whence
\begin{align}
     \left.\frac{d}{dt}\right|_{t=0}H_t(\Phi_t(p))&=\trace\left(-\left(\left.\frac{d}{dt}\right|_{t=0}\G(t)\right)\S(0)+\left.\frac{d}{dt}\right|_{t=0}\S(t)\right),\label{unatracciaderiprimh}\\
      \left.\frac{d^2}{dt^2}\right|_{t=0}H_t(\Phi_t(p))&=\underbrace{\trace\left(2\left(\left.\frac{d}{dt}\right|_{t=0}\G(t)\right)^2\S(0)\right)}_{\mathrm{I}}+\underbrace{\trace\left(-\left(\left.\frac{d^2}{dt^2}\right|_{t=0}\G(t)\right)\S(0)\right)}_{\mathrm{II}}\label{quattrotraccedersech}\\
      &\quad+\underbrace{\trace\left(-2\left(\left.\frac{d}{dt}\right|_{t=0}\G(t)\right)\left(\left.\frac{d}{dt}\right|_{t=0}\S(t)\right)\right)}_{\mathrm{III}}+\underbrace{\trace\left(\left.\frac{d^2}{dt^2}\right|_{t=0}\S(t)\right)}_{\mathrm{IV}}\nonumber.
\end{align}
First, by \eqref{unatracciaderiprimh},
\begin{equation*}
    \begin{split}
         &\left.\frac{d}{dt}\right|_{t=0}H_t(\Phi_t(p))=-\sum_{i,j=1}^n\left(\left.\frac{d}{dt}\right|_{t=0}\G(t)_{ij}\right)\S(0)_{ij}+\sum_{i=1}^n\left.\frac{d}{dt}\right|_{t=0}\S(t)_{ii}\\
         &\qquad\overset{\eqref{prounoaux},\eqref{protreaux}}{=}-\sum_{i,j=1}^nh_{ij}\left(\left\langle\nabla_{e_i} \X,e_j\right\rangle+\left\langle \nabla_{e_j} \X,e_i\right\rangle\right)+\sum_{i=1}^n\left\langle \nabla_{e_i}V',e_i\right\rangle-\ric(\X,\n)+\sum_{i=1}^n\left\langle\nabla_{e_i}\n ,\nabla_{e_i}\X\right\rangle.
    \end{split}
\end{equation*}
Notice that, for $i,j=1,\ldots,n,$
\begin{equation}\label{pezzoconsyminproof}
    \begin{split}
        \left\langle\nabla_{e_i} \X,e_j\right\rangle+\left\langle \nabla_{e_j} \X,e_i\right\rangle&=2\varphi h_{ij}+ \left\langle\nabla_{e_i}  \X^T,e_j\right\rangle+\left\langle \nabla_{e_j}  \X^T,e_i\right\rangle=2\varphi h_{ij}+2\sym\left(\nabla^S  \X^T\right)_{ij}
    \end{split}
\end{equation}
and
\begin{equation}\label{auxprimadiprg}
    \left\langle\nabla_{e_i}V',e_j\right\rangle\overset{\eqref{lemma125ritore}}{=}-\left(\hess^S\varphi\right)_{ij}+\left\langle \nabla_{e_i}A\left( \X^T\right),e_j\right\rangle.
\end{equation}
Moreover,
\begin{equation}\label{solopernecessitàdisbatti}
    \begin{split}
        \sum_{i=1}^n\left\langle\nabla_{e_i}\n ,\nabla_{e_i}\X\right\rangle=\varphi\sum_{i=1}^n\left\langle\nabla_{e_i}\n ,\nabla_{e_i}\n \right\rangle+\sum_{i=1}^n\left\langle\nabla_{e_i}\n ,\nabla_{e_i} \X^T\right\rangle=\varphi|h|^2+\left\langle h,\sym\left(\nabla^S  \X^T\right)\right\rangle
    \end{split}
\end{equation}
Therefore, by \eqref{pezzoconsyminproof}, \eqref{auxprimadiprg} and \eqref{solopernecessitàdisbatti},
\begin{equation*}
     \left.\frac{d}{dt}\right|_{t=0}H_t(\Phi_t(p))=-\Delta^S\varphi-\varphi\left(|h|^2+\ric(\n,\n)\right)-\left\langle h,\sym\left(\nabla^S  \X^T\right)\right\rangle+\divv^SA\left( \X^T\right)-\ric( \X^T,N).
\end{equation*}
Next, extend $ \X^T$ to a smooth vector field with compact support in $M$. Denote by $\Psi_t:M\to M$ its flow. Then $\Psi$ is a variation of $S$ with velocity $ \X^T$. By the above formula,
\begin{equation*}
     \X^T H(p)=  \left.\frac{d}{dt}\right|_{t=0}H_t(\Psi_t(p))=-\left\langle h,\sym\left(\nabla^S  \X^T\right)\right\rangle+\divv^S A\left( \X^T\right)-\ric( \X^T,\n).
\end{equation*}
Therefore, \eqref{derprimhinstatement} follows.
%\footnote{questo stesso ragionamento (che ho preso dal libro di manuel) dovrebbe permettere di semplificare anche le parti esclusivamente tangenti della variazione seconda}
Next, we prove \eqref{dersechinstatement}. First we compute $\mathrm{I}$. Indeed,
\begin{equation*}
    \begin{split}
        \mathrm{I}&=2\sum_{i,j,k=1}^n\left(\left.\frac{d}{dt}\right|_{t=0}\G(t)_{ik}\right)\left(\left.\frac{d}{dt}\right|_{t=0}\G(t)_{jk}\right)\S(0)_{ij}\\
        \overset{\eqref{prounoaux}}&{=}2\sum_{i,j,k=1}^nh_{ij}\left(\left\langle\nabla_{e_i}\X,e_k\right\rangle+\left\langle\nabla_{e_k}\X,e_i\right\rangle\right)\left( \left\langle\nabla_{e_j}\X,e_k\right\rangle+\left\langle\nabla_{e_k}\X,e_j\right\rangle\right).
  %      &=8\sum_{i,j,k=1}^nh_{ij}\sym\left(\nabla^S \X\right)_{ik}\sym\left( \nabla^S\X\right)_{jk}\\
    %    &=8\left\langle h,\sym\left(\nabla^S \X\right)^2\right\rangle.
    \end{split}
\end{equation*}
%But
%\begin{equation*}
 %   \begin{split}
        %\sum_{k=1}^n\left(\left\langle\nabla_{e_i}X,e_k\right\rangle+\left\langle\nabla_{e_k}X,e_i\right\rangle\right)^2&=\sum_{k=1}^n\left\langle\nabla_{e_i}X,e_k\right\rangle^2+\sum_{k=1}^n\left\langle\nabla_{e_k}X,e_i\right\rangle^2+2\sum_{k=1}^n\left\langle\nabla_{e_i}X,e_k\right\rangle \left\langle\nabla_{e_k}X,e_i\right\rangle\\
 %       &=\sum_{k=1}^n\left(\varphi h_{ik}+\left\langle\nabla_{e_i} \X^T,e_k\right\rangle\right)^2+\sum_{k=1}^n\left(\varphi h_{ik}+\left\langle\nabla_{e_k} \X^T,e_i\right\rangle\right)^2\\
   %     &\quad+2\sum_{k=1}^n\left(\varphi h_{ik}+\left\langle\nabla_{e_i} \X^T,e_k\right\rangle\right)\left(\varphi h_{ik}+\left\langle\nabla_{e_k} \X^T,e_i\right\rangle\right)\\
    %    &\overset{\eqref{semplificazioneortoprinci}}{=}4\varphi^2\k_i^2+8\varphi\k_i\left\langle\nabla_{e_i} \X^T,e_i\right\rangle+\sum_{k=1}^n\left(\left\langle\nabla_{e_i} \X^T,e_k\right\rangle+\left\langle\nabla_{e_k} \X^T,e_i\right\rangle\right)^2, 
   % \end{split}
%\end{equation*}
%whence
%\begin{equation*}
 %   \mathrm{I}=8\varphi^2\trace\left(A^3\right)+16\varphi\sum_{i=1}^n\k_i^2 \left\langle\nabla_{e_i} \X^T,e_i\right\rangle+2\sum_{i=1}^n\k_i\sum_{k=1}^n\left(\left\langle\nabla_{e_i} \X^T,e_k\right\rangle+\left\langle\nabla_{e_k} \X^T,e_i\right\rangle\right)^2.
%\end{equation*}
Next we compute $\mathrm{II}$. Noticing that
\begin{equation}\label{componentidivprimoinproof}
    \left\langle\nabla_{e_i}\X,\n\right\rangle=e_i(\varphi)+\left\langle\nabla_{e_i}\X^T,\n\right\rangle=e_i(\varphi)-\left\langle\nabla_{e_i}\n,\X^T\right\rangle=-\left\langle V',e_i\right\rangle\qquad i=1,\ldots,n,
\end{equation}
we deduce that
\begin{equation*}
    \begin{split}
        \mathrm{II}&=-\sum_{i,j=1}^n\left(\left.\frac{d^2}{dt^2}\right|_{t=0}\G(t)_{ij}\right)\S(0)_{ij}\\
        &=-\sum_{i,j=1}^n h_{ij}\left.\frac{d^2}{dt^2}\right|_{t=0}\langle E_i(t),E_j(t)\rangle\\
        &=-2\sum_{i,j=1}^n h_{ij}\left.\frac{d}{dt}\right|_{t=0}\left\langle \frac{D}{dt}E_i(t),E_j(t)\right\rangle\\
        &=-2\sum_{i,j=1}^n h_{ij}\left\langle \left.\frac{D^2}{dt^2}\right|_{t=0}E_i(t),E_j(t)\right\rangle -2\sum_{i,j=1}^n h_{ij}\left\langle \left.\frac{D}{dt}\right|_{t=0}E_i(t),\left.\frac{D}{dt}\right|_{t=0}E_j(t)\right\rangle\\
        \overset{\eqref{prounoaux},\eqref{produeaux}}&{=}-2\sum_{i,j=1}^n h_{ij}\left\langle \nabla_{e_i}\Z,e_j\right\rangle-2\sum_{i,j=1}^n h_{ij}\R(\X,e_i,\X,e_j)-2\sum_{i,j=1}^n h_{ij}\left\langle\nabla_{e_i}\X,\nabla_{e_j}\X\right\rangle\\
        &=-2\left\langle h,\nabla^S\Z\right\rangle-2\left\langle h,\j(\X,\X)\right\rangle-2\sum_{i,j=1}^n h_{ij}\left\langle \nabla_{e_i}\X,e_k\right\rangle\left\langle \nabla_{e_j}\X,e_k\right\rangle-2\sum_{i,j=1}^n h_{ij}\left\langle \nabla_{e_i}\X,\n\right\rangle\left\langle \nabla_{e_j}\X,\n\right\rangle\\
        \overset{\eqref{componentidivprimoinproof}}&{=}-2\left\langle h,\nabla^S\Z\right\rangle-2\left\langle h,\j(\X,\X)\right\rangle-2\sum_{i,j=1}^n h_{ij}\left\langle \nabla_{e_i}\X,e_k\right\rangle\left\langle \nabla_{e_j}\X,e_k\right\rangle-2h\left(V',V'\right).\\
    \end{split}
\end{equation*}
%Notice that  
%\begin{equation*}
%\begin{split}
%     \mathrm{II.1}&=-2\sum_{i,j=1}^n h_{ij}\left(e_i(\varphi)-\left\langle \nabla_{e_i}\n,\X^T\right\rangle\right)   \left(e_j(\varphi)-\left\langle \nabla_{e_j}\n,\X^T\right\rangle\right) \\
 %    &=-2\sum_{i,j=1}^n h_{ij}e_i(\varphi)e_j(\varphi)+4\sum_{i,j,k=1}^nh_{ij}h_{jk}e_i(\varphi)\left(\X^T\right)^k-2\sum_{i,j,k,l=1}^nh_{ij}h_{ik}h_{jl}\left(\X^T\right)^k\left(\X^T\right)^l\\
 %    &=-2h\left(\nabla^S\varphi,\nabla^S\varphi\right)+4 h^2\left(\nabla^S\varphi,\X^T\right)-2h^3\left(\X^T,\X^T\right).
%\end{split}
%\end{equation*}
%Therefore
%\begin{equation*}
 %   \begin{split}
  %      \mathrm{II} =-2&\left\langle h,\nabla^S\Z\right\rangle-2\left\langle h,\j(\X,\X)\right\rangle-2\sum_{i,j=1}^n h_{ij}\left\langle \nabla_{e_i}\X,e_k\right\rangle\left\langle \nabla_{e_j}\X,e_k\right\rangle\\
  %      &\qquad-2h\left(\nabla^S\varphi,\nabla^S\varphi\right)+4 h^2\left(\nabla^S\varphi,\X^T\right)-2h^3\left(\X^T,\X^T\right)
   % \end{split}
%\end{equation*}
We compute $\mathrm{III}$. To this aim,
\begin{equation*}
    \begin{split}
        \mathrm{III}&=-2\sum_{i,j=1}^n\left(\left.\frac{d}{dt}\right|_{t=0}\G(t)_{ij}\right)\left(\left.\frac{d}{dt}\right|_{t=0}\S(t)_{ij}\right)\\
        \overset{\eqref{prounoaux},\eqref{protreaux}}&{=}-2\sum_{i,j=1}^n\left(\left\langle\nabla_{e_i}\X,e_j\right\rangle+\left\langle \nabla_{e_j}\X,e_i\right\rangle\right)\left(\left\langle\nabla_{e_i}V'+\R(\X,e_i)\n,e_j\right\rangle+\left\langle\nabla_{e_i}\n ,\nabla_{e_j}\X\right\rangle\right)\\
        &=-2\sum_{i,j=1}^n\left(\left\langle\nabla_{e_i}\X,e_j\right\rangle+\left\langle \nabla_{e_j}\X,e_i\right\rangle\right)\left\langle\nabla_{e_i}V',e_j\right\rangle-2\sum_{i,j,k=1}^nh_{ij}\left(\left\langle\nabla_{e_i}\X,e_k\right\rangle+\left\langle \nabla_{e_k}\X,e_i\right\rangle\right)\left\langle\nabla_{e_k}\X,e_j\right\rangle\\
    %   \overset{\eqref{pezzoconsyminproof}}&{=}-4\sum_{i,j=1}^n\sym\left(\nabla^S  \X\right)_{ij}\left(\left\langle\nabla_{e_i}V',e_j\right\rangle+\left\langle\nabla_{e_i}\n ,\nabla_{e_j}\X\right\rangle\right)\\
   %    &\quad-2\sum_{i,j=1}^n\left(\left\langle\nabla_{e_i}\X,e_j\right\rangle+\left\langle \nabla_{e_j}\X,e_i\right\rangle\right)\R(\X,e_i,\n,e_j)\\
       &\quad-2\left\langle \nabla^S \X,\j(\X,\n)\right\rangle-2\left\langle \nabla^S \X,\j(\n,\X)\right\rangle.
    \end{split}
\end{equation*}
Finally, we compute $\mathrm{IV}$. Indeed, by \eqref{produeaux}, \eqref{protreaux} and \eqref{proquattroaux}, 
\begin{equation*}
    \begin{split}
        \mathrm{IV}&=\sum_{i=1}^n\left.\frac{d^2}{dt^2}\right|_{t=0}\S(t)_{ii}\\
        &=\underbrace{\divv^S V''}_{\mathrm{IV.1}}+\underbrace{\sum_{i=1}^n\left\langle \R\left(\X,\nabla_{e_i}\X\right)\n+2\R(\X,e_i)V'+\R\left(\Z,e_i\right)\n+\left(\nabla_\X \R\right)(\X,e_i)\n,e_i\right\rangle}_{\mathrm{IV.2}}\\
        &\quad+\underbrace{2\sum_{i=1}^n\left\langle  \nabla_{e_i}V'+\R(\X,e_i)\n, \nabla_{e_i}\X\right\rangle}_{\mathrm{IV.3}}+\underbrace{\sum_{i=1}^n\left\langle \nabla_{e_i}\n , \nabla_{e_i}\Z+\R(\X,e_i)\X\right\rangle}_{\mathrm{IV.4}}.
    \end{split}
\end{equation*}
%Let $E_1,\ldots,E_n$ be a local geodesic frame centered at $p$. Define, close to $p$,
%\begin{equation*}
%  \tilde V''\coloneqq \sum_{j=1}^n\left(-\left \langle \n ,\nabla_{\E_i}\nabla_{\X ^T} \X^T+\R( \X^T,\E_i) \X^T\right\rangle-2\left\langle A\left( \X^T\right) ,\nabla_{\E_i}  \X^T\right\rangle\right)\E_i.
%\end{equation*}
First, 
\begin{equation*}
    \begin{split}
        \mathrm{IV.1}\overset{\eqref{secondordernormalbehave}}&{=}\left(-\left|\nabla^S\varphi\right|^2+2h\left(\nabla^S\varphi, \X^T\right)-h^2\left( \X^T, \X^T\right)\right)H+\divv^S\left( V''\right)^T.
    \end{split}
\end{equation*}
Moreover,
\begin{equation*}
    \begin{split}
        \mathrm{IV.2}&= \sum_{i,j=1}^n\left\langle\nabla_{e_i}\X,e_j\right\rangle \R(\X,e_j,\n,e_i)+\sum_{i=1}^n\left\langle\nabla_{e_i}\X,\n\right\rangle \R(\X,\n,\n,e_i)\\
        &\quad-2\ric(\X,V')-2 \R(\X,\n,V',\n)-\ric(\Z,\n)-\left(\nabla_\X \ric\right)(\X,\n)\\
   %     &=\left \langle\nabla^S \X,\j(\n,\X)\right\rangle+\R\left(\X,\n,\n,\nabla^S\varphi\right)-\R\left(\X,\n,\n,A(\X^T)\right)\\
     %   &\quad -2\ric(\X,V')-2 \R(\X,\n,V',\n)-\ric(\Z,\n)-\left(\nabla_\X R\right)(\X,\n)\\
        \overset{\eqref{componentidivprimoinproof}}&{=}\left \langle\nabla^S \X,\j(\n,\X)\right\rangle-2\ric(\X,V')- \R(\X^T,\n,V',\n)-\ric(\Z,\n)-\left(\nabla_\X \ric\right)(\X,\n).
      %  &\overset{\eqref{secondordernormalbehave}}{=}\left(-\left|\nabla^S\varphi\right|^2+2h\left(\nabla^S\varphi, \X^T\right)-h^2\left( \X^T, \X^T\right)\right)H+\divv^S\left(V''-\tilde V''\right)^T+\divv^S\left(\tilde V''\right)^T\\
    \end{split}
\end{equation*}
In addition,
\begin{equation*}
    \begin{split}
        \mathrm{IV.3}&=2\sum_{i,j=1}^n\left\langle\nabla_{e_i}V',e_j\right\rangle\left\langle\nabla_{e_i}\X,e_j\right\rangle+2\sum_{i=1}^n\left\langle\nabla_{e_i}V',\n\right\rangle\left\langle\nabla_{e_i}\X,\n\right\rangle+2\sum_{i,j=1}^n\R(\X,e_i,\n,e_j)\left\langle\nabla_{e_i}\X,e_j\right\rangle\\
  %      &=2\left\langle\nabla^S\X,\nabla^S V'\right\rangle-2\sum_{i=1}^n\left\langle\nabla_{V'}\n,e_i\right\rangle\left( e_i(\varphi)-\left\langle\nabla_{\X^T}\n,e_i\right\rangle\right)+2\left \langle\nabla^S X,\j(\X,\n)\right\rangle\\
    %    &=2\left\langle\nabla^S\X,\nabla^S V'\right\rangle-2h\left(\nabla^S\varphi,V'\right)+2h^2\left(\X^T,V'\right)+2\left \langle\nabla^S X,\j(\X,\n)\right\rangle\\
        \overset{\eqref{componentidivprimoinproof}}&{=}2\sum_{i,j=1}^n\left\langle\nabla_{e_i}V',e_j\right\rangle\left\langle\nabla_{e_i}\X,e_j\right\rangle+2h\left(V',V'\right)+2\left \langle\nabla^S X,\j(\X,\n)\right\rangle
    \end{split}
\end{equation*}
Finally,
\begin{equation*}
    \begin{split}
        \mathrm{IV.4}&=\sum_{i,j=1}^nh_{ij}\left\langle \nabla_{e_i}\Z,e_j\right\rangle+\sum_{i,j=1}^nh_{ij}\R(\X,e_i,\X,e_j)=\left\langle h,\nabla^S\Z\right\rangle+\left\langle h,\j(\X,\X)\right\rangle.
    \end{split}
\end{equation*}
Recalling \eqref{quattrotraccedersech} and combining the above computations,
\begin{equation*}
    \begin{split}
        \left.\frac{d^2}{dt^2}\right|_{t=0}H_t\left(\Phi_t(p)\right)&=2\left\langle h,\left(\nabla^S X\right)^2\right\rangle-2\left\langle \nabla^S X,\left(\nabla^SV'\right)^t\right\rangle-\left\langle h,\j(\X,\X)\right\rangle-\left\langle\nabla^S \X,\j(\n,\X)\right\rangle\\
         &\quad+\left(-\left|\nabla^S\varphi\right|^2+2h\left(\nabla^S\varphi, \X^T\right)-h^2\left( \X^T, \X^T\right)\right)H+\divv^S\left( V''\right)^T\\
         &\quad-2\ric(\X,V')-\R\left(\X^T,\n,V',\n\right)-\left(\nabla_\X \ric\right)(\X,\n)-\left\langle h,\nabla^S \Z\right\rangle-\ric(\Z,\n).
    \end{split}
\end{equation*}
 Next, fix $p\in S$ and a geodesic frame $\E_1,\ldots,\E_n$ at $p$. Recall that 
    \begin{equation*}
        \left(V''\right)^T=-\sum_{i=1}^n\left \langle \n ,\nabla_{\E_i}\Z\right\rangle\E_i+\left(\tilde V''\right)^T.
    \end{equation*}
    Therefore
    \begin{equation*}
        \begin{split}
            \divv^S\left(V''\right)^T&=-\sum_{i=1}^n\E_i\left \langle \n ,\nabla_{\E_i}\Z\right\rangle+\divv^S \left(\tilde V''\right)^T\\
            &=-\sum_{i=1}^n\E_i\left \langle \n ,\nabla_{\E_i}\left(\left\langle\Z,\n\right\rangle\n\right)\right\rangle-\sum_{i=1}^n\E_i\left \langle \n ,\nabla_{\E_i}\Z^T\right\rangle+\divv^S \left(\tilde V''\right)^T\\
            &=-\sum_{i=1}^n\E_i\left(\E_i\left \langle\Z,\n\right\rangle\right)+\sum_{i=1}^n\E_i\left \langle A\left(\Z^T\right) ,\E_i\right\rangle+\divv^S \left(\tilde V''\right)^T\\
            &=-\Delta^S\left \langle\Z,\n\right\rangle+\divv^S A\left(\Z^T\right)+\divv^S \left(\tilde V''\right)^T.
        \end{split}
    \end{equation*}
Moreover, 
      \begin{equation*}
        \begin{split}
            -\left\langle h,\nabla^S \Z\right\rangle-\ric(\Z,\n)\overset{\eqref{corollarioditracedcodazzi}}&{=}-\divv^S A\left(\Z^T\right)+\Z^T H-\left\langle\Z,\n\right\rangle\left(|h|^2+\ric(\n,\n)\right).
        \end{split}
    \end{equation*}
    Therefore
%\begin{equation*}
 %   \begin{split}
  %      \left.\frac{d^2}{dt^2}\right|_{t=0}H_t\left(\Phi_t(p)\right)&=2\left\langle h,\left(\nabla^S X\right)^2\right\rangle-2\left\langle \nabla^S X,\left(\nabla^SV'\right)^t\right\rangle-\left\langle h,\j(\X,\X)\right\rangle-\left\langle\nabla^S \X,\j(\n,\X)\right\rangle\\
     %    &\quad+\left(-\left|\nabla^S\varphi\right|^2+2h\left(\nabla^S\varphi, \X^T\right)-h^2\left( \X^T, \X^T\right)\right)H+\divv^S\left( \tilde V''\right)^T\\
     %    &\quad-2\ric(\X,V')-R\left(\X^T,\n,V',\n\right)-\left(\nabla_\X \ric\right)(\X,\n)\\
    %     &\quad-\Delta^S\left \langle\Z,\n\right\rangle-\left\langle\Z,\n\right\rangle\left(|h|^2+\ric(\n,\n)\right)+\Z^T H,
  %  \end{split}
%\end{equation*}
%and 
\eqref{dersechinstatement} follows. Finally, assume that $\Phi$ is normal. Then $\X=\varphi\n$, $\X^T=0$ and $V'=-\nabla^S\varphi$, so that
\begin{equation*}
    \begin{split}
         \left.\frac{d^2}{dt^2}\right|_{t=0}&H_t\left(\Phi_t(p)\right)=2\left\langle h,\left(\nabla^S \left(\varphi\n\right)\right)^2\right\rangle+2\left\langle \nabla^S \left(\varphi\n\right),\left(\nabla^S\left(\nabla^S\varphi\right)\right)^t\right\rangle-\varphi^2\left\langle h,\j(\n,\n)\right\rangle\\
         &\quad-\varphi\left\langle\nabla^S \left(\varphi\n\right),\j(\n,\n)\right\rangle-\left|\nabla^S\varphi\right|^2H+\divv^S\left( \tilde V''\right)^T+2\varphi\ric(\n,\nabla^S\varphi)-\varphi^2\left(\nabla_\n \ric\right)(\n,\n)\\
         &\quad -\Delta^S\left \langle\Z,\n\right\rangle-\left\langle\Z,\n\right\rangle\left(|h|^2+\ric(\n,\n)\right)+\Z^T H.
    \end{split}
\end{equation*}
Noticing that $\nabla^S\left(\varphi\n\right)=\varphi h$ and $\left(\nabla^S\left(\nabla^S\varphi\right)\right)^t=\hess^S\varphi$, we deduce that
\begin{equation*}
        \begin{split}
                \left.\frac{d^2}{dt^2}\right|_{t=0}&H_t\left(\Phi_t(p)\right)=2\varphi^2\trace\left(h^3\right)+2\varphi\left\langle h,\hess^S\varphi\right\rangle-2\varphi^2\left\langle h,\j(\n,\n)\right\rangle-\left|\nabla^S\varphi\right|^2H+\divv^S\left( \tilde V''\right)^T\\
         &\quad+2\varphi\ric\left(\n,\nabla^S\varphi\right)-\varphi^2\left(\nabla_\n \ric\right)(\n,\n)-\Delta^S\left \langle\Z,\n\right\rangle-\left\langle\Z,\n\right\rangle\left(|h|^2+\ric(\n,\n)\right)+\Z^T H.
        \end{split}
    \end{equation*}
    Moreover, 
    \begin{equation*}
    \begin{split}
          2\varphi\left\langle h,\hess^S\varphi\right\rangle+2\varphi\ric\left(\n,\nabla^S\varphi\right)\overset{\eqref{corollarioditracedcodazzi}}&{=}2\varphi\divv^S A\left(\nabla^S\varphi\right)-2\varphi\left\langle \nabla^S H,\nabla^S\varphi\right\rangle.
      %    &=2\divv^S \left(\varphi A\left(\nabla^S\varphi\right)\right)-2 h\left(\nabla^S\varphi,\nabla^S\varphi\right)-\left\langle \nabla^S H,\nabla^S\left(\varphi^2\right)\right\rangle\\
  %        &=\divv^S \left( A\left(\nabla^S\left(\varphi^2\right)\right)\right)-\divv^S\left(\varphi^2\nabla^SH\right)-2 h\left(\nabla^S\varphi,\nabla^S\varphi\right)+\varphi^2\Delta^S H.
    \end{split}
    \end{equation*}
    Finally, since
    \begin{equation*}
        \left(\tilde V''\right)^T=2\sum_{i=1}^n\left\langle \nabla^S\varphi,\varphi A(e_i)\right\rangle e_i=\sum_{i=1}^n\left\langle \nabla^S\left(\varphi^2\right), A(e_i)\right\rangle e_i=\sum_{i=1}^n\left\langle  A\left(\nabla^S\left(\varphi^2\right)\right), e_i\right\rangle e_i=A\left(\nabla^S\left(\varphi^2\right)\right),
    \end{equation*}
    \eqref{secvarhnorm} follows.
%    we conclude that 
  %  \begin{equation*}
 %       \begin{split}
   %             \left.\frac{d^2}{dt^2}\right|_{t=0}H_t\left(\Phi_t(p)\right)&=\divv^S \left( A\left(\nabla^S\varphi^2\right)\right)+2\varphi\divv^S A\left(\nabla^S\varphi\right)-2\varphi\left\langle \nabla^S H,\nabla^S\varphi\right\rangle-\left|\nabla^S\varphi\right|^2H\\
   %             &\quad+2\varphi^2\trace\left(h^3\right)-2\varphi^2\left\langle h,\j(\n,\n)\right\rangle-\varphi^2\left(\nabla_\n \ric\right)(\n,\n)\\
   %      &\quad-\Delta^S\left \langle\Z,\n\right\rangle-\left\langle\Z,\n\right\rangle\left(|h|^2+\ric(\n,\n)\right)+\Z^T H.
   %     \end{split}
  %  \end{equation*}
\end{proof}
\subsection{Variation formulas}
 If 
$f:\rr\to\rr$ is smooth in a neighborhood of $\left\{H(p)\,:\,p\in S\right\}$, set
\begin{equation}\label{genfunriemappend}
    \tmc_f(S)=\int_S f(H)\,dS.
\end{equation}
Let $\Phi$ be a variation. The area formula implies that 
\begin{equation}\label{areaformulaapplicata}
     \tmc_f\left(\Phi_t(S)\right)=\int_S f\left(H_t\left(\Phi_t(p)\right)\right)\jac\left(\Phi_t|_{S}\right)\left(p\right)\,dS.
\end{equation}
%In particular, \Cref{lemmaderivatejacobiano} and \Cref{lemmavariazioneprimacurvaturaeseconda} grant that the first variation of $\tmc_f$ at $S$ along $\Phi$ depends only on $\X|_S$, while the second variation depends only on $\X|_S$ and on $\Z|_S$. In order to emphasize this dependence, we adopt the notation
%\begin{equation*}
 %%   \delta\tmc_f(S)[\X]\coloneqq \left.\frac{d}{dt}\right|_{t=0}\tmc\left(\Phi_t(S)\right),\qquad   \delta^2\tmc_f(S)[\X,\Z]\coloneqq \left.\frac{d^2}{dt^2}\right|_{t=0}\tmc\left(\Phi_t(S)\right).
%\end{equation*}
Denote by $\jacobi$ the \emph{Jacobi operator} associated to $S$, 
\begin{equation}\label{def_jacobi_operator}
    \jacobi \varphi=-\Delta^S\varphi-\varphi\left(|h|^2+\ric(\n,\n)\right),\qquad\varphi\in C^\infty(S).
\end{equation}
Recall that $\jacobi$ is self-adjoint, namely
\begin{equation}\label{jacobiselfadjoint}
    \int_S\varphi\jacobi\psi\,dS=\int_S\psi\jacobi\varphi\,dS,\qquad\varphi,\psi\in C^\infty(S).
\end{equation}
The first and second variation formulas for \eqref{genfunriemappend} read as follows.
\begin{theorem}\label{maintheoremappendixA}
    Let $S\subseteq M$ be a smooth, embedded, closed, two-sided hypersurface. Fix a function $f:\rr\to\rr$ which is smooth in a neighborhood of $\left\{H(p)\,:\,p\in S\right\}$. Let $\Phi$ be a variation. Denote by $\X$ and $\Z$ its velocity and acceleration respectively. On $S$, decompose $\X$ as $\X|_S=\varphi \n+\X^T$. Set 
        \begin{equation*}
 \delta\tmc_f(S)[\Phi]\coloneqq \left.\frac{d}{d t}\right|_{t=0}\tmc_f\left(\Phi_t(S)\right),\qquad   \delta^2\tmc_f(S)[\Phi]\coloneqq \left.\frac{d^2}{dt^2}\right|_{t=0}\tmc_f\left(\Phi_t(S)\right).
\end{equation*}     
    Then 
    \begin{align}
        \delta\tmc_f(S)[\varphi]\coloneqq  \delta\tmc_f(S)[\Phi]&=\int_S\varphi\Big(-\Delta^S f'(H)-f'(H)\left(|h|^2+\ric(\n,\n)\right)+ f(H) H\Big)\,dS,\label{primavarfunzgen}\\
        \delta^2\tmc_f(S)[\Phi]&=\delta\tmc_f(S)\left[\left\langle\Z|_S,\n\right\rangle-2\X^T\varphi+h\left(\X^T,\X^T\right)\right]+\int_S\varphi\,\mathcal L\varphi\,dS\label{variazsecondanormalemaarbitrariaperilresto},
    \end{align}
    where
    \begin{equation*}
        \begin{split}
            \mathcal L\varphi&=\jacobi\left(f''(H)\jacobi\varphi\right)+2f'(H)\divv^S A\left(\nabla^S\varphi\right)-\left(f'(H)H+f(H)\right)\Delta^S\varphi\\
            &\quad-2 f''(H)\left\langle\nabla^S H,A\left(\nabla^S\varphi\right)\right\rangle+\left(f''(H)H-2f'(H)\right)\left\langle \nabla^S H,\nabla^S\varphi\right\rangle\\
            &\quad+\varphi f'(H)\left(2\trace\left(h^3\right)-2\left\langle h,\j(\n,\n)\right\rangle-\left(\nabla_\n \ric\right)(\n,\n)\right)\\
            &\quad+\varphi \left ( f(H)H^2-\left(2f'(H) H+f(H)\right)\left(|h|^2+ \ric(\n,\n)\right)\right).
        \end{split}
    \end{equation*}
\end{theorem}
\begin{proof}
    First,
    \begin{equation*}
        \begin{split}
              &\frac{d\tmc_f\left(\Phi_t(S)\right)}{dt}\overset{\eqref{areaformulaapplicata}}{=}\int_S\frac{d}{dt}\Big(f\left(H_t\left(\Phi_t(p)\right)\right)\jac\left(\Phi_t|_{S}\right)\left(p\right)\Big)\,dS\\
              &\qquad=\int_S \left(f'\left(H_t\left(\Phi_t(p)\right)\right)\frac{d H_t\left(\Phi_t(p)\right)}{dt}\jac\left(\Phi_t|_{S}\right)\left(p\right)+f\left(H_t\left(\Phi_t(p)\right)\right)\frac{d \jac\left(\Phi_t|_{S}\right)\left(p\right)}{dt}\right)\,dS.\\
        \end{split}
    \end{equation*}
    In particular, by \eqref{derijacuno} and \eqref{derprimhinstatement},
    \begin{equation*}
        \delta\tmc_f(S)[\Phi]=\int_S\Big(f'(H)\left(-\Delta^S\varphi-\varphi\left(|h|^2+\ric(\n,\n)\right)+ \X^T H\right)+f(H)\left(\varphi H+\divv^S\X^T\right)\Big)\,dS.
    \end{equation*}
    Notice that, being $S$ closed,
    \begin{equation}\label{teoremadivergenzainproofvariazioneseconda}
        \int_S\Big( f'(H)\X^T H+f(H)\divv^S\X^T\Big)\,dS=\int_S\divv^S\left(f(H)\X^T\right)\,dS=0. 
    \end{equation}
The divergence theorem and \eqref{teoremadivergenzainproofvariazioneseconda} imply \eqref{primavarfunzgen}. We prove \eqref{variazsecondanormalemaarbitrariaperilresto}.
    We reduce to deal with the case of normal variations. Let $\W\in\Gamma(TM)$ be such that $\W|_S=\X^T$. Denote by $\Psi$ the flow of $-\W$. Since $\W|_S\in\Gamma(TS)$, then $\Psi(t,S)=S$ for any small $t$. Set $\tilde\Phi(t,p)\coloneqq\Phi(t,\Psi(t,p))$. Then $\tilde \Phi$ is a smooth variation, and moreover
    \begin{equation}\label{auxfromnormaltogeneral1}
        \tilde S_t\coloneqq\tilde\Phi(t,S)=\Phi(t,\Psi(t,S))=\Phi(t,S)=S_t
    \end{equation}
    for any $t$ sufficiently small. The area formula and \eqref{auxfromnormaltogeneral1} imply that 
    \begin{equation}\label{auxfromnormaltogeneraltre}
        \delta^2\tmc_f(S)[\Phi]=\delta^2\tmc_f(S)[\tilde \Phi].
    \end{equation} We compute the velocity $\tilde X$ and the acceleration $\tilde Z$ of $\tilde\Phi$. Fix $p\in M$. Fix local normal coordinates $x_1,\ldots,x_{n+1}$ centered at $p$. Set $q(t)=\Phi_t^{-1}(p)$ and $r(t)=\Psi^{-1}_t(q(t)).$ Then, since $\Psi$ is the flow of $-\W$, 
    \begin{equation}\label{auxfromnormaltogeneral2}
        \begin{split}
            \tilde\Y(t,p)&=\left.\frac{\partial}{\partial s}\right|_{s=0}\Phi(t+s,\Psi(t+s,r(t)))\\
            &=\sum_{i=1}^{n+1}\frac{\partial\Phi^i}{\partial t}(t,q(t))\frac{\partial}{\partial x_i}+\sum_{i,j=1}^{n+1}\frac{\partial\Phi^i}{\partial x_j}(t,q(t))\frac{\partial\Psi^j}{\partial t}(t,r(t))\frac{\partial}{\partial x_i}\\
            &=\Y(t,p)-\sum_{i,j=1}^{n+1}\frac{\partial\Phi^i}{\partial x_j}(t,q(t))\W(q(t))\frac{\partial}{\partial x_i}.
        \end{split}
    \end{equation}
    In particular, $\tilde X=\X-\W$, and $\tilde \X|_S=\varphi \n$, whence $\tilde \Phi$ is a normal variation. Moreover, as $   \Phi\left(t,q(t)\right)=p$,
\begin{equation*}
    \begin{split}
        0        %=\left.\frac{\partial}{\partial \tau}\right|_{\tau=0}\left(   \Phi\left(\tau,\Phi_\tau^{-1}(p)\right)\right)\\
=\frac{\partial\Phi}{\partial t}(0,p)+\sum_{i,j=1}^{n+1}\frac{\partial\Phi_i}{\partial z_j}(0,p)\dot q^j(0)\frac{\partial}{\partial x_i}        =\frac{\partial\Phi}{\partial t}(0,p)+\dot q(0),
    \end{split}
\end{equation*}
whence $\dot q(0)=-\X(p)$. Therefore

    
    \begin{equation*}
        \begin{split}
            \tilde \X'(p)\overset{\eqref{auxfromnormaltogeneral2}}&{=}\X'(p)-\sum_{i,j,k=1}^{n+1}\left(\delta_{ik}\frac{\partial^2\Phi^i}{\partial t\partial x_j}(0,p)\W^j(p)+\frac{\partial^2\Phi^i}{\partial x_k\partial x_j}(0,p)\dot q^k(0)\W^j(p)+\frac{\partial\Phi^i}{\partial x_j}(0,p)\frac{\partial\W^j}{\partial x_k}(p)\dot q^k(0)\right)\frac{\partial}{\partial x_i}\\
            &=\X'(p)-\sum_{i,j=1}^{n+1}\frac{\partial\X^i}{\partial x_j}(p)\W^j(p)\frac{\partial}{\partial x_i}+\sum_{i,k=1}^{n+1}\frac{\partial\W^i}{\partial x_k}(p)\X^k(p)\frac{\partial}{\partial x_i}\\
            &=\X'(p)-\left(\nabla_\W\X\right)(p)+\left(\nabla_\X\W\right)(p).
        \end{split}
    \end{equation*}
    Finally,
    \begin{equation*}
        \tilde\Z=\tilde\X'+\nabla_{\tilde\X}\tilde\X=\X'-\nabla_\W\X+\nabla_\X\W+\nabla_{\X-\W}\left(\X-\W\right)=\Z-2\nabla_\W\X+\nabla_\W\W,
    \end{equation*}
    so that, recalling that $\W|_S=\X^T$ ,
    \begin{equation}\label{auxfromnormaltogeneral4}
        \left\langle\tilde\Z|_S,\n\right\rangle=\left\langle\Z|_S,\n\right\rangle-2\X^T\varphi+h\left(\X^T,\X^T\right).
    \end{equation}   
    Then, by \eqref{derijacuno}, \eqref{derjduenormal},  \eqref{derprimhinstatement}, \eqref{secvarhnorm} and \eqref{auxfromnormaltogeneraltre},
    \begin{equation*}
        \begin{split}
            \delta^2\tmc_f(S)[\Phi]&=\int_Sf''(H)\left(\Delta^S\varphi+\varphi\left(|h|^2+\ric(\n,\n)\right)\right)^2\,dS\\
            &\quad+\int_Sf'(H)\Big( \divv^S \left( A\left(\nabla^S\varphi^2\right)\right)+2\varphi\divv^S A\left(\nabla^S\varphi\right)-2\varphi\left\langle \nabla^S H,\nabla^S\varphi\right\rangle-\left|\nabla^S\varphi\right|^2H\Big)\,dS\\
            &\quad+\int_Sf'(H)\Big(2\varphi^2\trace\left(h^3\right)-2\varphi^2\left\langle h,\j(\n,\n)\right\rangle-\varphi^2\left(\nabla_\n \ric\right)(\n,\n)\Big)\,dS\\
            &\quad+\int_Sf'(H)\Big(-\Delta^S\left \langle\tilde\Z,\n\right\rangle-\left\langle\tilde\Z,\n\right\rangle\left(|h|^2+\ric(\n,\n)\right)+\tilde\Z^T H\Big)\,dS\\
            &\quad+\int_S 2\varphi Hf'(H)\Big(-\Delta^S\varphi-\varphi\left(|h|^2+\ric(\n,\n)\right)\Big)\\
            &\quad+\int_Sf(H)\Big (\varphi^2H^2-\varphi^2|h|^2-\varphi^2\ric(\n,\n)+\left|\nabla^S\varphi\right|^2+\divv^S\tilde\Z\Big)\,dS.
        \end{split}
    \end{equation*}
    First, combining \eqref{primavarfunzgen} and \eqref{teoremadivergenzainproofvariazioneseconda}, and recalling \eqref{def_jacobi_operator},
        \begin{equation*}
        \begin{split}
            \delta^2&\tmc_f(S)[\Phi]=\underbrace{\int_S\left(f''(H)\jacobi\varphi\right)\jacobi\varphi\,dS}_\mathrm{I}\\
            &\quad+\underbrace{\int_Sf'(H)\Big( \divv^S \left( A\left(\nabla^S\varphi^2\right)\right)-\left|\nabla^S\varphi\right|^2H\Big)\,dS}_\mathrm{II}+\int_Sf'(H)\Big(2\varphi\divv^S A\left(\nabla^S\varphi\right)-2\varphi\left\langle \nabla^S H,\nabla^S\varphi\right\rangle\Big)\,dS\\
            &\quad+\int_Sf'(H)\Big(2\varphi^2\trace\left(h^3\right)-2\varphi^2\left\langle h,\j(\n,\n)\right\rangle-\varphi^2\left(\nabla_\n \ric\right)(\n,\n)\Big)\,dS\\
            &\quad+\int_S 2\varphi Hf'(H)\Big(-\Delta^S\varphi-\varphi\left(|h|^2+\ric(\n,\n)\right)\Big)\\
            &\quad+\underbrace{\int_Sf(H)\left|\nabla^S\varphi\right|^2\,dS}_{\mathrm{III}}+\int_Sf(H)\Big (\varphi^2H^2-\varphi^2|h|^2-\varphi^2\ric(\n,\n)\Big)\,dS+\delta\tmc_f(S)\left[\left\langle\tilde\Z|_S,\n\right\rangle\right].
        \end{split}
    \end{equation*}
    First,
    \begin{equation*}
        \mathrm{I}\overset{\eqref{jacobiselfadjoint}}{=}\int_S\varphi\,\jacobi\left(f''(H)\jacobi\varphi\right)\,dS.
    \end{equation*}
    Moreover,
    \begin{equation*}
 \begin{split}
     \mathrm{II}&=\int_S\Big(-\left\langle\nabla^S f'(H),A\left(\nabla^S\varphi^2\right)\right\rangle-\left\langle f'(H) H\nabla^S\varphi,\nabla^S\varphi\right\rangle\Big)\,dS\\
     &=\int_S\varphi\Big(-2 f''(H)\left\langle\nabla^S H,A\left(\nabla^S\varphi\right)\right\rangle+\divv^S\left(f'(H)H\nabla^S\varphi\right)\Big)\,dS\\
     &=\int_S\varphi\Big(-2 f''(H)\left\langle\nabla^S H,A\left(\nabla^S\varphi\right)\right\rangle+\left(f'(H)+f''(H)H\right)\left\langle \nabla^S H,\nabla^S\varphi\right\rangle+f'(H)H\Delta^S\varphi\Big)\,dS.\\
 \end{split}       
    \end{equation*}
    Finally,
    \begin{equation*}
        \begin{split}
            \mathrm{III}&=-\int_S \varphi\divv^S\left( f(H)\nabla^S\varphi\right)\,dS=-\int_S\varphi \Big(f'(H)\left\langle\nabla^S H,\nabla^S\varphi\right\rangle+f(H)\Delta^S\varphi\Big)\,dS.
        \end{split}
    \end{equation*}
    In particular,
    \begin{equation*}
        \mathrm{II}+\mathrm{III}=\int_S\varphi\Big(-2 f''(H)\left\langle\nabla^S H,A\left(\nabla^S\varphi\right)\right\rangle+f''(H)H\left\langle \nabla^S H,\nabla^S\varphi\right\rangle+\left(f'(H)H-f(H)\right)\Delta^S\varphi\Big)\,dS.
    \end{equation*}
The thesis follows combining the above computations with \eqref{auxfromnormaltogeneral4}.
 %   In conclusion,
  %   \begin{equation*}
  %      \begin{split}
  %          \delta^2\tmc_f(S)[X,Z]&=\delta\tmc_f(S)\left[\left\langle\Z,\n\right\rangle\right]+\int_S\varphi\,\jacobi\left(f''(H)\jacobi\varphi\right)\,dS\\
  %          &\quad+\int_S\varphi\Big(2f'(H)\divv^S A\left(\nabla^S\varphi\right)-\left(f'(H)H+f(H)\right)\Delta^S\varphi\Big)\,dS\\
   %         &\quad+\int_S\varphi\Big(-2 f''(H)\left\langle\nabla^S H,A\left(\nabla^S\varphi\right)\right\rangle+\left(f''(H)H-2f'(H)\right)\left\langle \nabla^S H,\nabla^S\varphi\right\rangle\Big)\,dS\\
    %        &\quad+\int_S\varphi ^2f'(H)\Big(2\trace\left(h^3\right)-2\left\langle h,\j(\n,\n)\right\rangle-\left(\nabla_\n \ric\right)(\n,\n)\Big)\,dS\\
    %        &\quad+\int_S\varphi^2 \Big ( f(H)H^2-\left(2f'(H) H+f(H)\right)\left(|h|^2+ \ric(\n,\n)\right)\Big)\,dS.
     %   \end{split}
 %   \end{equation*}
    \end{proof}
%    \subsection{Total mean curvature} When $f(H)=H$, $\tmc\coloneqq\tmc_f$ is the total mean curvature of $S$. In this case
% \begin{equation*}
 %       \begin{split}
  %          \mathcal L\varphi&=2\divv^S \Big(A\left(\nabla^S\varphi\right)-2H\nabla^S\varphi\Big)+\varphi \left(2\trace\left(h^3\right)-2\left\langle h,\j(\n,\n)\right\rangle-\left(\nabla_\n \ric\right)(\n,\n)\right)\\
     %       &\quad+\varphi \left ( H^3-3H\left(|h|^2+ \ric(\n,\n)\right)\right)\\
      %      &=2\divv^S \Big(A\left(\nabla^S\varphi\right)-2H\nabla^S\varphi\Big)+\varphi \Big(2\trace\left(h^3\right)-3H|h|^2+H^3\Big)\\
  %          &\quad-\varphi \Big (3H\ric(\n,\n)+2\left\langle h,\j(\n,\n)\right\rangle+\left(\nabla_\n \ric\right)(\n,\n) \Big).\\
      %  \end{split}
  %  \end{equation*}
  %  \begin{lemma}\label{lemmaaccatre}
 %       For any $k=1,\ldots,n$, denote by $H_k$ the mean curvature of order $k$. Namely, if $p\in S$ and $\k_1,\ldots,\k_n$ are the principal directions at $p$,
   %     \begin{equation*}
      %      H_k(p)\coloneqq \sum_{i_1<\ldots<i_k=1}^n\k_{i_1}\cdots\k_{i_k}.
    %    \end{equation*}
     %   Then
     %   \begin{equation*}
      %      2\trace\left(h^3\right)-3H|h|^2+H^3=6 H_3.
     %   \end{equation*}
  %  \end{lemma}
  %  \begin{proof}
    %    By a direct computation,
    %    \begin{equation*}
   %         \begin{split}
   %             2\trace\left(h^3\right)-3H|h|^2+H^3&=2\left(\trace\left(h^3\right)-H|h|^2\right)+H\left(H^2-|h|^2\right)\\
      %          &=2\left(\sum_{i=1}^n\k_i^3-\left(\sum_{i=1}^n\k_i\right)\left(\sum_{j=1}^n\k_j^2\right)\right)+\left(\sum_{i=1}^n\k_i\right)\left(\left(\sum_{j=1}^n\k_j\right)^2-\sum_{j=1}^n\k_j^2\right)\\
       %         &-2\sum_{\substack{i,j=1 \\ i\neq j}}^n\k_i\k_j^2+\left(\sum_{i=1}^n\k_i\right)\left(\sum_{\substack{j,k=1 \\ j\neq k}}^n\k_j\k_k\right)\\
         %       &-2\sum_{\substack{i,j=1 \\ i\neq j}}^n\k_i\k_j^2+\sum_{\substack{i,j,k=1 \\ i\neq j,i\neq k,j\neq k}}^n\k_i\k_j\k_k+\sum_{\substack{j,k=1 \\ j\neq k}}^n\k_j^2\k_k+\sum_{\substack{j,k=1 \\ j\neq k}}^n\k_j\k_k^2\\
        %        &=\sum_{\substack{i,j,k=1 \\ i\neq j,i\neq k,j\neq k}}^n\k_i\k_j\k_k\\
        %        &=6H_3.
         %   \end{split}
      %  \end{equation*}
  %  \end{proof}
  %  In particular, by \Cref{lemmaaccatre},
   %  \begin{equation*}
   %     \begin{split}
   %         \mathcal L\varphi&=2\divv^S \Big(A\left(\nabla^S\varphi\right)-2H\nabla^S\varphi\Big)+\varphi \Big(6H_3-3H\ric(\n,\n)-2\left\langle h,\j(\n,\n)\right\rangle-\left(\nabla_\n \ric\right)(\n,\n) \Big).\\
     %   \end{split}
 %   \end{equation*}
\section{Variation formulas for sub-Riemannian mean curvature functionals}\label{appendicesubriem}
   In this section, we apply the results of \Cref{appendiceriem} to establish variation formulas in the sub-Riemannian Heisenberg group $\hh^n$, for $n\geq 1$. Namely, we approximate it with a sequence of Riemannian manifolds $(\hh^n,\langle\cdot,\cdot\rangle_\eps)_{\eps>0}$, and we exploit \Cref{maintheoremappendixA} to deduce the analogous sub-Riemannian variation formulas. 
   
%In this section we collect some basic preliminaries on Heisenberg groups and on the geometry of its hypersurfaces.
\subsection{Heisenberg groups}\label{subsec_hg} The Heisenberg group $(\hn,\cdot)$ is $\mathbb R^{2n+1}$ endowed with a group law
that realizes it as stratified Lie group. Its Lie algebra is generated by the (global frame of) left-invariant vector fields 
\begin{equation*}
    Z_i=X_i=\frac{\partial}{\partial x_i}+y_i\frac{\partial}{\partial t},\qquad Z_{n+i}=Y_i=\frac{\partial}{\partial y_i}-x_i\frac{\partial}{\partial t},\qquad Z_{2n+1}=T=\frac{\partial}{\partial t},\qquad i=1,\ldots,n.
\end{equation*}
The only nontrivial commutation relations among $Z_1,\ldots,Z_{2n+1}$ are
\begin{equation*}
    [X_i,Y_i]=-2T,\qquad i=1,\ldots,n.
\end{equation*}
Accordingly, the \emph{horizontal distribution} $\hhh$ generated by $Z_1,\ldots,Z_{2n}$ is \emph{bracket-generating}. A vector field which is tangent to $\hhh$ at every point is called \emph{horizontal}. %Accordingly, the \emph{horizontal gradient} of a sufficiently regular function is the horizontal vector field defined by 
 %\begin{equation*}
 %    \nabla^\hhh f=\sum_{i=1}^{2n}Z_ifZ_i.
% \end{equation*}
The \emph{complex structure} $J:\Gamma(T\hh^n)\longrightarrow\Gamma(T\hh^n)$ is the unique $C^\infty(\hh^n)$-linear map which satisfies
\begin{equation*}
    J(X_i)=Y_i,\qquad J(Y_i)=-X_i,\qquad J(T)=0,\qquad i=1,\ldots,n.
\end{equation*}
 The triple $(\hh^n,\hhh,J)$ is a prototype of \emph{pseudohermitian manifold}  \cite[Appendix]{MR2165405}. 
The restriction to $\hhh$ of the unique Riemannian metric $\langle\cdot,\cdot\rangle$ making $Z_1,\ldots,Z_{2n+1}$ orthonormal endows $\hh^n$ with the sub-Riemannian structure $\left(\hh^n,\hhh,\langle\cdot,\cdot\rangle|_\hhh\right).$ 
%We denote the Riemannian gradient by 
%\begin{equation*}
%    \nabla f=\sum_{i=1}^{2n+1}Z_i f Z_i.
%\end{equation*}
 The \emph{pseudohermitian connection} $\nabla$ \cite{MR4193432} is the unique metric connection for which
\begin{equation}\label{pseudotorsion}
        \nabla_\A\B-\nabla_\B \A-[\A,\B]=2\langle J(\A),\B\rangle T,\qquad \A,\B\in\Gamma(T\hh^n).
\end{equation}
 The Riemannian volume induced by $\langle\cdot,\cdot\rangle$ is the Haar measure of the group, i.e. the standard Lebesgue measure.  Therefore, the Riemannian divergence induced by $\left\langle\cdot,\cdot\right\rangle$ is the Euclidean divergence.
\subsection{Hypersurfaces}\label{subsechypersurf}  Throughout this section,  $S\subseteq\hh^n$ is an embedded, closed, two-sided hypersurface of class $C^2$. A point $p\in S$ is called \emph{characteristic} if $\hhh_p=T_p S$. The set of characteristic points of $S$ is denoted by $S_0$. We will always assume that $S\setminus S_0$ is smooth. At non-characteristic points, the \emph{horizontal tangent space} $\hhh TS$ is the smooth, $(2n-1)$-dimensional distribution defined by 
\begin{equation*}
    \hhh T_p S=\hhh_p\cap T_p S,\qquad p\in S\setminus S_0.
\end{equation*}
Denote by $\n$ the Riemannian unit normal to $S$, and by $\n^\hhh$ its orthogonal projection onto $\hhh$. 
Then, the \emph{horizontal unit normal }  $$\vh=\frac{\n^\hhh}{|\n^\hhh|}$$ 
is well-defined on $S\setminus S_0$, and is the unique, up to sign, horizontal unit vector field orthogonal to $\hhh TS$. Notice that $p$ is a characteristic point if and only if $|\n^\hhh|=0$. Close to every non-characteristic point, it is always possible to extend $\vh$ to a full neighborhood in $\hh^n$ by setting 
   \begin{equation}\label{normcondist}
    \vh=\nabla^\hhh d,
\end{equation}
   where $d$ is the signed \emph{Carnot-Carathéodory distance} from $S$ (cf. \cite{simons,MR4193432}). %Indeed, locally near $p$, $d$ is smooth and satisfies the eikonal equation $|\nabla^\hhh d|=1$ . 
   Henceforth, $\vh$ is always extended as in \eqref{normcondist}.
   The \emph{fundamental function} $\alpha$ is defined on $S\setminus S_0$ as the unique smooth function such that 
   \begin{equation*}\label{firstdefinitionofthevecotrs}
       \s\coloneqq T-\alpha\vh\in\Gamma(TS).
   \end{equation*}
   It is known (cf. \cite{simons,MR4923606}) that $\alpha=Td$, and moreover
  \begin{equation}\label{propvh5}
    \nabla_{\vh}\vh=-2\alpha J(\vh).
\end{equation}
   Denote by $\hhh' TS$ the distribution defined by
\begin{equation*}
    \hhh'T_p S=\hhh T_pS\cap J\left(\hhh T_p S\right),\qquad p\in S\setminus S_0.
\end{equation*}
Then, $\hhh'TS$ is a $(2n-2)$-dimensional sub-bundle of $\hhh TS$, and the latter can be orthogonally decomposed as $
    \hhh TS=\hhh' TS\oplus\spann J(\vh). $
Notice that, in the first Heisenberg group, $\hhh' TS=\{0\}$. 
It is easy to check that 
\begin{equation*}
    \n=\frac{1}{\sqrt{1+\alpha^2}}\vh+\frac{\alpha}{\sqrt{1+\alpha^2}}T,\qquad \alpha=\frac{\left\langle \n,T\right\rangle}{|\n^\hhh|}.
\end{equation*}
%\begin{equation}
  %  TS=\hhh' TS\oplus\spann J(\vh)\oplus\spann \s,
%\end{equation}
%where 
%\begin{equation}\label{formadiesseepsilon}
 %   \s=-\frac{\alpha}{\sqrt{1+\eps^2\alpha^2}}\vh+\frac{1}{\sqrt{1+\eps^2\alpha^2}}\eps T.
%\end{equation}
%Moreover, set
%\begin{equation*}
%    \s\coloneqq\lim_{\eps\to 0}\frac{1}{\eps}\s^\eps=-\alpha\vh+T.
%\end{equation*}
%We recall that the \emph{pseudohermitian connection} of $\hh^n$, say $\nabla$, is the unique metric connection such that
%\begin{equation*}
%    \nabla_{Z_i}Z_j=0,\qquad i,j=1,\ldots,2n+1.
%\end{equation*}
%Accordingly, if $\A,\B,\C\in\Gamma(\hhh)$, \eqref{levi_civita} implies that
%\begin{equation}\label{relazionetraleconnessioni}
%    \left\langle\nabla^\eps_\A\B,\C\right\rangle_\eps=\left\langle\nabla_\A\B,\C\right\rangle.
%\end{equation}
The \emph{horizontal second fundamental form} $h^\hhh$ and the \emph{symmetric horizontal second fundamental form} (cf. \cite{simons,MR4193432}) $\tilde h^\hhh$ are defined on $S\setminus S_0$ respectively by
\begin{equation*}
    h^\hhh(\A,\B)\coloneqq\left\langle A^\hhh(\A),\B\right\rangle,\qquad \tilde h^\hhh(\A,\B)\coloneqq\frac{h^\hhh(\A,\B)+h^\hhh(\B,\A)}{2},\qquad\A,\B\in\Gamma(\hhh TS),
\end{equation*}
 where $ A^\hhh(\A)\coloneqq \nabla_\A\vh$ is the \emph{horizontal shape operator}. The forms $h^\hhh$ and $\tilde h^\hhh$ are related (cf. \cite{simons}) by
\begin{equation}\label{rapptrahetildacca2026}
    \tilde h^\hhh(\A,\B)=h^\hhh(\A,\B)+\alpha\left\langle J(\A),\B\right\rangle,\qquad\A,\B\in\Gamma(\hhh TS).
\end{equation}
%For any $\eps>0$, denote by $\sigma^\eps$ the Riemannian surface measure induced by $g_\eps$. It is easy to check that 
%\begin{equation}\label{areaelementeps}
 %   \sigma^\eps=\frac{\sqrt{1+\eps^2\alpha^2}}{\eps\sqrt{1+\alpha^2}}\sigma^1.
%\end{equation}
The \emph{horizontal mean curvature} is then defined on $S\setminus S_0$ by
\begin{equation*}
    H^\hhh=\trace h^\hhh=\trace \tilde h^\hhh.
\end{equation*}
%We say that $S$ is \emph{minimal} if $H^\hhh=0$ on $S\setminus S_0$, \emph{mean convex} if $H^\hhh\geq 0$ on $S\setminus S_0$, and \emph{strictly mean convex} if $H^\hhh>0$ on $S\setminus S_0$. 
Finally, the relevant sub-Riemannian surface measure $\sigma^\hhh$ is defined (cf. \cite{MR2354992, MR2262196}) by
\begin{equation}\label{areaelementorizz}
    \sigma^\hhh\coloneqq \frac{1}{\sqrt{1+\alpha^2}}\sigma,
\end{equation}
where $\sigma$ is the Riemannian surface measure induced by $\langle\cdot,\cdot\rangle$. Since $S$ is of class $C^2$, then $\sigma^\hhh(S_0)=\sigma(S_0)=0$ \cite{MR2021034}: accordingly, $S_0$ will be omitted in the forthcoming integrals.
The \emph{tangent pseudohermitian connection} $\nabla^S$ is the affine connection defined on $S$ by
\begin{equation*}
    \nabla^S _\A \B=\nabla_\A\B-\langle\nabla_\A\B,\vh\rangle\vh,\qquad \A,\B\in\Gamma(\hhh TS).
\end{equation*}
Let $\E_1,\ldots,\E_{2n-1}$ be any local orthonormal frame of $\hhh TS$. If $\varphi\in C^\infty(S\setminus S_0)$ and $\A\in\Gamma(\hhh TS)$ is supported in $S\setminus S_0$, the \emph{horizontal tangential gradient} of $\varphi$ and the  \emph{horizontal tangential divergence} of $\A$ are defined by 
\begin{equation*}
    \nabla ^{\hhh,S}\varphi=\sum_{i=1}^{2n-1}\left(\E_i\varphi\right)\E_i,\qquad \divv^{\hhh,S}\A\coloneqq\sum_{i=1}^{2n-1}\left\langle \nabla^S_{\E_i}\A,\E_i\right\rangle.
\end{equation*}
The \emph{horizontal tangential Laplacian} and the \emph{modified horizontal tangential Laplacian} are then defined by
\begin{equation*}
    \Delta^{\hhh,S} \varphi=\divv^{\hhh,S}\nabla^{\hhh,S}\varphi,\qquad \hat\Delta^{\hhh,S}\varphi\coloneqq \Delta^{\hhh,S}\varphi +2\alpha J(\vh)\varphi.
\end{equation*}
   Finally, we denote by $\jacobi^\hhh$ the \emph{horizontal Jacobi operator}
    \begin{equation*}\label{horizontaljacopdefmain}
    \jacobi^\hhh \varphi=-\hat\Delta^{\hhh,S}\varphi-\varphi\left(|\tilde h^\hhh|^2+4J(\vh)\alpha+(2n+2)\alpha^2\right),\qquad\varphi\in C^\infty(S\setminus S_0).
    \end{equation*}
    Unlike $\Delta^{\hhh,S}$  (cf. \cite{MR2354992}), both $\hat\Delta^{\hhh,S}$ and $\jacobi^\hhh$ are self-adjoint on $C^\infty_c(S\setminus S_0)$ (cf. \Cref{proposizioneibpsubriemsurf}).   
   \subsection{Riemannian approximation}
   For any $\eps>0$, denote by $\left\langle\cdot,\cdot\right\rangle_\eps$ the unique Riemannian metric on $\hh^n$ such that $(Z_1^\eps,\ldots,Z_{2n+1}^\eps)=(X_1,\ldots,X_n,Y_1,\ldots,Y_n,\eps T)$ is a global orthonormal frame. If $\A\in\Gamma(T\hh^n)$, denote by $A^{1,\eps},\ldots,A^{2n,\eps},A^{2n+1,\eps}$ its components with respect to the above frame. Clearly $A^j\coloneqq A^{j,\eps}$ is independent of $\eps$ for $j=1,\ldots,2n$. We may write $A^{2n+1}\coloneqq A^{2n+1,\eps}$ when there is no ambiguity.
   Since $\left\langle\A,\B\right\rangle_\eps$ does not depend on $\eps$ when $\A,\B\in\Gamma(\hhh)$, we set $\left\langle\A,\B\right\rangle\coloneqq \left\langle\A,\B\right\rangle_\eps.$    Denote by $\nabla^\eps$ and $\ \R^\eps $ the induced Levi-Civita connection and curvature tensor respectively. We recall the following relations (cf. \cite{MR5029815}):
   \begin{alignat}{2}
&\notag \nabla_{X_i}^\eps X_j=0, \qquad \qquad  \ \, \nabla^\eps_{Y_i}Y_j=0, \qquad \ \,\nabla^\eps_{T}T=0, \\
\label{levi_civita}&\nabla^\eps_{X_i}Y_j=-\delta_{i,j}T, \qquad \, \nabla^\eps_{X_i}\eps T=\frac{Y_i}{\eps}, \quad \ \ \nabla^\eps_{Y_i}\eps T=-\frac{X_i}{\eps}, \\
&\notag \nabla^\eps_{Y_i}X_j=\delta_{i,j}T, \qquad \ \ \,\,\nabla^\eps_{\eps T}X_i=\frac{Y_i}{\eps}, \  \ \quad \nabla^\eps_{\eps T}Y_i=-\frac{X_i}{\eps}.
\end{alignat}
%Denote by $J$ the linear map acting on $\Gamma(T\hh^n)$ as
%\begin{equation*}
%    J(X_i)=Y_i,\qquad J(Y_i)=-X_i,\qquad J(T)=0,\qquad i=1,\ldots,n.
%\end{equation*}
The Levi-Civita connection $\nabla^\eps$ relates to the pseudohermitian connection $\nabla$ as follows.
\begin{lemma}
    Let $\A,\B\in\Gamma(T\hh^n)$. Then
    \begin{equation}\label{levicivitavspseudohermitian}
        \nabla^\eps_\A\B=\nabla_\A\B+\frac{\langle\A,\eps T\rangle_\eps}{\eps}J(\B)+\frac{\langle\B,\eps T\rangle_\eps}{\eps}J(\A)+\langle \A,J(\B)\rangle T.
    \end{equation}
    In particular,
    \begin{equation}\label{connessioneerotazione}
    \nabla^\eps_\A \eps T=\frac{1}{\eps}J(\A).
\end{equation}
Moreover, if $\A,\B,\C\in\Gamma(\hhh)$, then
\begin{equation}\label{relazionetraleconnessioni}
    \left\langle\nabla^\eps_\A\B,\C\right\rangle_\eps=\left\langle\nabla_\A\B,\C\right\rangle.
\end{equation}
\end{lemma}
\begin{proof}
    By a direct computation,
    \begin{equation*}
        \begin{split}
            \nabla^\eps_\A\B&=\sum_{j,k=1}^{2n+1}A^{j,\eps}Z_j^\eps B^{k,\eps}Z_k^\eps+\sum_{j,k=1}^{2n+1}A^{j,\eps}B^{k,\eps}\nabla^\eps_{Z_j^\eps}Z_k^\eps\\
            \overset{\eqref{levi_civita}}&{=}\nabla_\A\B+\sum_{j=1}^nA^jB^{n+j}\nabla^\eps_{X_j}Y_j+\langle \B,\eps T\rangle_\eps\sum_{j=1}^nA^j\nabla^\eps_{X_j}\eps T+\sum_{j=1}^nA^{n+j}B^{j}\nabla^\eps_{Y_j}X_j\\
            &\quad+\langle \B,\eps T\rangle_\eps\sum_{j=1}^nA^{n+j}\nabla^\eps_{Y_j}\eps T+\langle \A,\eps T\rangle_\eps\sum_{j=1}^nB^{j}\nabla^\eps_{\eps T}X_j+\langle \A,\eps T\rangle_\eps\sum_{j=1}^nB^{n+j}\nabla^\eps_{\eps T}Y_j\\
            \overset{\eqref{levi_civita}}&{=}\nabla_\A\B+\frac{\langle\A,\eps T\rangle_\eps}{\eps}J(\B)+\frac{\langle\B,\eps T\rangle_\eps}{\eps}J(\A)+\langle \A,J(\B)\rangle T.
        \end{split}
    \end{equation*}
\end{proof}
\subsection{Ambient curvatures} First, we compute the ambient curvatures of $(\hh^n,\langle\cdot,\cdot \rangle_\eps)$, starting from $\R^\eps$.
\begin{proposition}\label{riemepsexplicitprop}
    Let $\eps>0$. Let $\A,\B,\C\in\Gamma(T\hh^n)$. Then
    \begin{equation*}
        \begin{split}
             \R^\eps (\A,\B)\C&=\frac{2}{\eps^2}\left\langle J(\A),\B\right\rangle _\eps J(\C)+\frac{1}{\eps^2}\left\langle J(\A),\C\right\rangle_\eps J(\B)-\frac{1}{\eps^2}\left\langle J(\B),\C\right\rangle _\eps J(\A)\\
            &\quad+\frac{A^{2n+1,\eps}}{\eps^2}\langle\B,\C\rangle _\eps\eps T-\frac{B^{2n+1,\eps}}{\eps^2}\langle\A,\C\rangle_\eps \eps T+\frac{B^{2n+1,\eps}C^{2n+1,\eps}}{\eps^2}\A-\frac{A^{2n+1,\eps}C^{2n+1,\eps}}{\eps^2}\B.
        \end{split}
    \end{equation*}
    In particular, if $\D\in\Gamma(T\hh^n)$,
    \begin{equation}\label{riemannepsilonquattrozero}
        \begin{split}
            & \R^\eps (\A,\B,\C,\D)=\frac{2}{\eps^2}\left\langle J(\A),\B\right\rangle _\eps \left\langle J(\C),\D\right\rangle_\eps+\frac{1}{\eps^2}\left\langle J(\A),\C\right\rangle_\eps \left\langle J(\B),\D\right\rangle_\eps-\frac{1}{\eps^2}\left\langle J(\B),\C\right\rangle _\eps \left\langle J(\A),\D\right\rangle_\eps\\
            &\quad+\frac{A^{2n+1,\eps}D^{2n+1,\eps}}{\eps^2}\langle\B,\C\rangle _\eps-\frac{B^{2n+1,\eps}D^{2n+1,\eps}}{\eps^2}\langle\A,\C\rangle_\eps +\frac{B^{2n+1,\eps}C^{2n+1,\eps}}{\eps^2}\left\langle \A,\D\right\rangle_\eps-\frac{A^{2n+1,\eps}C^{2n+1,\eps}}{\eps^2}\langle\B,\D\rangle_\eps.
        \end{split}
    \end{equation}
\end{proposition}
\begin{proof}
Fix $i,j,k=1,\ldots,n$. Then, by \eqref{levi_civita},
\begin{align*}
     \R^\eps (X_i,X_j)X_k&=0,\\
     \R^\eps (X_i,X_j)Y_k&=\nabla^{\eps}_{X_i}\nabla^\eps_{X_j}Y_k-\nabla^{\eps}_{X_j}\nabla^\eps_{X_i}Y_k-\nabla^\eps_{[X_i,X_j]}Y_k=-\frac{\delta_{jk}}{\eps}\nabla^{\eps}_{X_i}\eps T+\frac{\delta_{ik}}{\eps}\nabla^{\eps}_{X_j}\eps T=\frac{1}{\eps^2}\left(\delta_{ik}Y_j-\delta_{jk}Y_i\right),\\
     \R^\eps (X_i,X_j)\eps T&=\nabla^{\eps}_{X_i}\nabla^\eps_{X_j}\eps T-\nabla^{\eps}_{X_j}\nabla^\eps_{X_i}\eps T-\nabla^\eps_{[X_i,X_j]}\eps T=\frac{1}{\eps}\nabla^{\eps}_{X_i}Y_j-\frac{1}{\eps}\nabla^{\eps}_{X_j}Y_i=0.
\end{align*}
Moreover, 
\begin{align*}
     \R^\eps (X_i,Y_j)X_k&=\nabla^{\eps}_{X_i}\nabla^\eps_{Y_j}X_k-\nabla^{\eps}_{Y_j}\nabla^\eps_{X_i}X_k-\nabla^\eps_{[X_i,Y_j]}X_k=\frac{\delta_{jk}}{\eps}\nabla^{\eps}_{X_i}\eps T+\frac{2\delta_{ij}}{\eps}\nabla^\eps_{\eps T}X_k=\frac{1}{\eps^2}\left(\delta_{jk}Y_i+2\delta_{ij}Y_k\right),\\
     \R^\eps (X_i,Y_j)Y_k&=\nabla^{\eps}_{X_i}\nabla^\eps_{Y_j}Y_k-\nabla^{\eps}_{Y_j}\nabla^\eps_{X_i}Y_k-\nabla^\eps_{[X_i,Y_j]}Y_k=\frac{\delta_{ik}}{\eps}\nabla^\eps_{Y_j}\eps T+\frac{2\delta_{ij}}{\eps}\nabla^\eps_{\eps T}Y_k=-\frac{1}{\eps^2}\left(\delta_{ik}X_j+2\delta_{ij}X_k\right),\\
     \R^\eps (X_i,Y_j)\eps T&=\nabla^{\eps}_{X_i}\nabla^\eps_{Y_j}\eps T-\nabla^{\eps}_{Y_j}\nabla^\eps_{X_i}\eps T-\nabla^\eps_{[X_i,Y_j]}\eps T=0.
\end{align*}
In addition,
\begin{align*}
     \R^\eps (X_i,\eps T)X_k&=\nabla^{\eps}_{X_i}\nabla^\eps_{\eps T}X_k-\nabla^{\eps}_{\eps T}\nabla^\eps_{X_i}X_k-\nabla^\eps_{[X_i,\eps T]}X_k=\frac{1}{\eps}\nabla^{\eps}_{X_i}Y_k=-\frac{\delta_{ik}}{\eps^2}\eps T,\\
      \R^\eps (X_i,\eps T)Y_k&=\nabla^{\eps}_{X_i}\nabla^\eps_{\eps T}Y_k-\nabla^{\eps}_{\eps T}\nabla^\eps_{X_i}Y_k-\nabla^\eps_{[X_i,\eps T]}Y_k=0,\\
       \R^\eps (X_i,\eps T)\eps T&=\nabla^{\eps}_{X_i}\nabla^\eps_{\eps T}\eps T-\nabla^{\eps}_{\eps T}\nabla^\eps_{X_i}\eps T-\nabla^\eps_{[X_i,\eps T]}\eps T=-\frac{1}{\eps}\nabla^{\eps}_{\eps T}Y_i=\frac{1}{\eps^2}X_i.
\end{align*}
Furthermore,
\begin{align*}
     \R^\eps (Y_i,Y_j)X_k&=\nabla^{\eps}_{Y_i}\nabla^\eps_{Y_j}X_k-\nabla^{\eps}_{Y_j}\nabla^\eps_{Y_i}X_k-\nabla^\eps_{[Y_i,Y_j]}X_k=\frac{\delta_{jk}}{\eps}\nabla^{\eps}_{Y_i}\eps T-\frac{\delta_{ik}}{\eps}\nabla^{\eps}_{Y_j}\eps T=\frac{1}{\eps^2}\left(\delta_{ik}X_j-\delta_{jk}X_i\right),\\
     \R^\eps (Y_i,Y_j)Y_k&=0,\\
     \R^\eps (Y_i,Y_j)\eps T&=\nabla^{\eps}_{Y_i}\nabla^\eps_{Y_j}\eps T-\nabla^{\eps}_{Y_j}\nabla^\eps_{Y_i}\eps T-\nabla^\eps_{[Y_i,Y_j]}\eps T=-\frac{1}{\eps}\nabla^{\eps}_{Y_i}X_j+\frac{1}{\eps}\nabla^{\eps}_{Y_j}X_i=0.
\end{align*}
Finally,
\begin{align*}
      \R^\eps (Y_i,\eps T)X_k&=\nabla^{\eps}_{Y_i}\nabla^\eps_{\eps T}X_k-\nabla^{\eps}_{\eps T}\nabla^\eps_{Y_i}X_k-\nabla^\eps_{[Y_i,\eps T]}X_k=0,\\
        \R^\eps (Y_i,\eps T)Y_k&=\nabla^{\eps}_{Y_i}\nabla^\eps_{\eps T}Y_k-\nabla^{\eps}_{\eps T}\nabla^\eps_{Y_i}Y_k-\nabla^\eps_{[Y_i,\eps T]}Y_k=-\frac{1}{\eps}\nabla^{\eps}_{Y_i}X_k=-\frac{\delta_{ik}}{\eps^2}\eps T,\\
         \R^\eps (Y_i,\eps T)\eps T&=\nabla^{\eps}_{Y_i}\nabla^\eps_{\eps T}\eps T-\nabla^{\eps}_{\eps T}\nabla^\eps_{Y_i}\eps T-\nabla^\eps_{[Y_i,\eps T]}\eps T=\frac{1}{\eps}\nabla^\eps_{\eps T}X_i=\frac{1}{\eps^2}Y_i.
\end{align*}
    Therefore 
    \begin{equation*}
        \begin{split}
            & \R^\eps (\A,\B)\C=\sum_{i,j,k=1}^{2n+1}A^iB^jC^k \R^\eps (Z_i,Z_j)Z_k\\
            &=\sum_{i=1}^n\sum_{j,k=1}^{2n+1}A^iB^jC^k \R^\eps (X_i,Z_j)Z_k+\sum_{i=1}^n\sum_{j,k=1}^{2n+1}A^{n+i}B^jC^k \R^\eps (Y_i,Z_j)Z_k+A^{2n+1}\sum_{j,k=1}^{2n+1}B^jC^k \R^\eps (\eps T,Z_j)Z_k\\
            &=\sum_{i,j=1}^n\sum_{k=1}^{2n+1}A^iB^jC^k \R^\eps (X_i,X_j)Z_k+\sum_{i,j=1}^n\sum_{k=1}^{2n+1}A^iB^{n+j}C^k \R^\eps (X_i,Y_j)Z_k+B^{2n+1}\sum_{i=1}^n\sum_{k=1}^{2n+1}A^iC^k \R^\eps (X_i,\eps T)Z_k\\
            &\quad+\sum_{i,j=1}^n\sum_{k=1}^{2n+1}A^{n+i}B^jC^k \R^\eps (Y_i,X_j)Z_k+\sum_{i,j=1}^n\sum_{k=1}^{2n+1}A^{n+i}B^{n+j}C^k \R^\eps (Y_i,Y_j)Z_k\\
            &\quad+B^{2n+1}\sum_{i=1}^n\sum_{k=1}^{2n+1}A^{n+i}C^k \R^\eps (Y_i,\eps T)Z_k+A^{2n+1}\sum_{j=1}^n\sum_{k=1}^{2n+1}B^jC^k \R^\eps (\eps T,X_j)Z_k\\
            &\quad +A^{2n+1}\sum_{j=1}^n\sum_{k=1}^{2n+1}B^{n+j}C^k \R^\eps (\eps T,Y_j)Z_k.           
        \end{split}
    \end{equation*}
    By the symmetry of $\ \R^\eps ,$ 
    \begin{equation*}
        \begin{split}
             \R^\eps (\A,\B)\C&=\underbrace{\sum_{i,j=1}^n\sum_{k=1}^{2n+1}A^iB^jC^k \R^\eps (X_i,X_j)Z_k}_\mathrm{I}+\underbrace{\sum_{i,j=1}^n\sum_{k=1}^{2n+1}\left(A^iB^{n+j}-A^{n+j}B^i\right)C^k \R^\eps (X_i,Y_j)Z_k}_\mathrm{II}\\
            &\quad+\underbrace{\sum_{i=1}^n\sum_{k=1}^{2n+1}\left(A^iB^{2n+1}-A^{2n+1}B^i\right)C^k \R^\eps (X_i,\eps T)Z_k}_\mathrm{III}+\underbrace{\sum_{i,j=1}^n\sum_{k=1}^{2n+1}A^{n+i}B^{n+j}C^k \R^\eps (Y_i,Y_j)Z_k}_\mathrm{IV}\\
            &\quad+\underbrace{\sum_{i=1}^n\sum_{k=1}^{2n+1}\left(A^{n+i}B^{2n+1}-A^{2n+1}B^{n+i}\right)C^k \R^\eps (Y_i,\eps T)Z_k}_\mathrm{V}.
        \end{split}
    \end{equation*}
    First,
    \begin{equation*}
        \begin{split}
            \mathrm{I}=\frac{1}{\eps^2}\sum_{i,j=1}^nA^iB^jC^{n+i}Y_j-\frac{1}{\eps^2}\sum_{i,j=1}^nA^iB^jC^{n+j}Y_i=\frac{1}{\eps^2}\sum_{i,j=1}^nA^jB^iC^{n+j}Y_i-\frac{1}{\eps^2}\sum_{i,j=1}^nA^iB^jC^{n+j}Y_i.
        \end{split}
    \end{equation*}
    Moreover,
    \begin{equation*}
        \begin{split}
            \mathrm{II}&=\sum_{i,j,k=1}^n\left(A^iB^{n+j}-A^{n+j}B^i\right)C^k \R^\eps (X_i,Y_j)X_k+\sum_{i,j,k=1}^n\left(A^iB^{n+j}-A^{n+j}B^i\right)C^{n+k} \R^\eps (X_i,Y_j)Y_k\\
            &=\frac{1}{\eps^2}\sum_{i,j=1}^n\left(A^iB^{n+j}-A^{n+j}B^i\right)C^jY_i+\frac{2}{\eps^2}\sum_{i,k=1}^n\left(A^iB^{n+i}-A^{n+i}B^i\right)C^kY_k\\
            &\quad-\frac{1}{\eps^2}\sum_{i,j=1}^n\left(A^iB^{n+j}-A^{n+j}B^i\right)C^{n+i}X_j-\frac{2}{\eps^2}\sum_{i,k=1}^n\left(A^iB^{n+i}-A^{n+i}B^i\right)C^{n+k}X_i\\
            &=\frac{1}{\eps^2}\sum_{i,j=1}^n\left(A^iB^{n+j}-A^{n+j}B^i\right)C^jY_i-\frac{1}{\eps^2}\sum_{i,j=1}^n\left(A^jB^{n+i}-A^{n+i}B^j\right)C^{n+j}X_i+\frac{2}{\eps^2}\left\langle J(\A),\B\right\rangle J(\C).
        \end{split}
    \end{equation*}
    In addition,
    \begin{equation*}
        \begin{split}
            \mathrm{III}&=\sum_{i,k=1}^n\left(A^iB^{2n+1}-A^{2n+1}B^i\right)C^k \R^\eps (X_i,\eps T)X_k+C^{2n+1}\sum_{i=1}^n\left(A^iB^{2n+1}-A^{2n+1}B^i\right) \R^\eps (X_i,\eps T)\eps T\\
            &=-\frac{1}{\eps^2}\sum_{i=1}^n\left(A^iB^{2n+1}-A^{2n+1}B^i\right)C^i\eps T+\frac{C^{2n+1}}{\eps^2}\sum_{i=1}^n\left(A^iB^{2n+1}-A^{2n+1}B^i\right)X_i\\
            &=-\frac{B^{2n+1}}{\eps^2}\sum_{j=1}^nA^jC^j\eps T+\frac{A^{2n+1}}{\eps^2}\sum_{j=1}^nB^jC^j\eps T+\frac{B^{2n+1}C^{2n+1}}{\eps^2}\sum_{i=1}^nA^iX_i-\frac{A^{2n+1}C^{2n+1}}{\eps^2}\sum_{i=1}^nB^iX_i.
        \end{split}
    \end{equation*}
    Furthermore,
    \begin{equation*}
        \mathrm{IV}=\frac{1}{\eps^2}\sum_{i,j=1}^nA^{n+i}B^{n+j}C^iX_j-\frac{1}{\eps^2}\sum_{i,j=1}^nA^{n+i}B^{n+j}C^jX_i=\frac{1}{\eps^2}\sum_{i,j=1}^nA^{n+j}B^{n+i}C^jX_i-\frac{1}{\eps^2}\sum_{i,j=1}^nA^{n+i}B^{n+j}C^jX_i,
    \end{equation*}
    Finally,
    \begin{equation*}
        \begin{split}
            \mathrm{V}&=\sum_{i,k=1}^n\left(A^{n+i}B^{2n+1}-A^{2n+1}B^{n+i}\right)C^{n+k} \R^\eps (Y_i,\eps T)Y_k\\
            &\quad+C^{2n+1}\sum_{i=1}^n\left(A^{n+i}B^{2n+1}-A^{2n+1}B^{n+i}\right) \R^\eps (Y_i,\eps T)\eps T\\
            &=-\frac{1}{\eps^2}\sum_{i=1}^n\left(A^{n+i}B^{2n+1}-A^{2n+1}B^{n+i}\right)C^{n+i}\eps T+\frac{C^{2n+1}}{\eps^2}\sum_{i=1}^n\left(A^{n+i}B^{2n+1}-A^{2n+1}B^{n+i}\right)Y_i\\
            &=-\frac{B^{2n+1}}{\eps^2}\sum_{i=1}^nA^{n+i}C^{n+i}\eps T+\frac{A^{2n+1}}{\eps^2}\sum_{i=1}^nB^{n+i}C^{n+i}\eps T\\
            &\quad+\frac{B^{2n+1}C^{2n+1}}{\eps^2}\sum_{i=1}^nA^{n+i}Y_i-\frac{A^{2n+1}C^{2n+1}}{\eps^2}\sum_{i=1}^nB^{n+i}Y_i.
        \end{split}
    \end{equation*}
    By the above computations,
    \begin{equation*}
\begin{split}
    \mathrm{I}+\mathrm{II}+\mathrm{IV}&=\frac{2}{\eps^2}\left\langle J(\A),\B\right\rangle J(\C)+\frac{1}{\eps^2}\left\langle J(\A),\C\right\rangle \sum_{i=1}^nB^i Y_i-\frac{1}{\eps^2}\left\langle J(\A),\C\right\rangle \sum_{i=1}^nB^{n+i} X_i\\
    &\quad+\frac{1}{\eps^2}\left\langle J(\B),\C\right\rangle \sum_{i=1}^nA^{n+i} X_i-\frac{1}{\eps^2}\left\langle J(\B),\C\right\rangle \sum_{i=1}^nA^{i} Y_i\\
    &=\frac{2}{\eps^2}\left\langle J(\A),\B\right\rangle J(\C)+\frac{1}{\eps^2}\left\langle J(\A),\C\right\rangle J(\B)-\frac{1}{\eps^2}\left\langle J(\B),\C\right\rangle J(\A)
\end{split}        
    \end{equation*}
    and 
    \begin{equation*}
        \begin{split}
            \mathrm{III}+\mathrm{V}&=\frac{A^{2n+1}}{\eps^2}\langle\B,\C\rangle \eps T-\frac{B^{2n+1}}{\eps^2}\langle\A,\C\rangle \eps T+\frac{B^{2n+1}C^{2n+1}}{\eps^2}\A-\frac{A^{2n+1}C^{2n+1}}{\eps^2}\B.
        \end{split}
    \end{equation*}
    The thesis follows combining the above computations.
\end{proof}
The explicit expression of the Ricci curvature $\ric^\eps$ is a simple consequence of \Cref{riemepsexplicitprop}.
\begin{proposition}
    Let $\eps>0$. Let $\A,\B,\C\in\Gamma(T\hh^n)$. Then
    \begin{equation}\label{ricciepsilon}
        \begin{split}
            \ric^\eps(\B,\C)&=-\frac{2}{\eps^2}\left\langle \B,\C\right\rangle_\eps +\frac{(2n+2)}{\eps^2}\left\langle \B,\eps T\right\rangle_\eps\left\langle \C,\eps T\right\rangle_\eps
        \end{split}
    \end{equation}
    and 
    \begin{equation*}
        \left(\nabla^\eps_\A\ric^\eps\right)(\B,\C)=\frac{2n+2}{\eps^3}\left\langle J(\A),\B\right\rangle_\eps\left\langle \C,\eps T\right\rangle_\eps+\frac{2n+2}{\eps^3}\left\langle J(\A),\C\right\rangle_\eps\left\langle \B,\eps T\right\rangle_\eps
    \end{equation*}
    In particular, 
    \begin{equation}\label{nabaepsriccizerosudiagonale}
        \left(\nabla^\eps_\A\ric^\eps\right)(\A,\A)=0.
    \end{equation}
\end{proposition}
\begin{proof}
    By definition,
    \begin{equation*}
        \begin{split}
            \eps^2\ric(\B,\C)&
            %=\sum_{i=1}^{2n+1}\eps^2 \R^\eps (Z_i,\B,\C,Z_i)\\
    =\sum_{i=1}^{2n}\Big(2\left\langle J(Z_i),\B\right\rangle _\eps \left\langle J(\C),Z_i\right\rangle_\eps+\left\langle J(Z_i),\C\right\rangle_\eps \left\langle J(\B),Z_i\right\rangle_\eps-\left\langle J(\B),\C\right\rangle _\eps \left\langle J(Z_i),Z_i\right\rangle_\eps\Big)\\
            &\quad+\langle\B,\C\rangle _\eps-B^{2n+1,\eps}C^{2n+1,\eps}+(2n+1)B^{2n+1,\eps}C^{2n+1,\eps}-B^{2n+1,\eps}C^{2n+1,\eps}.\\
            &=-3\left\langle J(\B),J(\C)\right\rangle +\langle\B,\C\rangle _\eps+(2n-1)B^{2n+1,\eps}C^{2n+1,\eps}\\
            &=-2\left\langle \B,\C\right\rangle +(2n+2)B^{2n+1,\eps}C^{2n+1,\eps}.
        \end{split}
    \end{equation*}
    Therefore, since $\nabla^\eps\left\langle\cdot,\cdot\right\rangle_\eps\equiv 0$,
    \begin{equation*}
        \begin{split}
            \eps^2\left(\nabla_\A\ric^\eps\right)(\B,\C)&=(2n+2)\big(\A\left(\left\langle \B,\eps T\right\rangle_\eps\left\langle \C,\eps T\right\rangle_\eps\right)-\left\langle \nabla^\eps_\A\B,\eps T\right\rangle_\eps\left\langle \C,\eps T\right\rangle_\eps-\left\langle \B,\eps T\right\rangle_\eps\left\langle \nabla^\eps_\A\C,\eps T\right\rangle_\eps\big)\\
            &=(2n+2)\big(\left\langle \B,\nabla^\eps_\A\eps T\right\rangle_\eps\left\langle \C,\eps T\right\rangle_\eps+\left\langle \B,\eps T\right\rangle_\eps\left\langle \C,\nabla^\eps_\A\eps T\right\rangle_\eps\big)\\
            \overset{\eqref{connessioneerotazione}}&{=}\frac{2n+2}{\eps}\left\langle J(\A),\B\right\rangle_\eps\left\langle \C,\eps T\right\rangle_\eps+\frac{2n+2}{\eps}\left\langle J(\A),\C\right\rangle_\eps\left\langle \B,\eps T\right\rangle_\eps.
        \end{split}
    \end{equation*}
    Finally, \eqref{nabaepsriccizerosudiagonale} is straightforward.
       \end{proof}
   \subsection{Extrinsic curvatures}\label{sec_extrinsiccurvaturesriemapp} Fix $\eps>0$. We adapt some aspects of \Cref{subsechypersurf} to the Riemannian structure $(\hh^n,\langle\cdot,\cdot\rangle_\eps).$
We assign an upper index $\eps$ to the geometric quantities associated with $S$ and induced by $\left\langle\cdot,\cdot\right\rangle_\eps$. Accordingly,
\begin{equation}\label{normaleriemannianaepsilon}
    \n^\eps=\frac{1}{\sqrt{1+\eps^2\alpha^2}}\vh+\frac{\eps\alpha}{\sqrt{1+\eps^2\alpha^2}}\eps T.
\end{equation}
Therefore, $TS$ admits the $\left\langle\cdot,\cdot\right\rangle_\eps$-orthonormal  decomposition $TS=\hhh' TS\oplus\spann J(\vh)\oplus\spann \s^\eps,$ 
where 
\begin{equation}\label{formadiesseepsilon}
    \s^\eps=-\frac{\eps\alpha}{\sqrt{1+\eps^2\alpha^2}}\vh+\frac{1}{\sqrt{1+\eps^2\alpha^2}}\eps T.
\end{equation}
%Moreover, by definition (cf. \eqref{firstdefinitionofthevecotrs}),
%\begin{equation*}
%    \s\coloneqq\lim_{\eps\to 0}\frac{1}{\eps}\s^\eps.
%\end{equation*}
If $\sigma^\eps$ is the Riemannian surface measure induced by $\langle\cdot,\cdot\rangle_\eps$, it is easy to check that 
\begin{equation}\label{areaelementeps}
    \sigma^\eps=\frac{\sqrt{1+\eps^2\alpha^2}}{\eps\sqrt{1+\alpha^2}}\sigma^1.
\end{equation}
In the rest of this section, we express some Riemannian extrinsic quantities with respect to the relevant sub-Riemannian geometric objects. We begin with the second fundamental form.
\begin{lemma}\label{lemmasecondaformafondriem}
    Let $\eps>0$. Let $\A,\B\in\Gamma(\hhh T S)$. Let $\C\in\Gamma(\hhh'TS)$. Then
    \begin{equation}\label{formaepsenonepstuttoorizz}
        \begin{alignedat}{2}                  h^\eps(\A,\B)&=\frac{1}{\sqrt{1+\eps^2\alpha^2}}\tilde h^\hhh(\A,\B),&\qquad
        h^\eps(\C,\s^\eps)&=\frac{\eps\C\alpha}{1+\eps^2\alpha^2},\\
        h^\eps(J(\vh),\s^\eps)&=\frac{1}{\eps}+\frac{\eps J(\vh)\alpha}{1+\eps^2\alpha^2},&\qquad h^\eps(\s^\eps,\s^\eps)&=\frac{\eps^2\s\alpha}{\left(1+\eps^2\alpha^2\right)^{\frac{3}{2}}}.
        \end{alignedat}
    \end{equation}
\end{lemma}
\begin{proof}
    First, since
    \begin{equation}\label{utilenellaproofmadatogliere}
        \left\langle \vh,\B\right\rangle=\left\langle\eps T,\B\right\rangle_\eps=0,
    \end{equation}
    then
    \begin{equation*}
        \begin{split}
            h^\eps(\A,\B)%&=\left\langle\nabla^\eps_\A\n^\eps,\B\right\rangle_\eps\\
            \overset{\eqref{normaleriemannianaepsilon}}&{=}\left\langle\nabla^\eps_\A\left(\frac{1}{\sqrt{1+\eps^2\alpha^2}}\vh\right),\B\right\rangle_\eps+\left\langle\nabla^\eps_\A\left(\frac{\eps\alpha}{\sqrt{1+\eps^2\alpha^2}}\eps T\right),\B\right\rangle_\eps\\
            \overset{\eqref{utilenellaproofmadatogliere}}&{=}\frac{1}{\sqrt{1+\eps^2\alpha^2}}\Big(\left\langle\nabla^\eps_\A\vh,\B\right\rangle_\eps+\eps\alpha\left\langle\nabla^\eps_\A\eps T,\B\right\rangle_\eps\Big)\\
            \overset{\eqref{connessioneerotazione},\eqref{relazionetraleconnessioni}}&{=}\frac{1}{\sqrt{1+\eps^2\alpha^2}}\Big(h^\hhh(\A,\B)+\alpha\left\langle J(\A),\B\right\rangle\Big)\\
            \overset{\eqref{rapptrahetildacca2026}}&{=}\frac{1}{\sqrt{1+\eps^2\alpha^2}}\tilde h^\hhh(\A,\B).
        \end{split}
    \end{equation*}
    Let $\D\in\Gamma(TS)$. Then
    \begin{equation*}
        \begin{split}
            h^\eps&(\D,\s^\eps)=\left\langle\nabla^\eps_\D\n^\eps,\s^\eps\right\rangle_\eps\\
            \overset{\eqref{normaleriemannianaepsilon}}&{=}\frac{1}{\sqrt{1+\eps^2\alpha^2}}\left\langle\nabla^\eps_\D\left(\frac{1}{\sqrt{1+\eps^2\alpha^2}}\vh\right),-\eps\alpha \vh+\eps T\right\rangle_\eps+\frac{1}{\sqrt{1+\eps^2\alpha^2}}\left\langle\nabla^\eps_\D\left(\frac{\eps\alpha}{\sqrt{1+\eps^2\alpha^2}}\eps T\right),-\eps\alpha \vh+\eps T\right\rangle_\eps\\
            &=-\frac{\eps\alpha}{\sqrt{1+\eps^2\alpha^2}}\D\left(\frac{1}{\sqrt{1+\eps^2\alpha^2}}\right)+\frac{1}{1+\eps^2\alpha^2}\Big\langle\nabla^\eps_\D\vh,-\eps\alpha \vh+\eps T\Big\rangle_\eps\\
            &+\frac{\eps\alpha}{\sqrt{1+\eps^2\alpha^2}}\D\left(\frac{1}{\sqrt{1+\eps^2\alpha^2}}\right)+\frac{\eps\D\alpha}{1+\eps^2\alpha^2}+\frac{\eps\alpha}{1+\eps^2\alpha^2}\Big\langle\nabla^\eps_\D\eps T,-\eps\alpha \vh+\eps T\Big\rangle_\eps\\
            &=\frac{1}{1+\eps^2\alpha^2}\Big\langle\nabla^\eps_\D\vh,\eps T\Big\rangle_\eps-\frac{\eps^2\alpha^2}{1+\eps^2\alpha^2}\Big\langle\nabla^\eps_\D\eps T, \vh\Big\rangle_\eps+\frac{\eps\D\alpha}{1+\eps^2\alpha^2}\\
            \overset{\eqref{connessioneerotazione}}&{=}-\frac{1}{\eps(1+\eps^2\alpha^2)}\left\langle J(\D),\vh\right\rangle_\eps-\frac{\eps^2\alpha^2}{\eps(1+\eps^2\alpha^2)}\left\langle J(\D), \vh\right\rangle_\eps+\frac{\eps\D\alpha}{1+\eps^2\alpha^2}\\
            &=\frac{1}{\eps}\left\langle \D,J(\vh)\right\rangle+\frac{\eps\D\alpha}{1+\eps^2\alpha^2}.
        \end{split}
    \end{equation*}
\end{proof}
A first trivial consequence of \Cref{lemmasecondaformafondriem} is the behavior of the mean curvature.
\begin{corollary}
    Let $p\in S\setminus S_0$. Then
    \begin{equation}\label{curvepsecurvh}
    \begin{split}
         H^\eps&=\frac{1}{\sqrt{1+\eps^2\alpha^2}}H^\hhh+\frac{\eps^2\s\alpha}{\left(1+\eps^2\alpha^2\right)^{\frac{3}{2}}}.
    %    \A H^\eps&=\frac{\A H^\hhh}{\sqrt{1+\eps^2\alpha^2}}+\frac{\eps^2\A\s\alpha}{\left(1+\eps^2\alpha^2\right)^{\frac{3}{2}}}-\frac{\eps^2H^\hhh\alpha\A\alpha}{\left(1+\eps^2\alpha^2\right)^\frac{3}{2}}-\frac{3\eps^4\alpha\s\alpha\A\alpha}{\left(1+\eps^2\alpha^2\right)^\frac{5}{2}}.
    \end{split}
    \end{equation}
    In particular,
    \begin{equation}\label{convhepstoh}
        \A_1\cdots\A_k H^\eps\xrightarrow[\eps\to 0]{}\A_1\cdots\A_k H^\hhh\text{ locally uniformly on $S\setminus S_0$ for any $k\in\mathbb N$, $\A_1,\ldots,\A_k\in\Gamma(TS)$.}
    \end{equation}
   \end{corollary}
   In the next result, the approximation occurs both in the second fundamental form and in its entries.
\begin{lemma}\label{lemmasecofondformcondueparametri}
    Let $p\in S\setminus S_0$. Let $f,g\in C^\infty(S\setminus S_0)$. Then 
    \begin{equation}\label{hepsgradepsprimadilimieteffbgyhygtfr}
        \begin{split}
            h^\eps\left(\nabla^{\eps,S}f,\nabla^{\eps,S}g\right)&=\frac{\tilde h^\hhh\left(\nabla^{\hhh,S}f,\nabla^{\hhh,S}g\right)}{\sqrt{1+\eps^2\alpha^2}}+\frac{J(\vh) f\s g+J(\vh)g\s f}{\sqrt{1+\eps^2\alpha^2}}\\
            &\quad+\frac{\eps^2\left\langle \nabla^{\hhh,S}\alpha,\s g\nabla^{\hhh,S}f+\s f\nabla^{\hhh,S}g\right\rangle}{\left(1+\eps^2\alpha^2\right)^\frac{3}{2}}+\frac{\eps^4\s\alpha\s f\s g}{\left(1+\eps^2\alpha^2\right)^\frac{5}{2}}.
        \end{split}
    \end{equation}
    In particular, if $(f^\eps)_\eps,\,(g^\eps)_\eps$ converge smoothly to $f$ and $g$ as in \eqref{convhepstoh}, then    \begin{equation}\label{convhepsgradeps}
        \begin{split}
        h^\eps\left(\nabla^{\eps,S}f^\eps,\nabla^{\eps,S}g^\eps\right)&\xrightarrow[\eps\to 0]{}\tilde h^\hhh\left(\nabla^{\hhh,S}f,\nabla^{\hhh,S}g\right)+ J(\vh)f\,\s g+J(\vh)g\,\s f \text{ locally uniformly on $S\setminus S_0$.}
        \end{split}
    \end{equation}
    
\end{lemma}
\begin{proof}
       Fix $p\in S\setminus S_0$. Let $\E_1,\ldots,\E_{n-1},\E_{n+1},\ldots,\E_{2n-1}$ be a local orthonormal frame of $\hhh'TS$. Set $\E_n=J(\vh)$ and $\E_{2n}=\s^\eps$. In this way, $\E_1,\ldots,\E_{2n}$ is a local orthonormal frame of $TS$. Then
       \begin{equation*}
           \begin{split}
               h^\eps&\left(\nabla^{\eps,S}f,\nabla^{\eps,S}g\right)=\sum_{i,j=1}^{2n}h^\eps(\E_i,\E_j)\E_i f\E_i g\\
               &=\sum_{i,j=1}^{2n-1}h^\eps(\E_i,\E_j)\E_i f\E_i g+\sum_{i=1}^{2n-1}h^\eps(\E_i,\s^\eps)\left(\E_i f\s^\eps g+\E_i g\s^\eps f\right)+h^\eps(\s^\eps,\s^\eps)\s^\eps f\s^\eps g\\
               \overset{\eqref{formaepsenonepstuttoorizz}}&{=}\frac{\tilde h^\hhh\left(\nabla^{\hhh,S}f,\nabla^{\hhh,S}g\right)}{\sqrt{1+\eps^2\alpha^2}}+\frac{\eps^2}{\left(1+\eps^2\alpha^2\right)^\frac{3}{2}}\sum_{\substack{i=1\\ i\neq n}}^{2n-1}\E_i\alpha\left(\E_i f\s g+\E_i g\s f\right)\\
               &\quad+\frac{1}{\sqrt{1+\eps^2\alpha^2}}\left(1+\frac{\eps^2J(\vh)\alpha}{1+\eps^2\alpha^2}\right)\left(J(\vh) f\s g+J(\vh)g\s f\right)+\frac{\eps^4\s\alpha\s f\s g}{\left(1+\eps^2\alpha^2\right)^\frac{5}{2}}\\
               &=\frac{\tilde h^\hhh\left(\nabla^{\hhh,S}f,\nabla^{\hhh,S}g\right)}{\sqrt{1+\eps^2\alpha^2}}+\frac{J(\vh) f\s g+J(\vh)g\s f}{\sqrt{1+\eps^2\alpha^2}}+\frac{\eps^2\left\langle \nabla^{\hhh,S}\alpha,\s g\nabla^{\hhh,S}f+\s f\nabla^{\hhh,S}g\right\rangle}{\left(1+\eps^2\alpha^2\right)^\frac{3}{2}}+\frac{\eps^4\s\alpha\s f\s g}{\left(1+\eps^2\alpha^2\right)^\frac{5}{2}}.
           \end{split}
       \end{equation*}
       In particular, \eqref{convhepsgradeps} follows by \eqref{hepsgradepsprimadilimieteffbgyhygtfr}.
\end{proof}
The next convergence result has been achieved, through a different approach, in \cite{MR5029815}. 
\begin{lemma}
    Let $p\in S\setminus S_0$. Then
    \begin{equation*}
        \left|h^\eps\right|^2+\ric^\eps(\n^\eps,\n^\eps)=\frac{|\tilde h^\hhh|^2+4 J(\vh)\alpha+(2n+2)\alpha^2}{1+\eps^2\alpha^2}+\frac{2\eps^2\left|\nabla^{\hhh,S}\alpha\right|^2}{\left(1+\eps^2\alpha^2\right)^2}+\frac{\eps^4\left(\s\alpha\right)^2}{\left(1+\eps^2\alpha^2\right)^3}
    \end{equation*}
    In particular,
    \begin{equation}\label{convergenzadelpotenzialejacobi}
        \left|h^\eps\right|^2+\ric^\eps(\n^\eps,\n^\eps)\xrightarrow[\eps \to 0]{}|\tilde h^\hhh|^2+4 J(\vh)\alpha+(2n+2)\alpha^2\qquad\text{  locally uniformly on $S\setminus S_0$.}
    \end{equation}
\end{lemma}
\begin{proof}
     Let $\E_1,\ldots,\E_{2n-1}$ be as in the proof of \Cref{lemmasecofondformcondueparametri}. First, by \eqref{ricciepsilon} and \eqref{normaleriemannianaepsilon}, 
     \begin{equation*}
         \ric^\eps(\n^\eps,\n^\eps)=-\frac{2}{\eps^2}+\frac{(2n+2)\alpha^2}{1+\eps^2\alpha^2}.
     \end{equation*}
     Moreover,
     \begin{equation*}
         \begin{split}
             \left|h^\eps\right|^2&=\sum_{i,j=1}^{2n-1}h^\eps(\E_i,\E_j)^2+2h^\eps(J(\vh),\s^\eps)^2+2\sum_{\substack{i=1 \\ i\neq n}}^{2n-1}h^\eps(\E_i,\s^\eps)^2+h^\eps(\s^\eps,\s^\eps)^2\\
         \overset{\eqref{formaepsenonepstuttoorizz}}&{=}\frac{|\tilde h^\hhh|^2}{1+\eps^2\alpha^2}+\frac{2}{\eps^2}+\frac{4 J(\vh)\alpha}{1+\eps^2\alpha^2}+\frac{2\eps^2\left(J(\vh)\alpha\right)^2}{\left(1+\eps^2\alpha^2\right)^2}+\frac{2\eps^2}{\left(1+\eps^2\alpha^2\right)^2}\sum_{\substack{i=1 \\ i\neq n}}^{2n-1}\left(\E_i\alpha\right)^2+\frac{\eps^4\left(\s\alpha\right)^2}{\left(1+\eps^2\alpha^2\right)^3}\\
         &=\frac{2}{\eps^2}+\frac{|\tilde h^\hhh|^2+4 J(\vh)\alpha}{1+\eps^2\alpha^2}+\frac{2\eps^2\left|\nabla^{\hhh,S}\alpha\right|^2}{\left(1+\eps^2\alpha^2\right)^2}+\frac{\eps^4\left(\s\alpha\right)^2}{\left(1+\eps^2\alpha^2\right)^3}.
         \end{split}
     \end{equation*}
     The thesis follows combining the above computations.
\end{proof}
\Cref{lemmasecondaformafondriem} allows to compare the Laplace-Beltrami operator with $\hat \Delta^{\hhh,S}$. We need the following lemma.
\begin{lemma}
    Let $p\in S\setminus S_0$. Then 
    \begin{equation}\label{nabepssepsseps}
        \nabla^\eps_{\s^\eps}\s^\eps=-2\alpha J(\vh)-\frac{\eps^2\alpha}{1+\eps^2\alpha^2}\nabla^{\hhh,S}\alpha-\frac{\eps^2\s\alpha}{\left(1+\eps^2\alpha^2\right)^{\frac{3}{2}}}\n^\eps.
    \end{equation}
    \end{lemma}
    \begin{proof}
       First,
        \begin{equation*}
            \left\langle \nabla^\eps_{\s^\eps}\s^\eps,\n^\eps\right\rangle_\eps=-h^\eps(\s^\eps,\s^\eps)\overset{\eqref{formaepsenonepstuttoorizz}}{=}-\frac{\eps^2\s\alpha}{\left(1+\eps^2\alpha^2\right)^{\frac{3}{2}}}.
        \end{equation*}
        Moreover, $\left\langle \nabla^\eps_{\s^\eps}\s^\eps,\s^\eps\right\rangle_\eps=0$. 
        Finally, fix $\E\in\Gamma(\hhh T S)$. Then
        \begin{equation*}
            \begin{split}
                   \Big\langle &\nabla^\eps_{\s^\eps}\s^\eps,\E\Big\rangle_\eps=\frac{1}{\sqrt{1+\eps^2\alpha^2}}\left\langle \nabla^\eps_{-\eps\alpha\vh+\eps T}\left(\frac{-\eps\alpha\vh+\eps T}{\sqrt{1+\eps^2\alpha^2}}\right),\E\right\rangle_\eps\\
                   &=\frac{1}{1+\eps^2\alpha^2}\left\langle \nabla^\eps_{-\eps\alpha\vh+\eps T}\left(-\eps\alpha\vh\right),\E\right\rangle_\eps+\frac{1}{1+\eps^2\alpha^2}\left\langle \nabla^\eps_{-\eps\alpha\vh+\eps T}\eps T,\E\right\rangle_\eps\\
                   \overset{\eqref{connessioneerotazione}}&{=}\frac{\eps^2\alpha^2}{1+\eps^2\alpha^2}\left\langle \nabla^\eps_{\vh}\vh,\E\right\rangle_\eps-\frac{\eps\alpha}{1+\eps^2\alpha^2}\left\langle \nabla^\eps_{\eps T}\vh,\E\right\rangle_\eps-\frac{\alpha}{1+\eps^2\alpha^2}\left\langle \E,J(\vh)\right\rangle\\
                \overset{\eqref{levi_civita},\eqref{propvh5}}&{=}-\frac{2\eps^2\alpha^3}{1+\eps^2\alpha^2}\left\langle \E,J(\vh)\right\rangle-\frac{\eps^2\alpha}{1+\eps^2\alpha^2}\sum_{i=1}^{2n}T(\vh^i)E^i-\frac{\alpha}{1+\eps^2\alpha^2}\sum_{i=1}^n\left(\vh^i Y_i-\vh^{n+i}X_i\right)-\frac{\alpha}{1+\eps^2\alpha^2}\left\langle \E,J(\vh)\right\rangle\\
                \overset{\eqref{normcondist}}&{=}-2\alpha\left\langle \E,J(\vh)\right\rangle-\frac{\eps^2\alpha\E\alpha}{1+\eps^2\alpha^2}.
            \end{split}
        \end{equation*}
        The thesis follows by the above computations.
    \end{proof}
    \begin{proposition}
        Let $ \A\in\Gamma(\hhh T S)$. Let $f\in C^\infty(S\setminus S_0)$. Let $p\in S\setminus S_0$. Then
        \begin{equation}\label{diverepsilonespressa}
            \divv^{\eps,S}\left(\A+f\s\right)=\divv^{\hhh,S}\A+2\alpha \left\langle \A,J(\vh)\right\rangle+\s f-f H^\hhh\alpha+\frac{\eps^2\alpha}{1+\eps^2\alpha^2}\left(\A\alpha+f\s\alpha\right).
        \end{equation}
        In particular, if $\varphi\in\C^\infty(S\setminus S_0)$ and $p\in S\setminus S_0$,
        \begin{equation}\label{laplaepsilonespresso}
            \Delta^{\eps,S}\varphi=\hat \Delta^{\hhh,S}\varphi+\frac{\eps^2\s\s\varphi}{1+\eps^2\alpha^2}-\frac{\eps^4\alpha\s\alpha\s\varphi}{\left(1+\eps^2\alpha^2\right)^2}-\frac{\eps^2H^\hhh\alpha\s\varphi}{1+\eps^2\alpha^2}+\frac{\eps^2\alpha}{1+\eps^2\alpha^2}\nabla^{\hhh,S}\alpha.
        \end{equation}
        Therefore, if $(\varphi^\eps)_\eps$ converges to $\varphi$ smoothly as in \eqref{convhepstoh}, then
        \begin{equation}\label{jacepstojach}
        \jacobi^\eps\varphi^\eps\xrightarrow[\eps\to 0]{}\jacobi^\hhh\varphi,\qquad\text{ locally uniformly on $S\setminus S_0$.}
    \end{equation}
    \end{proposition}
    \begin{proof}
        Let $\E_1,\ldots,\E_{2n-1}$ be as in the proof of \Cref{lemmasecofondformcondueparametri}. First,
        \begin{equation*}
            \begin{split}
                \divv^{\eps,S}\A
                %&=\sum_{i=1}^{2n-1}\left\langle\nabla^\eps_{\E_i}\A,\E_i\right\rangle_\eps+\left\langle\nabla^\eps_{\s^\eps}\A,\s^\eps\right\rangle_\eps\\
                \overset{\eqref{relazionetraleconnessioni}}{=}\sum_{i=1}^{2n-1}\left\langle\nabla_{\E_i}\A,\E_i\right\rangle-\left\langle\nabla^\eps_{\s^\eps}\s^\eps,\A\right\rangle_\eps
                \overset{\eqref{nabepssepsseps}}{=}\divv^{\hhh,S}\A+2\alpha \left\langle \A,J(\vh)\right\rangle+\frac{\eps^2\alpha}{1+\eps^2\alpha^2}\A\alpha.
            \end{split}
        \end{equation*}
        Moreover,
        \begin{equation*}
            \begin{split}
                \divv^{\eps,S}\left(f\s\right)
                %&=\s f+f \divv^{\eps,S}\s\\
                &=\s f+f\sum_{i=1}^{2n-1}\left\langle\nabla^\eps_{\E_i}\s,\E_i\right\rangle_\eps+f\left\langle\nabla^\eps_{\s^\eps}\s,\s^\eps\right\rangle_\eps\\
                &=\s f-f\alpha\sum_{i=1}^{2n-1}\left\langle\nabla^\eps_{\E_i}\vh,\E_i\right\rangle_\eps+f \sum_{i=1}^{2n-1}\left\langle\nabla^\eps_{\E_i} T,\E_i\right\rangle_\eps+f\left\langle\nabla^\eps_{\s^\eps}\left(\frac{\sqrt{1+\eps^2\alpha^2}}{\eps}\s^\eps\right),\s^\eps\right\rangle_\eps\\
                \overset{\eqref{connessioneerotazione},\eqref{relazionetraleconnessioni}}&{=}\s f-f\alpha\sum_{i=1}^{2n-1}\left\langle\nabla_{\E_i}\vh,\E_i\right\rangle+\frac{f}{\eps^2} \sum_{i=1}^{2n-1}\left\langle J(\E_i),\E_i\right\rangle+\frac{f}{\sqrt{1+\eps^2\alpha^2}}\s\left(\sqrt{1+\eps^2\alpha^2}\right)\\
                &=\s f-f H^\hhh\alpha+\frac{\eps^2f\alpha\s\alpha}{1+\eps^2\alpha^2}.
            \end{split}
        \end{equation*}
        To prove \eqref{laplaepsilonespresso} it suffices to apply \eqref{diverepsilonespressa}, noticing that 
        \begin{equation*}
            \nabla^{\eps,S}\varphi=\nabla^{\hhh,S}\varphi+\frac{\eps^2\s\varphi}{1+\eps^2\alpha^2}\s.
        \end{equation*}
        Finally, \eqref{jacepstojach} follows by \eqref{convergenzadelpotenzialejacobi} and \eqref{laplaepsilonespresso}.
    \end{proof}
   The above results and the divergence theorem yield the following integration-by-parts formula (cf. \cite{MR2354992}).
    \begin{proposition}\label{proposizioneibpsubriemsurf}
        Let $\varphi\in C^2(S)\cap C^\infty(S\setminus S_0)$.  Let $\A\in\Gamma(\hhh TS)$. If $\supp \left(\varphi\A\right)\subseteq S\setminus S_0,$ then 
        \begin{align}\label{ibpsubriemequation}
            \int_S\varphi\divv^{\hhh,S}\A\,d\sigma^\hhh+2\int_S\varphi\alpha\left\langle \A,J(\vh)\right\rangle\,d\sigma^\hhh=-\int_S\A\varphi\,d\sigma^\hhh.
        \end{align}
        In addition, if $\psi\in C^2(S)\cap C^\infty(S\setminus S_0)$ and $\supp(\varphi\psi)\subseteq S \setminus S_0$, then 
        \begin{equation}\label{ibpverticalformula}
            \int_S\varphi\,\s\psi\,d\sigma^\hhh=\int_S\varphi\,\psi H^\hhh\alpha\,d\sigma^\hhh-\int_S\s\varphi\,\psi\,d\sigma^\hhh.
        \end{equation}
    \end{proposition}
Finally, we show the convergence of the cubic curvature term appearing in \eqref{variazsecondanormalemaarbitrariaperilresto}.
\begin{lemma}\label{lemmatracehepsterza}
    Let $p\in S\setminus S_0$. Then
    \begin{equation*}
        \begin{split}
            \trace&\left(\left(h^\eps\right)^3\right)=\frac{3\tilde h^\hhh(J(\vh),J(\vh))}{\eps^2\left(1+\eps^2\alpha^2\right)^\frac{1}{2}}+\frac{1}{\left(1+\eps^2\alpha^2\right)^\frac{3}{2}}\left(\trace\left(\left(\tilde h^\hhh\right)^3\right)+6\tilde h^\hhh\left(\nabla^{\hhh,S}\alpha,J(\vh)\right)+3\s\alpha\right)\\
            &\quad+\frac{3\eps^2}{\left(1+\eps^2\alpha^2\right)^\frac{5}{2}}\left(\tilde h^\hhh\left(\nabla^{\hhh,S}\alpha,\nabla^{\hhh,S}\alpha\right)+2J(\vh)\alpha\s\alpha\right)+\frac{3\eps^4}{\left(1+\eps^2\alpha^2\right)^\frac{7}{2}}\s\alpha\left|\nabla^{\hhh,S}\alpha\right|^2+\frac{\eps^6}{\left(1+\eps^2\alpha^2\right)^\frac{9}{2}}\left(\s\alpha\right)^3
        \end{split}
    \end{equation*}
\end{lemma}
\begin{proof}
     Let $\E_1,\ldots,\E_{2n-1}$ be as in the proof of \Cref{lemmasecofondformcondueparametri}. Set $h^\eps_{ij}=h^\eps(\E_i,\E_j)$ for $i,j=1,\ldots,2n$. Then
      \begin{equation*}
          \begin{split}
              \trace\left(\left(h^\eps\right)^3\right)%&=\sum_{i,j,k=1}^{2n}h^\eps_{ij}h^\eps_{jk}h^\eps_{ki}\\
             % &=\sum_{\substack{i,j,k=1 \\ k\neq 2n}}^{2n}h^\eps_{ij}h^\eps_{jk}h^\eps_{ki}+\sum_{i,j=1}^{2n}h^\eps_{ij}h^\eps_{j,2n}h^\eps_{2n,i}\\
             % &=\sum_{\substack{i,j,k=1 \\ j,k\neq 2n}}^{2n}h^\eps_{ij}h^\eps_{jk}h^\eps_{ki}+\sum_{\substack{i,k=1 \\ k\neq 2n}}^{2n}h^\eps_{i,2n}h^\eps_{2n,k}h^\eps_{ki}+\sum_{\substack{i,j=1 \\ j\neq 2n}}^{2n}h^\eps_{ij}h^\eps_{j,2n}h^\eps_{2n,i}+\sum_{i=1}^{2n}h^\eps_{i,2n}h^\eps_{2n,2n}h^\eps_{2n,i}\\
      %        &=\sum_{\substack{i,j,k=1 \\ i,j,k\neq 2n}}^{2n}h^\eps_{ij}h^\eps_{jk}h^\eps_{ki}+\sum_{\substack{i,j=1 \\ i,j\neq 2n}}^{2n}h^\eps_{ij}h^\eps_{j,2n}h^\eps_{2n,i}+\sum_{\substack{j,k=1 \\ j,k\neq 2n}}^{2n}h^\eps_{2n,j}h^\eps_{jk}h^\eps_{k,2n}+\sum_{\substack{i,k=1 \\ i,k\neq 2n}}^{2n}h^\eps_{i,2n}h^\eps_{2n,k}h^\eps_{ki}\\
      %        &\quad+\sum_{\substack{i=1 \\ i\neq 2n}}^{2n}h^\eps_{i,2n}h^\eps_{2n,2n}h^\eps_{2n,i}+\sum_{\substack{j=1 \\ j\neq 2n}}^{2n}h^\eps_{2n,j}h^\eps_{j,2n}h^\eps_{2n,2n}+\sum_{\substack{k=1 \\ k\neq 2n}}^{2n}h^\eps_{2n,2n}h^\eps_{2n,k}h^\eps_{k,2n}+\left(h^\eps_{2n,2n}\right)^3\\
              =\underbrace{\sum_{i,j,k=1 }^{2n-1}h^\eps_{ij}h^\eps_{jk}h^\eps_{ki}}_\mathrm{I}+\underbrace{3\sum_{i,j=1 }^{2n-1}h^\eps_{ij}h^\eps_{i,2n}h^\eps_{j,2n}}_\mathrm{II}+\underbrace{3h^\eps_{2n,2n}\sum_{i=1}^{2n-1}\left(h^\eps_{i,2n}\right)^2}_\mathrm{III}+\underbrace{\left(h^\eps_{2n,2n}\right)^3}_\mathrm{IV}.
          \end{split}
      \end{equation*}
      First,
      \begin{equation*}
          \mathrm{I}\overset{\eqref{formaepsenonepstuttoorizz}}{=}\frac{1}{\left(1+\eps^2\alpha^2\right)^\frac{3}{2}}\trace\left(\left(\tilde h^\hhh\right)^3\right).
      \end{equation*}
      Moreover,
      \begin{equation*}
          \begin{split}
              \mathrm{II}%&=3\sum_{\substack{i,j=1 \\  j\neq n} }^{2n-1}h^\eps_{ij}h^\eps_{i,2n}h^\eps_{j,2n}+3\sum_{i=1 }^{2n-1}h^\eps_{i,n}h^\eps_{i,2n}h^\eps_{n,2n}\\
  %            &=3\sum_{\substack{i,j=1 \\ i,j\neq n} }^{2n-1}h^\eps_{ij}h^\eps_{i,2n}h^\eps_{j,2n}+3\sum_{\substack{j=1 \\ j\neq n} }^{2n-1}h^\eps_{n,j}h^\eps_{n,2n}h^\eps_{j,2n}+3\sum_{\substack{i=1 \\  i\neq n}}^{2n-1}h^\eps_{i,n}h^\eps_{i,2n}h^\eps_{n,2n}+3h^\eps_{n,n}\left(h^\eps_{n,2n}\right)^2\\
              &=3\sum_{\substack{i,j=1 \\ i,j\neq n} }^{2n-1}h^\eps_{ij}h^\eps_{i,2n}h^\eps_{j,2n}+6h^\eps_{n,2n}\sum_{\substack{i=1  \\ i\neq n}}^{2n-1}h^\eps_{i,n}h^\eps_{i,2n}+3h^\eps_{n,n}\left(h^\eps_{n,2n}\right)^2\\
              \overset{\eqref{formaepsenonepstuttoorizz}}&{=}\frac{3\eps^2}{\left(1+\eps^2\alpha^2\right)^\frac{5}{2}}\sum_{\substack{i,j=1 \\ i,j\neq n} }^{2n-1}\tilde h^\hhh(\E_i,\E_j)\E_i\alpha\E_j\alpha+6\left(1+\frac{\eps^2J(\vh)\alpha}{1+\eps^2\alpha^2}\right)\frac{1}{\left(1+\eps^2\alpha^2\right)^\frac{3}{2}}\sum_{\substack{i=1  \\ i\neq n}}^{2n-1}\tilde h^\hhh (\E_i,J(\vh))\E_i\alpha\\
              &\quad+\frac{3}{\eps^2\left(1+\eps^2\alpha^2\right)^\frac{1}{2}}\tilde h^\hhh(J(\vh),J(\vh))\left(1+\frac{\eps^2J(\vh)\alpha}{1+\eps^2\alpha^2}\right)^2\\
              &=\frac{3\eps^2}{\left(1+\eps^2\alpha^2\right)^\frac{5}{2}}\tilde h^\hhh\left(\nabla^{\hhh,S}\alpha,\nabla^{\hhh,S}\alpha\right)+\frac{6}{\left(1+\eps^2\alpha^2\right)^\frac{3}{2}}\tilde h^\hhh\left(\nabla^{\hhh,S}\alpha,J(\vh)\right)+\frac{3}{\eps^2\left(1+\eps^2\alpha^2\right)^\frac{1}{2}}\tilde h^\hhh(J(\vh),J(\vh)).
          \end{split}
      \end{equation*}
      In addition,
      \begin{equation*}
          \begin{split}
              \mathrm{III}&=3h^\eps_{2n,2n}\sum_{\substack{i=1 \\ i\neq n}}^{2n-1}\left(h^\eps_{i,2n}\right)^2+3h^\eps_{2n,2n}\left(h^\eps_{n,2n}\right)^2\\
              \overset{\eqref{formaepsenonepstuttoorizz}}&{=}\frac{3\eps^4}{\left(1+\eps^2\alpha^2\right)^\frac{7}{2}}\s\alpha\sum_{\substack{i=1 \\ i\neq n}}^{2n-1}\left(\E_i\alpha\right)^2+\frac{3}{\left(1+\eps^2\alpha^2\right)^\frac{3}{2}}\s\alpha\left(1+\frac{\eps^2J(\vh)\alpha}{1+\eps^2\alpha^2}\right)^2\\
              &=\frac{3\eps^4}{\left(1+\eps^2\alpha^2\right)^\frac{7}{2}}\s\alpha\left|\nabla^{\hhh,S}\alpha\right|^2+\frac{6\eps^2}{\left(1+\eps^2\alpha^2\right)^\frac{5}{2}}\s\alpha J(\vh)\alpha+\frac{3}{\left(1+\eps^2\alpha^2\right)^\frac{3}{2}}\s\alpha.
          \end{split}
      \end{equation*}
      Finally,
      \begin{equation*}
          \mathrm{IV}\overset{\eqref{formaepsenonepstuttoorizz}}{=}\frac{\eps^6}{\left(1+\eps^2\alpha^2\right)^\frac{9}{2}}\left(\s\alpha\right)^3.
      \end{equation*}
      The thesis follows combining the above computations.
\end{proof}
\begin{lemma}\label{lemmadiscalprodhconcnn}
    Let $p\in S\setminus S_0$. Then
    \begin{equation*}
        \left\langle h^\eps,\j^\eps(\n^\eps,\n^\eps)\right\rangle_\eps=\frac{3\tilde h^\hhh(J(\vh),J(\vh))}{\eps^2\left(1+\eps^2\alpha^2\right)^\frac{3}{2}}-\frac{1}{\left(1+\eps^2\alpha^2\right)^\frac{3}{2}}\Big(\s\alpha+H^\hhh\alpha^2\Big).
    \end{equation*}
\end{lemma}
\begin{proof}
        Let $\E_1,\ldots,\E_{2n-1}$ be as in the proof of \Cref{lemmasecofondformcondueparametri}. By \eqref{riemannepsilonquattrozero} and \eqref{normaleriemannianaepsilon}, if $\B,\D\in\Gamma(TS)$,
         \begin{equation}\label{formadicnnepsilon}
        \begin{split}
            \j^\eps(\n^\eps,\n^\eps)(\B,\D)&=\frac{3}{\eps^2}\left\langle \n^\eps,J(\B)\right\rangle _\eps \left\langle \n^\eps,J(\D)\right\rangle_\eps-\frac{B^{2n+1,\eps}D^{2n+1,\eps}}{\eps^2}-\frac{\alpha^2}{1+\eps^2\alpha^2}\langle\B,\D\rangle_\eps.
        \end{split}
    \end{equation}
    Therefore, 
    \begin{equation*}
        \begin{split}
              \left\langle h^\eps,\j^\eps(\n^\eps,\n^\eps)\right\rangle_\eps&=\sum_{i,j=1}^{2n}h^\eps(\E_i,\E_j)\j^\eps(\n^\eps,\n^\eps)(\E_i,\E_j)\\
              \overset{\eqref{formadicnnepsilon}}&{=}\frac{3}{\eps^2}h^\eps (J(\vh),J(\vh))\left\langle \n^\eps,\vh\right\rangle_\eps^2-\frac{1}{\eps^2}h^\eps(\s^\eps,\s^\eps)\left\langle\s^\eps,\eps T\right\rangle_\eps^2-\frac{\alpha^2}{1+\eps^2\alpha^2}H^\eps\\
              \overset{\eqref{normaleriemannianaepsilon},\eqref{formadiesseepsilon},\eqref{formaepsenonepstuttoorizz},\eqref{curvepsecurvh}}&{=}\frac{3\tilde h^\hhh(J(\vh),J(\vh))}{\eps^2\left(1+\eps^2\alpha^2\right)^\frac{3}{2}}-\frac{\s\alpha}{\left(1+\eps^2\alpha^2\right)^\frac{5}{2}}-\frac{H^\hhh\alpha^2}{\left(1+\eps^2\alpha^2\right)^\frac{3}{2}}--\frac{\eps^2\alpha^2\s\alpha}{\left(1+\eps^2\alpha^2\right)^\frac{5}{2}}\\
              &=\frac{3\tilde h^\hhh(J(\vh),J(\vh))}{\eps^2\left(1+\eps^2\alpha^2\right)^\frac{3}{2}}-\frac{\s\alpha}{\left(1+\eps^2\alpha^2\right)^\frac{3}{2}}-\frac{H^\hhh\alpha^2}{\left(1+\eps^2\alpha^2\right)^\frac{3}{2}}.
        \end{split}
    \end{equation*}
\end{proof}
    \begin{theorem}
        Let $\eps>0$. Let $p\in S\setminus S_0$. 
        Then
        \begin{equation*}
            \begin{split}
                2&\trace\left(\left(h^\eps\right)^3\right)-2 \left\langle h^\eps,\j^\eps(\n^\eps,\n^\eps)\right\rangle_\eps-\left(\nabla^\eps_{\n^\eps}\ric^\eps\right)(\n^\eps,\n^\eps)\\
                &=\frac{1}{\left(1+\eps^2\alpha^2\right)^\frac{3}{2}}\left(2\trace\left(\left(\tilde h^\hhh\right)^3\right)+12\tilde h^\hhh\left(\nabla^{\hhh,S}\alpha,J(\vh)\right)+8\s\alpha+6\tilde h^\hhh(J(\vh),J(\vh))\alpha^2+2 H^\hhh\alpha^2\right)\\
                &\quad +\frac{6\eps^2}{\left(1+\eps^2\alpha^2\right)^\frac{5}{2}}\left(\tilde h^\hhh\left(\nabla^{\hhh,S}\alpha,\nabla^{\hhh,S}\alpha\right)+2J(\vh)\alpha\s\alpha\right)+\frac{6\eps^4}{\left(1+\eps^2\alpha^2\right)^\frac{7}{2}}\s\alpha\left|\nabla^{\hhh,S}\alpha\right|^2+\frac{2\eps^6}{\left(1+\eps^2\alpha^2\right)^\frac{9}{2}}\left(\s\alpha\right)^3.
            \end{split}
        \end{equation*}
        In particular,
        \begin{equation}\label{qbqequation}
            \begin{split}
              2&\trace\left(\left(h^\eps\right)^3\right)-2 \left\langle h^\eps,\j^\eps(\n^\eps,\n^\eps)\right\rangle_\eps-\left(\nabla^\eps_{\n^\eps}\ric^\eps\right)(\n^\eps,\n^\eps)\\
                &\xrightarrow[\eps\to 0]{}  2\trace\left(\left(\tilde h^\hhh\right)^3\right)+12\tilde h^\hhh\left(\nabla^{\hhh,S}\alpha,J(\vh)\right)+8\s\alpha+6\tilde h^\hhh(J(\vh),J(\vh))\alpha^2+2 H^\hhh\alpha^2
            \end{split}
        \end{equation}
        locally uniformly on $S\setminus S_0$.
    \end{theorem}
    \begin{proof}
    Notice that
    \begin{equation*}
        \begin{split}
            \frac{6\tilde h^\hhh(J(\vh),J(\vh))}{\eps^2\left(1+\eps^2\alpha^2\right)^\frac{1}{2}}-\frac{6\tilde h^\hhh(J(\vh),J(\vh))}{\eps^2\left(1+\eps^2\alpha^2\right)^\frac{3}{2}}=\frac{6\tilde h^\hhh(J(\vh),J(\vh))\alpha^2}{\left(1+\eps^2\alpha^2\right)^\frac{3}{2}}.
        \end{split}
    \end{equation*}
        The thesis follows by \eqref{nabaepsriccizerosudiagonale}, \Cref{lemmatracehepsterza} and \Cref{lemmadiscalprodhconcnn}.
    \end{proof}
    \subsection{Variation formulas}
      We establish the relevant variation formulas for \eqref{tmcfunsubriemderfgt5ryhythyh}. Accordingly, we adapt some notation of \Cref{subsec_variations} to the sub-Riemannian setting. If $\Phi$ is a variation and $\Y$ is as in \eqref{def_xechis}, we set 
      \begin{equation*}
    \X(q)=\Y(0,q),\qquad \Z(q)=\left.\frac{\partial }{\partial \tau}\right|_{\tau=0}\Y(\tau,q)+\left(\nabla_\X\X\right)(q)\qquad\text{for any $q\in \hh^n$,}
\end{equation*}
where in the above definition $\nabla$ is the pseudohermitian connection \eqref{pseudotorsion}.
According to the Riemannian terminology, we call $\X$ and $\Z$ respectively  \emph{variational velocity field} and \emph{variational acceleration field} of $\Phi$. If a hypersurface $S$ is fixed, we say that a variation is \emph{non-characteristic} whenever $S\cap K(\Phi)\Subset S\setminus S_0$. In the next result, we restrict ourselves to consider non-characteristic variations 
of closed hypersurfaces. If instead one computes variations on non-compact hypersurfaces, it suffices to restrict the relevant functionals to $K(\Phi)$. 
%As it will be clearer in a while, the first variation formula of $\tmc^\hhh_f$ depends only on $\X|_S$, while the second variation depends only on $\X|_S$ and $\Z|_S$, accordingly, we define
%\begin{equation*}
  %%  \delta\tmc^\hhh_f(S)[\X]\coloneqq \left.\frac{d}{d\tau}\right|_{\tau=0}\tmc^\hhh\left(\Phi_\tau(S)\right),\qquad   \delta^2\tmc^\hhh_f(S)[\X,\Z]\coloneqq \left.\frac{d^2}{d\tau^2}\right|_{\tau=0}\tmc\left(\Phi^\hhh_\tau(S)\right).
%\end{equation*}  
\begin{theorem}\label{teoremavariazionigeneralimainsubriem}
        Let $S\subseteq\hh^n$ be an embedded, closed, two-sided hypersurface of class $C^2$. Let $S\setminus S_0$ be smooth. Fix a function $f:\rr\to\rr$ which is smooth in a neighborhood of $\left\{H^\hhh(p)\,:\,p\in S\setminus S_0\right\}$. Let $\Phi$ be a non-characteristic variation.  Denote by $\X$ and $\Z$ its velocity and acceleration respectively. Define $\varphi,\psi\in C^\infty_c(S\setminus S_0)$ by
        \begin{equation*}\label{lefunzioniimportanti}
        \begin{split}        \varphi&\coloneqq\left\langle \X,\vh+\alpha T\right\rangle,\\
            \psi&\coloneqq \left\langle\Z,\vh+\alpha T\right\rangle-2\left(\X-\varphi\vh\right)\varphi-\Big\langle\nabla_{\X-\varphi\vh}\left(\X-\varphi\vh\right),\vh+\alpha T\Big\rangle-2\alpha\varphi\left\langle \X,J(\vh)\right\rangle.
        \end{split}
        \end{equation*}
               Set 
        \begin{equation*}
 \delta\tmc^\hhh_f(S)[\Phi]\coloneqq \left.\frac{d}{d\tau}\right|_{\tau=0}\tmc^\hhh_f\left(\Phi_\tau(S)\right),\qquad   \delta^2\tmc^\hhh_f(S)[\Phi]\coloneqq \left.\frac{d^2}{d\tau^2}\right|_{\tau=0}\tmc^\hhh_f\left(\Phi_\tau(S)\right).
\end{equation*} 
        Then 
            \begin{align}            \delta\tmc^\hhh_f(S)[\varphi]\coloneqq\delta\tmc^\hhh_f(S)[\Phi]&=  \int_S\varphi\Big(\jacobi^\hhh f'(H^\hhh)+f(H^\hhh) H^\hhh\Big)\,d\sigma^\hhh,\label{variazprimasubriemmain}\\
            \delta^2\tmc_f^\hhh(S)[\Phi]&=\delta\tmc^\hhh_f(S)\left[\psi\right]+\int_S\varphi\mathcal L^\hhh\varphi\,d\sigma^\hhh, \label{variazionesecondasubriemannianamain}
    \end{align}
    where $ \mathcal L^\hhh$ is the self-adjoint operator defined on $S\setminus S_0$ by
    \begin{equation*}
        \begin{split}
            \mathcal L^\hhh\varphi&=\jacobi^\hhh\left(f''(H^\hhh)\jacobi^\hhh\varphi\right)\\
            &\quad+2\divv^{\hhh,S} \left\langle f'(H^\hhh)A^\hhh\left(\nabla^{\hhh,S}\varphi\right)\right)+4 f'(H^\hhh)\s J(\vh)\varphi-\left(f'(H^\hhh)H^\hhh+f(H^\hhh)\right)\hat\Delta^{\hhh,S}\varphi\\
            &\quad+4\alpha f'(H^\hhh)\Big(\tilde h^\hhh\left(\nabla^{\hhh,S}\varphi,J(\vh)\right)-H^\hhh J(\vh)\varphi\Big)- 4 f''(H^\hhh)\Big(\tilde h^\hhh\left(\nabla^{\hhh,S}H^\hhh,\nabla^{\hhh,S}\varphi\right)+ \s \varphi J(\vh)H^\hhh\Big)\\
            &\quad+\left(f''(H^\hhh)H^\hhh-2f'(H^\hhh)\right)\left\langle \nabla^{\hhh,S} H^\hhh,\nabla^{\hhh,S}\varphi\right\rangle\\
            &\quad+\varphi f'(H^\hhh)\Big( 2\trace\left(\left(\tilde h^\hhh\right)^3\right)+12\tilde h^\hhh\left(\nabla^{\hhh,S}\alpha,J(\vh)\right)+8\s\alpha+6\tilde h^\hhh(J(\vh),J(\vh))\alpha^2+2 H^\hhh\alpha^2\Big)\\
            &\quad+\varphi \Big ( f(H^\hhh)\left( H^\hhh\right)^2-\left(2f'(H^\hhh) H^\hhh+f(H^\hhh)\right)\left(|\tilde h^\hhh|^2+4J(\vh)\alpha+(2n+2)\alpha^2\right)\Big).
        \end{split}
    \end{equation*}
   Here $\jacobi^\hhh$ and $A^\hhh$ are the horizontal Jacobi operator and the horizontal shape operator (cf. \Cref{subsechypersurf}).
    \end{theorem} 
    \begin{proof}
        Fix $\eps>0$. 
 Set $   \tmc^\eps_f(S)\coloneqq\int_S f\left(H^\eps\right)\,d\sigma^\eps. $
Then \begin{equation}\label{confrontoduefunzionali}
       \begin{split}
           \tmc^\eps_f(S)\overset{\eqref{areaelementeps},\eqref{curvepsecurvh}}&{=}\int_S f\left(\frac{H^\hhh}{\sqrt{1+\eps^2\alpha^2}}+\frac{\eps^2\s\alpha}{\left(1+\eps^2\alpha^2\right)^{\frac{3}{2}}}\right)\frac{\sqrt{1+\eps^2\alpha^2}}{\eps\sqrt{1+\alpha^2}}\,d\sigma^1,\\
       \tmc^\hhh_f(S)\overset{\eqref{areaelementorizz}}&{=}\int_S f\left(H^\hhh\right)\frac{1}{\sqrt{1+\alpha^2}}\,d\sigma^1.\\
       \end{split}
\end{equation}
Therefore, by \eqref{confrontoduefunzionali} and since $\Phi$ is supported on $S\setminus S_0$, we infer that 
\begin{equation}\label{convergenzadellevariazioniprime}
    \delta\tmc^\hhh_f(S)[\Phi]=\lim_{\eps\to 0}\eps  \delta\tmc^\eps_f(S)[\Phi],\qquad \delta^2\tmc^\hhh_f(S)[\Phi]=\lim_{\eps\to 0}\eps  \delta^2\tmc^\eps_f(S)[\Phi].
\end{equation}
      Set $\varphi^\eps=\left\langle\X,\n^\eps\right\rangle_\eps$. Then, by \eqref{primavarfunzgen},
\begin{equation*}
      \eps\delta\tmc_f^\eps(S)[\Phi]=  \int_S\Big(f'(H^\eps)\left(-\Delta^{\eps,S}\varphi^\eps-\varphi^\eps\left(|h^\eps|^2+\ric^\eps(\n^\eps,\n^\eps)\right)\right)+\varphi^\eps f(H^\eps) H^\eps\Big)\frac{\sqrt{1+\eps^2\alpha^2}}{\sqrt{1+\alpha^2}}\,d\sigma^1.
    \end{equation*}
    Notice that 
    \begin{equation*}
        \varphi^\eps=\left\langle\X,\n^\eps\right\rangle=\frac{\left\langle \X,\vh\right\rangle}{\sqrt{1+\eps^2\alpha^2}}+\frac{\eps\alpha\left\langle \X,\eps T\right\rangle_\eps}{\sqrt{1+\eps^2\alpha^2}}=\frac{\left\langle \X,\vh\right\rangle}{\sqrt{1+\eps^2\alpha^2}}+\frac{\alpha\left\langle \X, T\right\rangle_1}{\sqrt{1+\eps^2\alpha^2}}=\frac{\varphi}{\sqrt{1+\eps^2\alpha^2}}.
    \end{equation*}
    Therefore, \eqref{laplaepsilonespresso} implies that 
    \begin{equation*}
       \varphi^\eps\xrightarrow[\eps \to 0]{}\varphi,\qquad \Delta^{\eps,S}\varphi^\eps\xrightarrow[\eps\to 0]{}\hat\Delta^{\hhh,S}\varphi\qquad\text{uniformly on $S$}.
    \end{equation*}
    Then, by \eqref{convergenzadellevariazioniprime} and recalling \eqref{convhepstoh} and \eqref{convergenzadelpotenzialejacobi}, \eqref{variazprimasubriemmain} follows. Next, we prove \eqref{variazionesecondasubriemannianamain}. Set $$\X^{T,\eps}=\X-\varphi^\eps\n^\eps=\X-\frac{\varphi}{1+\eps^2\alpha^2}\left(\vh+\eps^2\alpha T\right),$$ and denote by $\Z^\eps$ the acceleration of $\Phi$ with respect to $\nabla^\eps$.
 % and $\Phi^\eps$ is supported in $S\setminus S_0$ uniformly with respect to $\eps>0$ (cf. \eqref{exponentialvariation}). Notice that  
 By \eqref{variazsecondanormalemaarbitrariaperilresto}, 
   \begin{equation*}
        \begin{split}
            \delta^2\tmc^\eps_f(S)&[\Phi]=\delta\tmc^\eps_f(S)\left[\left\langle\Z|_S,\n^\eps\right\rangle_\eps-2\X^{T,\eps}\varphi^\eps +h^\eps\left(\X^{T,\eps},\X^{T,\eps}\right)\right]+\int_S\varphi^\eps\,\jacobi^\eps\left(f''(H^\eps)\jacobi^\eps\varphi^\eps\right)\,d\sigma^\eps\\
            &\quad+\int_S2\varphi^\eps \divv^{\eps,S} \left(f'(H^\eps)A^\eps\left(\nabla^{\eps,S}\varphi^\eps\right)\right)\,d\sigma^\eps-\int_S\varphi^\eps\left(f'(H^\eps)H^\eps+f(H^\eps)\right)\Delta^{\eps,S}\varphi^\eps\,d\sigma^\eps\\
            &\quad-\int_S4\varphi ^\eps f''(H^\eps)h^\eps\left(\nabla^{\eps,S} H^\eps,\nabla^{\eps,S}\varphi^\eps\right)\,d\sigma^\eps+\int_S\varphi^\eps\left(f''(H^\eps)H^\eps-2f'(H^\eps)\right)\left\langle \nabla^{\eps,S} H^\eps,\nabla^{\eps,S}\varphi^\eps\right\rangle_\eps\,d\sigma^\eps\\
            &\quad+\int_S\left(\varphi ^\eps\right)^2f'(H^\eps)\Big(2\trace\left(\left(h^\eps\right)^3\right)-2\left\langle h^\eps,\j^\eps(\n^\eps,\n^\eps)\right\rangle-\left(\nabla^\eps_{\n^\eps} \ric^\eps\right)(\n^\eps,\n^\eps)\Big)\,d\sigma^\eps\\
            &\quad+\int_S\left(\varphi^\eps\right)^2 \Big ( f(H^\eps)\left( H^\eps\right)^2-\left(2f'(H^\eps) H^\eps+f(H^\eps)\right)\left(|h^\eps|^2+ \ric^\eps(\n^\eps,\n^\eps)\right)\Big)\,d\sigma ^\eps.
        \end{split}
    \end{equation*}
    Denote the terms on the right-hand side by $\mathrm{I^\eps},\ldots,\mathrm{VIII^\eps}$. First, 
    \begin{equation*}
        \left\langle\Z^\eps,\n^\eps\right\rangle_\eps=\frac{\left\langle\X',\vh+\alpha T\right\rangle_1}{\sqrt{1+\eps^2\alpha^2}}+\left\langle \nabla^\eps_\X\X,\n^\eps\right\rangle_\eps\overset{\eqref{levicivitavspseudohermitian}}{=}\frac{\left\langle\Z,\vh+\alpha T\right\rangle_1}{\sqrt{1+\eps^2\alpha^2}}+\frac{2\left\langle\X,\eps T\right\rangle_\eps}{\eps}\left\langle J(\X),\n^\eps\right\rangle_\eps.
    \end{equation*}
    Moreover, $\X-\varphi\vh\in\Gamma(TS)$, and 
    \begin{equation*}
        -2\X^{T,\eps}\varphi^\eps=-2\left(\X-\frac{\varphi}{1+\eps^2\alpha^2}\left(\vh+\eps^2\alpha T\right)\right)\left(\frac{\varphi}{\sqrt{1+\eps^2\alpha^2}}\right)\xrightarrow[\eps\to 0]{}-2\left(\X-\varphi\vh\right)\varphi\qquad\text{uniformly on $S$}.
    \end{equation*}
    In addition,
    \begin{equation*}
        h^\eps\left(\X^{T,\eps},\X^{T,\eps}\right)=-\left\langle\nabla^\eps_{\X^{T,\eps}}\X^{T,\eps},\n^\eps\right\rangle_\eps\overset{\eqref{levicivitavspseudohermitian}}{=}-\left\langle\nabla_{\X^{T,\eps}}\X^{T,\eps},\n^\eps\right\rangle_\eps-\frac{2\left\langle\X^{T,\eps},\eps T\right\rangle_\eps}{\eps}\left\langle J(\X),\n^\eps\right\rangle_\eps.
    \end{equation*}
Therefore
\begin{equation*}
    \begin{split}
        \left\langle\Z^\eps,\n^\eps\right\rangle_\eps+h^\eps\left(\X^{T,\eps},\X^{T,\eps}\right)&=\frac{\left\langle\Z,\vh+\alpha T\right\rangle_1}{\sqrt{1+\eps^2\alpha^2}}-\left\langle\nabla_{\X^{T,\eps}}\X^{T,\eps},\n^\eps\right\rangle_\eps+\frac{2\left\langle J(\X),\n^\eps\right\rangle_\eps}{\eps}\left\langle\X-\X^{T,\eps},\eps T\right\rangle_\eps\\
        &=\frac{\left\langle\Z,\vh+\alpha T\right\rangle_1}{\sqrt{1+\eps^2\alpha^2}}-\left\langle\nabla_{\X^{T,\eps}}\X^{T,\eps},\n^\eps\right\rangle_\eps+\frac{2\alpha\varphi_\epsilon\left\langle J(\X),\n^\eps\right\rangle_\eps}{\sqrt{1+\eps^2\alpha^2}}\\
        &\xrightarrow[\eps\to 0]{}\left\langle\Z,\vh+\alpha T\right\rangle_1-\Big\langle\nabla_{\X-\varphi\vh}\left(\X-\varphi\vh\right),\vh+\alpha T\Big\rangle_1-2\alpha\varphi\left\langle \X,J(\vh)\right\rangle
    \end{split}
\end{equation*}
uniformly on $S$. By the first part of the proof, we conclude that 
    \begin{equation*}
        \eps\mathrm{I^\eps}\xrightarrow[\eps\to 0]{}\delta\tmc^\hhh_f(S)\left[\left\langle\Z,\vh+\alpha T\right\rangle_1-2\left(\X-\varphi\vh\right)\varphi-\Big\langle\nabla_{\X-\varphi\vh}\left(\X-\varphi\vh\right),\vh+\alpha T\Big\rangle_1-2\alpha\varphi\left\langle \X,J(\vh)\right\rangle\right].
    \end{equation*}
    Moreover, by \eqref{jacepstojach},
    \begin{equation*}
        \eps\mathrm{II^\eps}=\eps \int_Sf''(H^\eps)\jacobi^\eps\varphi^\eps\,\jacobi^\eps\varphi^\eps\,d\sigma^\eps\xrightarrow[\eps\to 0]{}\int_Sf''(H^\hhh)\jacobi^\hhh\varphi\,\jacobi^\hhh\varphi\,d\sigma^\hhh=\int_S\varphi\,\jacobi^\hhh\left(f''(H^\hhh)\jacobi^\hhh\varphi\right)\,d\sigma^\hhh.
    \end{equation*}
    Next,
    \begin{equation*}
        \begin{split}
            \eps\mathrm{III^\eps}&=-\eps\int_S 2 f'(H^\eps)h^\eps\left(\nabla^{\eps,S}\varphi^\eps,\nabla^{\eps,S}\varphi^\eps\right)\,d\sigma^\eps\\
            \overset{\eqref{convhepsgradeps}}&{\xrightarrow[\eps\to 0]{}}-\int_S 2f'(H^\hhh)\Big(\tilde h^\hhh\left(\nabla^{\hhh,S}\varphi,\nabla^{\hhh,S}\varphi\right)+ 2J(\vh)\varphi\s\varphi\Big)\,d\sigma^\hhh\\
            \overset{\eqref{rapptrahetildacca2026}}&{=}
            -\int_S 2 \left\langle f'(H^\hhh)A^\hhh\left(\nabla^{\hhh,S}\varphi\right),\nabla^{\hhh,S}\varphi\right\rangle\,d\sigma^\hhh -\int_S 4f'(H^\hhh )J(\vh)\varphi\s\varphi\,d\sigma^\hhh\\
            \overset{\eqref{ibpsubriemequation},\eqref{ibpverticalformula}}&{=} \int_S 2\varphi\divv^{\hhh,S} \left\langle f'(H^\hhh)A^\hhh\left(\nabla^{\hhh,S}\varphi\right)\right)\,d\sigma^\hhh+\int_S4\varphi\alpha f'(H^\hhh)h^\hhh\left(\nabla^{\hhh,S}\varphi,J(\vh)\right)\,d\sigma^\hhh\\
            &\quad+\int_S 4\varphi J(\vh)\varphi f''(H^\hhh)\s H^\hhh\,d\sigma^\hhh+\int_S 4\varphi f'(H^\hhh)\s J(\vh)\varphi,d\sigma^\hhh-\int_S4\varphi\alpha f'(H^\hhh)H^\hhh J(\vh)\varphi\,d\sigma^\hhh\\
            \overset{\eqref{rapptrahetildacca2026}}&{=}\int_S 2\varphi\divv^{\hhh,S} \left\langle f'(H^\hhh)A^\hhh\left(\nabla^{\hhh,S}\varphi\right)\right)\,d\sigma^\hhh+\int_S4\varphi\alpha f'(H^\hhh)\Big(\tilde h^\hhh\left(\nabla^{\hhh,S}\varphi,J(\vh)\right)-H^\hhh J(\vh)\varphi\Big)\,d\sigma^\hhh\\
            &\quad+\int_S 4\varphi J(\vh)\varphi f''(H^\hhh)\s H^\hhh\,d\sigma^\hhh+\int_S 4\varphi f'(H^\hhh)\s J(\vh)\varphi,d\sigma^\hhh.
        \end{split}
    \end{equation*}
    Moreover,
    \begin{equation*}
        \eps\mathrm{IV}^\eps\xrightarrow[\eps\to 0]{}-\int_S\varphi\left(f'(H^\hhh)H^\hhh+f(H^\hhh)\right)\hat\Delta^{\hhh,S}\varphi\,d\sigma^\hhh.
    \end{equation*}
    In addition,
    \begin{equation*}
       \eps \mathrm{V}^\eps   \overset{\eqref{convhepsgradeps}}{\xrightarrow[\eps\to 0]{}}-\int_S 4\varphi f''(H^\hhh)\Big(\tilde h^\hhh\left(\nabla^{\hhh,S}H^\hhh,\nabla^{\hhh,S}\varphi\right)+ J(\vh)\varphi\s H^\hhh+ \s \varphi J(\vh)H^\hhh\Big)\,d\sigma^\hhh.
    \end{equation*}
    Moreover, 
    \begin{equation*}
        \eps\mathrm{VI^\eps}\xrightarrow[\eps\to 0]{}\int_S\varphi\left(f''(H^\hhh)H^\hhh-2f'(H^\hhh)\right)\left\langle \nabla^{\hhh,S} H^\hhh,\nabla^{\hhh,S}\varphi\right\rangle\,d\sigma^\hhh.
    \end{equation*}
    Furthermore, by \eqref{qbqequation},
    \begin{equation*}
        \eps\mathrm{VII}^\eps\xrightarrow[\eps\to 0]{}\int_S\varphi^2 f'(H^\hhh)\Big( 2\trace\left(\left(\tilde h^\hhh\right)^3\right)+12\tilde h^\hhh\left(\nabla^{\hhh,S}\alpha,J(\vh)\right)+8\s\alpha+6\tilde h^\hhh(J(\vh),J(\vh))\alpha^2+2 H^\hhh\alpha^2\Big)\,d\sigma^\hhh.
    \end{equation*}
    Finally, by \eqref{convergenzadelpotenzialejacobi},
    \begin{equation*}
        \eps\mathrm{VIII}^\eps\xrightarrow[\eps\to 0]{}\int_S\varphi^2 \Big ( f(H^\hhh)\left( H^\hhh\right)^2-\left(2f'(H^\hhh) H^\hhh+f(H^\hhh)\right)\left(|\tilde h^\hhh|^2+4J(\vh)\alpha+(2n+2)\alpha^2\right)\Big)\,d\sigma^\hhh.
    \end{equation*}
  The thesis follows combining the above computations.
 %   Then, combining the above computations,
  %  \begin{equation*}
   %       \begin{split}
     %       \delta^2&\tmc_f^\hhh(S)[\Phi]=\delta\tmc^\hhh_f(S)\left[\left\langle\Z,\vh\right\rangle+\alpha\left\langle\Z,T\right\rangle_1\right]+\int_S\varphi\,\jacobi^\hhh\left(f''(H^\hhh)\jacobi^\hhh\varphi\right)\,d\sigma^\hhh\\
     %       &\quad+\int_S 2\varphi\divv^{\hhh,S} \left\langle f'(H^\hhh)A^\hhh\left(\nabla^{\hhh,S}\varphi\right)\right)\,d\sigma^\hhh+\int_S4\varphi\alpha f'(H^\hhh)\Big(\tilde h^\hhh\left(\nabla^{\hhh,S}\varphi,J(\vh)\right)-H^\hhh J(\vh)\varphi\Big)\,d\sigma^\hhh\\
      %      &\quad+\int_S 4\varphi f'(H^\hhh)\s J(\vh)\varphi,d\sigma^\hhh-\int_S\varphi\left(f'(H^\hhh)H^\hhh+f(H^\hhh)\right)\hat\Delta^{\hhh,S}\varphi\,d\sigma^\hhh\\
     %       &\quad-\int_S 4\varphi f''(H^\hhh)\Big(\tilde h^\hhh\left(\nabla^{\hhh,S}H^\hhh,\nabla^{\hhh,S}\varphi\right)+ \s \varphi J(\vh)H^\hhh\Big)\,d\sigma^\hhh\\
      %      &\quad+\int_S\varphi\left(f''(H^\hhh)H^\hhh-2f'(H^\hhh)\right)\left\langle \nabla^{\hhh,S} H^\hhh,\nabla^{\hhh,S}\varphi\right\rangle\,d\sigma^\hhh\\
      %      &\quad+\int_S\varphi^2 f'(H^\hhh)\Big( 2\trace\left(\left(\tilde h^\hhh\right)^3\right)+12\tilde h^\hhh\left(\nabla^{\hhh,S}\alpha,J(\vh)\right)+8\s\alpha+6\tilde h^\hhh(J(\vh),J(\vh))\alpha^2+2 H^\hhh\alpha^2\Big)\,d\sigma^\hhh\\
      %      &\quad+\int_S\varphi^2 \Big ( f(H^\hhh)\left( H^\hhh\right)^2-\left(2f'(H^\hhh) H^\hhh+f(H^\hhh)\right)\left(|\tilde h^\hhh|^2+4J(\vh)\alpha+(2n+2)\alpha^2\right)\Big)\,d\sigma^\hhh.
     %   \end{split}
   % \end{equation*}
    \end{proof}
\section*{Declarations}
\footnotesize{
%\noindent{\it \textbf{Acknowledgements.}} The authors thank Fares Essebei, Andrea Pinamonti, Francesco Serra Cassano and Giacomo Vianello for interesting and valuable conversations on the topic of the paper.
%\smallskip
%
\noindent{%\textbf{Acknowledgments.}} The authors thank R. Monti, J. Pozuelo and M. Ritoré for fruitful discussion about the addressed topics. Part of this research was carried out while S. Verzellesi was visiting the Department of Mathematics at the University of Trento.}

%\smallskip
\noindent{\textbf{Conflict of interest.}} The author has no financial interests or conflicts of interest related to the subject matter.
\smallskip

\noindent \textbf{Data availability statement.} Data sharing not applicable as no datasets were generated or analyzed during the current study.
\bibliographystyle{abbrvmr_macro}
\bibliography{biblio}
\end{document}

